% paralleles.tex — f165r, vv. 209–251 : loci paralleli. Généré par .travail/f165r_lp_redac2/rediger.py
% à partir des collectes vérifiées (paralleles_209-230.json, paralleles_231-251.json) et des parallèles
% d'Austin–Olson 2004 retrouvés dans les sources ; règles : edition/CONVENTIONS.md § 4
\lp{209 ὦ τρισκακόδαιμον]}{Ar.~\textit{Ach.}~400 ὦ τρισμακάρι᾽ Εὐριπίδη ; Ar.~\textit{Thesm.}~875 ὦ τρισκακόδαιμον, ψεύδεται νὴ τὼ θεώ ; Ar.~\textit{Ran.}~19 ὢ τρισκακοδαίμων ἄρ᾽ ὁ τράχηλος οὑτοσί ; Eust.~725.20–40 τοιαῦτα δὲ καὶ ἕτερα οὐκ ὀλίγα, ἐν οἷς καὶ ὁ παρὰ τῷ κωμικῷ τρισκακοδαίμων καὶ ὁ ἀνάπαλιν τρισόλβιος. … παραδείγματα βλασφημιῶν τῶν ἀπὸ ἀριθμοῦ, οἷον, τρισεξώλης ὁ πάνυ ἐξώλης … Σοφοκλῆς δὲ τρίδουλον τὸν ἐκ προγόνων δοῦλον ἔφη, Ἀνακρέων δὲ τρισκεκορημένον τὸν πολλάκις ἐκσεσαρωμένον, Ἀρχίλοχος δὲ τρισοιζυρὴν πόλιν ἔφη τὴν Θάσον ; Ar.~\textit{Ach.}~1024 ὦ τρισκακόδαιμον εἶτα λευκὸν ἀμπέχει ; \textit{idem} Suid.~α~1652 (κακόδαιμον ; ἀμπέχῃ) ; Ar.~\textit{Eccl.}~1098 ὦ τρισκακοδαίμων εἰ γυναῖκα δεῖ σαπρὰν ; Ar.~fr.~339~K.-A. (=~324~Kock ; Θεσμοφοριάζουσαι β´ ; Phot.~α~1065 = Suid.~α~1446) οἴμοι κακοδαίμων τῆς τόθ᾽ ἡμέρας, ὅτε / εἶπέν μ᾽ ὁ κῆρυξ· οὗτος ἀλφάνει ; Ar.~fr.~426~K.-A. (=~416~Kock ; Ὁλκάδες ; Ath.~7.329b) ὦ κακοδαίμων ὅστις ἐν ἅλμῃ πρῶτον τριχίδων ἀπεβάφθη ; Men.~fr.~298.1~K.-A. (=~404~Kock ; Gell.~2.23.14 ; Stob.) ὦ τρισκακοδαίμων, ὅστις ὢν πένης γαμεῖ ; Plaut.~\textit{Aul.}~86 trivenefica ; Soph.~\textit{OC}~372 εἰσῆλθε τοῖν τρὶς ἀθλίοιν ἔρις κακή ; Soph.~\textit{OT}~1063 μητρὸς φανῶ τρίδουλος, ἐκφανεῖ κακή ; \textit{Od.}~5.306 τρὶς μάκαρες Δαναοὶ καὶ τετράκις, οἳ τότ᾽ ὄλοντο ; Verg.~\textit{Aen.}~1.94 O terque quaterque beati ; Aeschin.~1.59 οὓς ἠγάπα ὁ τρισκακοδαίμων ἄνθρωπος ; Σ Pi.~\textit{O.}~1.97g ὅθεν καὶ τρισκακοδαίμονας καὶ τρισαθλίους ἐν τῷ βίῳ τινὰς λέγεσθαι ; Hsch.~κ~327 κακοδαίμων· ὁ θεῷ ἀπεχθής (Dem. 19, 115) ; \textit{idem} Synag.~κ~43 (ἐπαχθής ; add. καὶ ὁ τὴν ψυχὴν ἔχων ἐμπαθῆ), Suid.~κ~156 (ἐπαχθής ; add. καὶ ὁ τὴν ψυχὴν ἔχων ἐμπαθῆ) ; Hsch.~τ~1425 τρίς· ἐπὶ τοῦ πολλάκις (Α 213 . .) ; Hsch.~τ~1428 τρισευδαίμων· τρισμακάριστος (Bacch. 3, 10. Aeschin. epist. 5, 5) ; Hsch.~τ~1429 τρισόλβιος· τρισμακάριστος (cf. Ar. Eccl. 1129. Soph. fr. 837, 1 R.) ; Hsch.~ω~2 ὤ· ἐπίρρημα σχετλιασμοῦ (Α 254 . .) ; Hsch.~ω~3 ὦ· ἄρθρον κλητικῆς πτώσεως. οἷον «ὦ ἄνθρωπε» (Ep. Rom. 2, 1 . .) ; Poll.~6.165 τρισκακοδαίμων, καὶ ἄλλα τοιαῦτα ; Suid.~ω~17 Ὧ δειλακρίων σύ: ἀντὶ τοῦ ὦ κακόδαιμον καὶ δειλότατε ; Ar.~\textit{Nub.}~166 ὦ τρισμακάριος τοῦ διεντερεύματος ; Ar.~\textit{Pax}~1271 καὶ ταῦτ᾽ ὦ τρισκακόδαιμον ; Ar.~\textit{Pax}~1333 ὦ τρὶς μάκαρ ὡς δικαίως ; Ar.~\textit{Av.}~1273 ὦ τρισμακάρι᾽ ; Ar.~\textit{Av.}~1707 ὦ τρισμακάριον πτηνὸν ὀρνίθων γένος ; Men.~fr.~734~K.-A. (=~626.3–4~Kock ; Stob.~3.16.9) ὦ τρισκακοδαίμων, ὅστις ἐκ φειδωλίας / κατέθετο μῖσος διπλάσιον τῆς οὐσίας ; Hor.~\textit{Carm.}~1.13.17 felices ter et amplius ; Prop.~3.12.15 ter quater in casta felix, o Postume, Galla! ; Hsch.~ζ~125 ζεῦγος τρίδουλον ; Hsch.~κ~326 κακοδαιμονίαν· ἀτυχίαν ; Hsch.~τ~1426 τρισαλιτήριος· ἁμαρτωλός (2. Macc. 8, 34 . .) ; Hsch.~τ~1427 τρὶς ἕξ· νικητήριος βόλος (Aesch. Ag. 33). καὶ παροιμία· τρὶς ἓξ ἢ τρεῖς κύβοι. ἐλέγετο δὲ ἐπὶ τοῦ εὐστόχου, ἀπὸ τῶν κυβευτῶν (Pherecr. fr. 129 K.–A.)}
\lp{209 ὡς ἀπόλωλ᾽ Εὐριπίδης]}{Ar.~\textit{Ach.}~473 οἴμοι κακοδαίμων, ὡς ἀπόλωλ᾽ ; \textit{idem} Ar.~\textit{Eq.}~752 ; Ar.~\textit{Pax}~242–243 ἰὼ Πρασιαὶ τρὶς ἄθλιαι καὶ πεντάκις / καὶ πολλοδεκάκις, ὡς ἀπολεῖσθε τήμερον ; \textit{idem} Suid.~π~2225 (τρισάθλιαι ; πολλῷ δεκάκις ; om. ὡς ; σήμερον) ; Ar.~\textit{Thesm.}~77 εἴτ᾽ ἔστ᾽ ἔτι ζῶν εἴτ᾽ ἀπόλωλ᾽ Εὐριπίδης ; Ar.~\textit{Plut.}~850–851 οἴμοι κακοδαίμων, ὡς ἀπόλωλα δείλαιος, / καὶ τρὶς κακοδαίμων καὶ τετράκις καὶ πεντάκις ; Eur.~\textit{Andr.}~71 ὦ δύστηνος, ὡς ἀπωλόμην ; \textit{idem} Eur.~\textit{Hel.}~862 ; Aesch.~\textit{Sept.}~6 Ἐτεοκλέης ἂν εἷς πολὺς κατὰ πτόλιν / ὑμνοῖθ᾽ ὑπ᾽ ἀστῶν ; Soph.~\textit{Aj.}~864 τοῦθ᾽ ὑμὶν Αἴας τοὔπος ὕστατον θροεῖ ; Ar.~\textit{Thesm.}~1212 οἴμ᾽ ὠς ἀπόλωλο ; Eur.~\textit{Cycl.}~665 ὤμοι μάλ᾽, ὡς ὑβρίσμεθ᾽, ὡς ὀλώλαμεν ; Eur.~\textit{Hipp.}~341 τρίτη δ᾽ ἐγὼ δύστηνος ὡς ἀπόλλυμαι ; Eur.~\textit{Hipp.}~1350 ἀπόλωλα τάλας, οἴμοι μοι ; Soph.~\textit{Phil.}~742 ἀπόλωλα, τέκνον, κοὐ δυνήσομαι κακὸν ; Trag. adesp. fr.~91 TrGF II (Ath.~3.107e) ἀπόλωλα· πέπλων μ᾽ ὤλεσαν περιπτυχαί ; Hsch.~α~6474 ἀπόλωλεν· ἀπώλετο (α 166 . .) ; Hsch.~ω~159 ὤλετο· ἀπώλετο (Ι 413 . .) ; Greg. Cypr.~1.80 Ἀπόλωλεν ἡ ὗς καὶ τὸ τάλαντον καὶ ἡ γνάθος}
\lp{210 ὦ φίλτατ᾽ ὦ κηδεστά]}{Ar.~\textit{Ach.}~475 Εὐριπίδιον ὦ φιλτάτιον καὶ γλυκύτατον ; Ar.~\textit{Pax}~1198 ὦ φίλτατ᾽ ὦ Τρυγαῖ᾽ ; Eur.~\textit{Or.}~1045 ὦ φίλτατ᾽, ὦ ποθεινὸν ἥδιστόν τ᾽ ἔχων ; Ar.~\textit{Eq.}~726 ὦ Δημίδιον ⟨ὦ⟩ φίλτατον ; Ar.~\textit{Nub.}~1464 νῦν οὖν ὅπως ὦ φίλτατε ; Ar.~\textit{Vesp.}~1512 ὦ Καρκίν᾽ ὦ μακάριε τῆς εὐπαιδίας ; Ar.~\textit{Av.}~1271–1273 ὦ Πισθέταιρ᾽ ὦ μακάρι᾽ ὦ σοφώτατε, / ὦ κλεινότατ᾽ ὦ σοφώτατ᾽ ὦ γλαφυρώτατε, / ὦ τρισμακάρι᾽ ὦ κατακέλευσον ; Ar.~\textit{Eccl.}~1129 ὦ δέσποτ᾽ ὦ μακάριε καὶ τρισόλβιε ; Eur.~\textit{Tro.}~740 ὦ φίλτατ᾽, ὦ περισσὰ τιμηθεὶς τέκνον ; Soph.~\textit{Phil.}~799–801 ὦ τέκνον ὦ γενναῖον, ἀλλὰ συλλαβὼν / τῷ Λημνίῳ τῷδ᾽ ἀνακαλουμένῳ πυρὶ / ἔμπρησον, ὦ γενναῖε ; Hsch.~φ~542 φίλτατοι· φιλικώτατοι (Ι 198 . .) ; Ar.~\textit{Eq.}~611 ὦ φίλτατ᾽ ἀνδρῶν καὶ νεανικώτατε ; Ar.~\textit{Lys.}~950 ὑπολύομαι γοῦν. ἀλλ᾽ ὅπως ὦ φίλτατε ; Eur.~\textit{Hel.}~1451 Φοίνισσα Σιδωνιὰς ὦ ; Eur.~\textit{Or.}~1246 Μυκηνίδες ὦ φίλαι ; Hsch.~φ~541 φίλτατα· τὰ ὄμματα καὶ προσφιλέστατα (Eur. IA 458 . .)}
\lp{210 κηδεστά]}{Ar.~\textit{Thesm.}~74 οὐ χρῆν σε κρύπτειν ὄντα κηδεστὴν ἐμόν ; Hsch.~κ~2463 κηδεσταί· πενθεροί, οἱ τῆς κόρης γονεῖς. καὶ οἱ τῶν γαμούντων οἰκεῖοι ἐκδίδοντες, καὶ συγγενεῖς, καταχρηστικῶς (Eur. Alc. 731) ; Moer.~κ~38 Hansen (p.~201,~14~Bekker) κηδεστὰς καὶ τοὺς διδόντας καὶ τοὺς λαμβάνοντας τὰς κόρας Ἀττικοί, πενθεροὺς καὶ γαμβροὺς Ἕλληνες ; Synag.~κ~310 κηδεστής· ὁ κατ᾽ ἐπιγαμβρίαν οἰκεῖος, ἢ πενθερός ; \textit{idem} Suid.~κ~1487 (add. οὐ μόνον δὲ ὁ διδούς, ἀλλὰ καὶ ὁ λαμβάνων. οὐ χρῆν σε κρύπτειν ὄντα κηδεστὴν ἐμόν) ; Ar.~\textit{Thesm.}~584 Εὐριπίδην φάσ᾽ ἄνδρα κηδεστήν τινα ; Ar.~\textit{Thesm.}~1165 ὅδ᾽ ἐστὶν οὑν τῇ σανίδι κηδεστὴς ἐμός ; Eur.~\textit{Alc.}~731 δίκας δὲ δώσεις σοῖσι κηδεσταῖς ἔτι ; Eur.~\textit{Hec.}~834 τοῦτον καλῶς δρῶν ὄντα κηδεστὴν σέθεν ; Σ Eur.~\textit{Hec.}~834 ἀντὶ τοῦ γαμβροῦ κηδεστὴν εἶπεν ; Hsch.~κ~2455 κηδεία· οἰκειότης ; Hsch.~κ~2459 κηδείους· τοὺς προγόνους, ἢ κατ᾽ ἐπιγαμίαν οἰκείους, ἢ συγγενεῖς (Τ 294) ; Hsch.~κ~2471 κῆδος· κηδεία. πένθος, λύπη (Ο 245). φθορά. συγγένεια (Ν 464). μνηστεία. θεραπεία. φροντίς ; Poll.~3.30 καὶ κηδεσταὶ μὲν ἀλλήλοις ἑκατέρωθεν, οἱ τὸ κῆδος συνάψαντες, καὶ πηοὶ κατὰ τοὺς ποιητάς ; D.~20.57 ἰδίᾳ μὲν γὰρ ἕκαστος ἡμῶν σκοπεῖ τίς ἄξιός ἐστιν ἑκάστου κηδεστὴς ; Hsch.~κ~2460 κηδεμόνες· φροντισταί, εὐεργέται. καὶ οἱ κατ᾽ ἐπιγαμίαν οἰκείοι (Ψ 163) ; Moer.~τ~19 Hansen (p.~210,~28~Bekker) τίνι κηδεύομεν Ἀττικοί, παρὰ τίνος γυναῖκα λαμβάνομεν Ἕλληνες}
\lp{210 μὴ σαυτὸν προδῷς]}{Ar.~\textit{Eq.}~241 ἀλλαντοπῶλα μὴ προδῷς τὰ πράγματα ; Pl. Com. fr.~66~K.-A. (Λάϊος ; Phot.~α~1168) ἀπαμβρακοῦ καὶ μὴ προδῷς σαυτήν ; Plaut.~\textit{Epid.}~1.1.97a–98 (numérotation Perseus) quia tu tete deseris. — quid faciam? men rogas? ; Eur.~\textit{Alc.}~250 ἔπαιρε σαυτήν, ὦ τάλαινα, μὴ προδῷς ; Eur.~\textit{Andr.}~191 ὅμως δ᾽ ἐμαυτὴν οὐ προδοῦσ᾽ ἁλώσομαι ; Pl.~\textit{Cri.}~45c σαυτὸν προδοῦναι ; X.~\textit{Eph.}~2.4.5 μὴ προδῷς σεαυτὸν ; Hsch.~α~5762 «ἀπαμβρακοῦ»· ἀνδρίζου. καρτέρει ; Synag. (Σ‴) α~1633 «ἀπαμβρακοῦ»· καρτέρει. λέγεται δὲ πρὸς ὠδίνουσαν γυναῖκα. οἱ δὲ ἀνδρίζου. ἄλλοι φυλάττου ; Eur.~\textit{Hel.}~901 τὴν εὐσέβειαν μὴ προδῷς τὴν σήν ποτε ; Aesch.~\textit{Sept.}~251 ὦ ξυντέλεια, μὴ προδῷς πυργώματα ; Soph.~\textit{Aj.}~588 μὴ προδοὺς ἡμᾶς γένῃ ; Lib.~\textit{Ep.}~115.6 μηδὲ προδοῦναι σαυτόν τε κἀμὲ}
\lp{211 πῶς οὖν ποιήσω δῆτα;]}{Soph.~\textit{Phil.}~757 τί δῆτα δράσω; — μή με ταρβήσας προδῷς ; Ar.~\textit{Eq.}~82 πῶς δῆτα πῶς γένοιτ᾽ ἂν ἀνδρικώτατα ; Ar.~\textit{Pax}~200 πῶς οὖν σὺ δῆτ᾽ ἐνταῦθα κατελείφθης μόνος ; Eur.~\textit{Phoen.}~734 τί δῆτα δράσω; πολεμίοισι δῶ πόλιν ; Ar.~\textit{Thesm.}~618 τί δῆτά μ᾽ ἕλκεις ἀσθενοῦσαν ; Hsch.~δ~929 δῆτα· οὕτως. ἢ τοίνυν}
\lp{211–212 τοῦτον μὲν μακρὰ / κλάειν κέλευ᾽]}{Ar.~\textit{Ach.}~1131 κλάειν κελεύων Λάμαχον τὸν Γοργάσου ; Ar.~\textit{Eq.}~433 κλάειν σε μακρὰ κελεύσας ; Ar.~\textit{Vesp.}~584 κλάειν ἡμεῖς μακρὰ τὴν κεφαλὴν εἰπόντες τῇ διαθήκῃ ; Ar.~\textit{Ran.}~34 ἦ τἄν σε κωκύειν ἂν ἐκέλευον μακρά ; Ar.~\textit{Ran.}~178–179 ὡς σεμνὸς ὁ κατάρατος· οὐκ οἰμώξεται; / ἐγὼ βαδιοῦμαι ; Ar.~\textit{Plut.}~612 σὲ δ᾽ ἐᾶν κλάειν μακρὰ τὴν κεφαλήν ; Hor.~\textit{Sat.}~1.10.91 Demetri, teque, Tigelli, / discipularum inter iubeo plorare cathedras ; Hdt.~4.127.4 ἀντὶ δὲ τοῦ ὅτι δεσπότης ἔφησας εἶναι ἐμός, κλαίειν λέγω ; Suid.~κ~1702 Κλάειν ἐῶσι: καταφρονοῦσι, παρ᾽ οὐδὲν λογίζονται ; Plu.~\textit{Prov.Alex.}~1.62 Ἡ Σκυθῶν ῥῆσις: πρὸς Δαρεῖον τὸν Πέρσην ἀπεκρίναντο οἱ Σκύθαι, κλαίειν αὐτόν. ἐπὶ τῶν ἀποτόμως οἰμώζειν τινὰ λεγόντων ; \textit{idem} Diogenian.~5.11 (Ἡ ἀπὸ Σκυθῶν ; ἐπὶ τῶν οἰμώζειν τινὰ λεγόντων praep. ; om. κλαίειν αὐτόν) ; Ar.~\textit{Pax}~255 κλαύσει μακρά ; Ar.~\textit{Av.}~692 Προδίκῳ παρ᾽ ἐμοῦ κλάειν εἴπητε τὸ λοιπόν ; Ar.~\textit{Av.}~1207 οἰμώξει μακρά ; \textit{idem} Ar.~\textit{Plut.}~111 ; Ar.~\textit{Eccl.}~425 δεῖπνον παρέχειν ἅπασιν ἢ κλάειν μακρά ; Ar.~\textit{Plut.}~62 κλάειν ἔγωγέ σοι λέγω ; Eup.~fr.~328.1~K.-A. (=~305~Kock ; Orus ap. Zonar. p.~605) τίς οὑξεγείρας μ᾽ ἐστίν; οἰμώξει μακρά ; Antiph.~fr.~217.6~K.-A. (=~218~Kock ; Φιλοθήβαιος ; Ath.~6.225f) τὸν μὲν γελᾶν, τὸν δ᾽ ἕτερον οἰμώζειν μακρά ; Hippon.~fr.~86.18 Degani ]κλαίειν κελεύ[ων Βού]παλο[ν ; Archestr. ap. Ath.~3.85 (116e–117c) σαπέρδῃ δ᾽ ἐνέπω κλαίειν μακρά, Ποντικῷ ὄψῳ ; Hsch.~ε~6052 ἐρρῶσθαι φράσας· οἰμώζειν εἰπών (Dem. 18, 152 . .) ; Hsch.~κ~2849 κλάει· δακρύει. κλαίει (Ar. Ran. 1066) ; Moer.~κ~46 Hansen (p.~201,~24~Bekker) κλάειν καὶ κάειν σὺν τῷ ᾱ Ἀττικοί, μετὰ τοῦ ῑ Ἕλληνες ; Suid.~κ~1701 Κλάειν: στυγνάζειν, ὀδύρεσθαι. ὁ δὲ θύσας Ἐπικούρῳ καὶ δόξαις ταῖς ἐκείνου κλάειν ἔλεγεν ; Luc.~\textit{Dial.\,meretr.}~10.4 Τὰ μὲν ἄλλα ἡ ἀπὸ Σκυθῶν ῥῆσις ; Ar.~\textit{Lys.}~520 ὀτοτύξεσθαι μακρὰ τὴν κεφαλήν}
\lp{211–212 τοῦτον μὲν … ἐμοὶ δ᾽]}{Ar.~\textit{Thesm.}~64 ὦ δαιμόνιε τοῦτον μὲν ἔα χαίρειν, σὺ δὲ ; Soph.~\textit{El.}~448–449 ἀλλὰ ταῦτα μὲν μέθες· σὺ δὲ / τεμοῦσα κρατὸς βοστρύχων ἄκρας φόβας ; Hdt.~5.24.4 Μίλητον μὲν ἔα καὶ τὴν νεόκτιστον ἐν Θρηίκῃ πόλιν, σὺ δέ μοι ἑπόμενος ἐς Σοῦσα ἔχε τά περ ἂν ἐγὼ ἔχω}
\lp{212 ἐμοὶ δ᾽ ὅ τι βούλει χρῶ]}{Ar.~\textit{Nub.}~439–441 νῦν οὖν ἀτεχνῶς ὅ τι βούλονται / τουτὶ τοὐμὸν σῶμ᾽ αὐτοῖσιν / παρέχω ; Ar.~\textit{Nub.}~453–456 δρώντων ἀτεχνῶς ὅ τι χρῄζουσιν, / κεἰ βούλονται / νὴ τὴν Δήμητρ᾽ ἔκ μου χορδὴν / τοῖς φροντισταῖς παραθέντων ; Herod.~5.6 Βίτιννα· δοῦλός εἰμι, χρῶ ὅτι βούλει ⟨μοι⟩ ; Antipho~5.90 χρῆσθαι ὅ τι ἂν βούλησθε ; \textit{idem} Lys.~6.32 (praem. παραδέδωκεν αὑτὸν ὑμῖν) ; Th.~2.4.7 ξυνέβησαν τοῖς Πλαταιεῦσι παραδοῦναι σφᾶς τε αὐτοὺς καὶ τὰ ὅπλα χρήσασθαι ὅτι ἂν βούλωνται ; Th.~7.85.1 Νικίας Γυλίππῳ ἑαυτὸν παραδίδωσι … καὶ ἑαυτῷ μὲν χρήσασθαι ἐκέλευεν ἐκεῖνόν τε καὶ Λακεδαιμονίους ὅτι βούλονται ; And.~1.11 εἰ δὲ μή, χρῆσθέ μοι ὅ τι ἂν ὑμῖν δοκῇ ; Theopomp. Hist. ap. Ath.~12.14 (518a) ὥσπερ παρ᾽ ἡμῖν οἱ κουρεῖς. παρ᾽ οὓς ὅταν εἰσέλθωσιν, παρέχουσιν ἑαυτοὺς πάντα τρόπον, οὐθὲν αἰσχυνόμενοι τοὺς ὁρῶντας ; [D.]~59.66 ἐπὶ δὲ τοῦ δικαστηρίου ἄνευ ἐγχειριδίου χρῆσθαι ὅ τι ἂν βουληθῇ, ὡς μοιχῷ ὄντι ; Aristid.~\textit{Or.}~49 (II, p.~532 Dindorf) ὅ τι βούλει χρῶ λαβὼν, οὐ παραιτοῦμαι ; \textit{idem} Philostr.~\textit{VA}~8.7 (ὅ τι βούλει χρῶ tantum) ; Lib.~\textit{Decl.}~23.79 νῦν δὲ σός εἰμι καὶ χρῶ μοι λαβών ; Lys.~7.40 παρέσχον ἐμαυτὸν ὅ τι βούλεσθε χρῆσθαι ; Ar.~\textit{Pax}~206 ὑμᾶς παραδόντες δρᾶν ἀτεχνῶς ὅ τι βούλεται ; Ar.~\textit{Thesm.}~751–752 μὴ δῆθ᾽, ἱκετεύω σ᾽· ἀλλ᾽ ἔμ᾽ ὅ τι χρῄζεις ποίει / ὑπέρ γε τούτου ; X.~\textit{An.}~4.8.11 καὶ τοῖς περιττοῖς χρήσονται ὅ τι ἂν βούλωνται ; X.~\textit{Cyr.}~8.3.46 λαβὼν γάρ, ἔφη, ταῦτα πάντα κέκτησο, καὶ χρῶ ὅπως βούλει αὐτοῖς ; X.~\textit{Eph.}~2.5.4 ὅ τι βούλει ποίει, καὶ χρῶ σώματι ὡς οἰκέτου ; Epict.~2.16.42 χρῶ μοι λοιπὸν εἰς ὃ ἂν θέλῃς ; Ael.~\textit{VH}~13.37 καὶ δίδωμι χρῆσθαι ὅ τι βούλεσθε ; Julian.~\textit{Or.}~7.22 (232) χρῆσθέ μοι πρὸς ὅ, τι βούλεσθε ; Hsch.~χ~723 χρῆσθαι· καυχᾶσθαι. πράττειν. προσφέρεσθαι, χρησμῷ χρῆσθαι ; Ar.~\textit{Vesp.}~723 ὅ τι βούλει σοι ; Ar.~\textit{Pax}~1204 καὶ τῶνδ᾽ ὅ τι βούλει προῖκα· καὶ ταυτὶ δέχου ; Ar.~\textit{Lys.}~98 ἐπερήσομαί τι μικρόν. — ὅ τι βούλει γε σύ ; Ar.~\textit{Ran.}~3 νὴ τὸν Δί᾽ ὅ τι βούλει γε, πλὴν πιέζομαι ; Epicur.~\textit{Sent.\,Vat.}~51 χρῶ ὡς βούλει τῇ σεαυτοῦ προαιρέσει ; LXX \textit{Es.}~3.11 τῷ δὲ ἔθνει χρῶ ὡς βούλει ; D.H.~7.62.2 παραδιδόντες ἑαυτοὺς ἀντ᾽ ἐκείνου χρῆσθαι, ὅ τι βούλονται ; Arr.~\textit{An.}~3.23.8 παραδιδόναι σφᾶς αὐτοὺς ἐπιτρέποντας Ἀλεξάνδρῳ χρῆσθαι ὅ τι βούλοιτο ; Hsch.~β~920 βούλει· βούλῃ}
\lp{212 λαβών]}{Ar.~\textit{Thesm.}~253 ὅ τι; τὸν κροκωτὸν πρῶτον ἐνδύου λαβών ; Ar.~\textit{Ran.}~616 βασάνιζε γὰρ τὸν παῖδα τουτονὶ λαβών ; Soph.~\textit{Ant.}~398–399 καὶ νῦν, ἄναξ, τήνδ᾽ αὐτός, ὡς θέλεις, λαβὼν / καὶ κρῖνε κἀξέλεγχ᾽ ; Ar.~\textit{Ach.}~188 αὗται μέν εἰσι πεντέτεις. γεῦσαι λαβών ; Ar.~\textit{Av.}~56 σὺ δ᾽ οὖν λίθῳ κόψον λαβών ; Ar.~\textit{Thesm.}~285 ὅπως λαβοῦσα θύσω ; Eur.~\textit{Alc.}~1020 λέξω. γυναῖκα τήνδε μοι σῷσον λαβών ; Soph.~\textit{OC}~475 οἰός γε νεαρᾶς νεοπόκῳ μαλλῷ λαβών}
\lp{213 ἄγε νυν ἐπειδή]}{Ar.~\textit{Vesp.}~211 ἄγε νυν, ἐπειδὴ τουτονὶ σεσοβήκαμεν ; Ar.~\textit{Lys.}~1273 ἄγε νυν ἐπειδὴ τἄλλα πεποίηται καλῶς ; Ar.~\textit{Ach.}~98 ἄγε δὴ σὺ βασιλεὺς ἅττα σ᾽ ἀπέπεμψεν φράσον ; Ar.~\textit{Eq.}~1011 ἄγε νυν ὅπως αὐτοὺς ἀναγνώσεσθέ μοι ; Ar.~\textit{Nub.}~489 ἄγε νυν ὅπως, ὅταν τι προβάλλω σοι σοφὸν ; Ar.~\textit{Pax}~1056 ἄγε νυν ἀπάρχου κᾆτα δὸς τἀπάργματα ; Ar.~\textit{Lys.}~1100 ἄγε δὴ Λάκωνες αὔθ᾽ ἕκαστα χρὴ λέγειν ; Ar.~\textit{Thesm.}~107 ἄγε νυν ὄλβιζε Μοῦσα ; Ar.~\textit{Thesm.}~652 ἄγε δὴ τί δρῶμεν ; Ar.~\textit{Thesm.}~765 ἄγε δὴ τίς ἔσται μηχανὴ σωτηρίας ; Ar.~\textit{Thesm.}~947 ἄγε νυν ἡμεῖς παίσωμεν ἅπερ νόμος ἐνθάδε ταῖσι γυναιξίν ; Hsch.~α~406 ἄγε· ἴθι, ἐπίρρημα κελευστικόν (Γ 192 . .)}
\lp{213 σαυτὸν ἐπιδίδως ἐμοί]}{Ar.~\textit{Eq.}~739–740 σαυτὸν δὲ λυχνοπώλαισι καὶ νευρορράφοις / καὶ σκυτοτόμοις καὶ βυρσοπώλαισιν δίδως ; Ar.~\textit{Vesp.}~1129–1130 ἐπειδήπερ γ᾽ ἅπαξ / ἐμοὶ σεαυτὸν παραδέδωκας εὖ ποιεῖν ; Ar.~\textit{Thesm.}~249 Ἀγάθων, ἐπειδὴ σαυτὸν ἐπιδοῦναι φθονεῖς ; Ar.~\textit{Plut.}~781 ἄκων ἐμαυτὸν τοῖς πονηροῖς ἐπεδίδουν ; Soph.~\textit{Phil.}~84 δός μοι σεαυτόν ; \textit{idem} Soph.~\textit{Trach.}~1117 ; Pl.~\textit{Euthd.}~285c παραδίδωμι ἐμαυτὸν Διονυσοδώρῳ τούτῳ ὥσπερ τῇ Μηδείᾳ τῇ Κόλχῳ. ἀπολλύτω με ; Hsch.~ε~1336 ἐκδίδως· προδίδως ; D.~9.66 Κλειτάρχῳ δ᾽ ἐνέδωκεν αὑτόν ; Luc.~\textit{Hermot.}~28 ἢν ἅπαξ ἐπιδῷ τις αὑτὸν τῇ πνεούσῃ ; Luc.~\textit{Hermot.}~86 ἐπιδιδόντα ἐμαυτὸν καὶ κατὰ ῥοῦν συρρέοντα τῷ ὕδατι ; Alciphr.~3.5 οὔ φησιν ἑαυτὴν ἐπιδώσειν ; Hsch.~ε~1338 ἐκδόται· προδόται, ἀποδόται ; Hsch.~ε~1339 ἔκδοτος· προδεδομένος, ἀποδεδομένος (Eur. Ion. 1251) ; Ar.~\textit{Pax}~333 ἐπιδίδωμι τοῦτό γ᾽ ὑμῖν ὥστε μὴ λυπεῖν ἔτι ; Cratin.~fr.~93~K.-A. (=~86~Kock ; Κλεοβουλῖναι ; Poll.~10.105) ἐπέδωκε βαλάνων ἄβακα τῶν ἐκ Φιττέως ; Hermipp.~fr.~43.2~K.-A. (=~42.2~Kock ; Μοῖραι ; Σ Ar.~\textit{Ran.}~320) ἐὰν τοσοῦτον ἐπιδιδῷ τῆς ἡμέρας ; Eur.~\textit{Med.}~186 μόχθου δὲ χάριν τήνδ᾽ ἐπιδώσω ; \textit{Il.}~23.559 Εὐμήλῳ ἐπιδοῦναι, ἐγὼ δέ κε καὶ τὸ τελέσσω ; Hes.~\textit{Op.}~396 ὡς καὶ νῦν ἐπ᾽ ἔμ᾽ ἦλθες· ἐγὼ δέ τοι οὐκ ἐπιδώσω ; Thgn.~561–562 εἴη μοι τὰ μὲν αὐτὸν ἔχειν, τὰ δὲ πόλλ᾽ ἐπιδοῦναι / χρήματα τῶν ἐχθρῶν τοῖσι φίλοισιν ἔχειν ; Hdt.~2.30.5 διδοῦσι σφέας αὐτοὺς τῷ Αἰθιόπων βασιλέι ; Th.~4.11.4 τὰς ναῦς τοῖς Λακεδαιμονίοις ἐν τῷ παρόντι ἐπιδοῦναι ; Plu.~\textit{Pomp.}~6.4 τῶν δὲ ἄλλων ἀσμένως ἐπιδιδόντων αὑτούς ; Luc.~\textit{Dom.}~10 ὅλον ἐπιδοὺς ἑαυτὸν τῷ τάχει ; J.~\textit{BJ}~4.164 τὴν ἐμαυτοῦ ψυχὴν ἐπιδώσω μόνην ὑπὲρ τοῦ θεοῦ ; Ach. Tat.~8.17.3 Δοῦλον οὖν με σεαυτῆς ἀπὸ ταύτης τῆς ἡμέρας νόμιζε· καί σοι προῖκα ἐπιδίδωμι τὸ μὲν πρῶτον ἐμαυτόν}
\lp{214 ἀπόδυθι τουτὶ θοἰμάτιον]}{Ar.~\textit{Vesp.}~1131–1132 τί οὖν κελεύεις δρᾶν με; — τὸν τρίβων᾽ ἄφες, / τηνδὶ δὲ χλαῖναν ἀναβαλοῦ τριβωνικῶς ; Ar.~\textit{Av.}~947 ἀπόδυθι· δεῖ γὰρ τὸν ποιητὴν ὠφελεῖν ; Cratin.~fr.~330~K.-A. (Phot.~α~2489) ἀπόδυθι τὴν στολήν ; Moer.~α~47 Hansen (p.~188,~24~Bekker) ἀπόδυθι Ἀττικοί, ἀπόδυσαι Ἕλληνες ; Synag. (Σ‴) α~1835 ἀπόδυσαι· μᾶλλον δὲ ἀπόδυθι οἱ Ἀττικοὶ λέγουσιν ; \textit{idem} Phot.~α~2489 (add. … παρατηρητέον δὲ ὡς οὕτω τὰ ῥήματα λέγοντες οὐκ ἐπισυνάπτουσι τὰ ἱμάτια ἢ τῶν ἱματίων ἀρκούμενοι μόνῃ τῇ ῥηματικῇ φωνῇ δηλῶσαι τὸ νόημα. πλὴν Κρατῖνος „ἀπόδυθι“, φησί, „τὴν στολήν“) ; Ar.~\textit{Thesm.}~731 σὺ δὲ τὸ Κρητικὸν / ἀπόδυθι ταχέως ; Luc.~\textit{Vit.Auct.}~6 ἀλλὰ ἀπόδυθι, καὶ γυμνὸν γάρ σε ἰδεῖν βούλομαι ; Hsch.~α~6322 ἀπόδυθι· ἀπόδυσαι (Ar. Av. 934 . .) ; Com. adesp. fr.~274~Kock (Bekker, Anecd.~427.26) ἀπόδυθι ταχέως, Πάμφιλ᾽· ἡμέτερος ὁ πλοῦς ; Ael.~\textit{VH}~12.22 καὶ θοιμάτιον ἀποδὺς}
\lp{214 τουτὶ θοἰμάτιον]}{Ar.~\textit{Nub.}~497 ἴθι νυν κατάθου θοἰμάτιον. — ἠδίκηκά τι ; Hsch.~θ~622 θοἰμάτιον· ⟨τὸ⟩ ἱμάτιον ; \textit{idem} Synag.~θ~92 (θοιμάτιον ; τὸ) ; Moer.~θ~8 Hansen (p.~198,~30~Bekker) θοιμάτιον Ἀττικοί, τὸ ἱμάτιον Ἕλληνες ; Ar.~\textit{Thesm.}~250 ἀλλ᾽ ἱμάτιον γοῦν χρῆσον ἡμῖν τουτῳί ; Ar.~\textit{Thesm.}~568 λαβὲ θοἰμάτιον Φιλίστη ; Ar.~\textit{Thesm.}~656 ξυζωσαμένας εὖ κἀνδρείως τῶν θ᾽ ἱματίων ἀποδύσας ; Ar.~\textit{Eccl.}~535 εἶτ᾽ οὐ τὸ σαυτῆς ἱμάτιον ἐχρῆν σ᾽ ἔχειν ; Ar.~\textit{Vesp.}~1056–1059 ἐσβάλλετέ τ᾽ ἐς τὰς κιβωτοὺς / μετὰ τῶν μήλων. κἂν ταῦτα ποιῆθ᾽, / ὑμῖν δι᾽ ἔτους τῶν ἱματίων / ὀζήσει δεξιότητος ; Ar.~\textit{Thesm.}~261 φέρ᾽ ἔγκυκλον· — τουτὶ λάβ᾽ ἀπὸ τῆς κλινίδος ; Ar.~\textit{Thesm.}~1181 φέρε θοἰμάτιον ἄνωθεν ὦ τέκνον τοδί ; Ar.~\textit{Eccl.}~527 ᾤχου σιωπῇ θοἰμάτιον λαβοῦσά μου ; Hsch.~ο~1908 οὑτοσί· διὰ τοῦ ι̅, καὶ ὁδί}
\lp{214 καὶ δὴ χαμαί]}{Ar.~\textit{Ach.}~342 οὑτοιί σοι χαμαί ; Ar.~\textit{Av.}~175 βλέψον κάτω. — καὶ δὴ βλέπω ; Ar.~\textit{Lys.}~637 ἀλλὰ θώμεσθ᾽ ὦ φίλαι γρᾶες ταδί πρῶτον χαμαί ; Ar.~\textit{Lys.}~925 κατάκεισο, καὶ δὴ ᾽κδύομαι ; Hsch.~κ~213 καὶ δή· ἐπὶ τοῦ ἤδη (Β 135) ; Ar.~\textit{Eq.}~22 λέγε δὴ μόλωμεν ξυνεχὲς ὡδὶ ξυλλαβών. — καὶ δὴ λέγω· μόλωμεν ; Ar.~\textit{Eq.}~155 ἄγε δὴ σὺ κατάθου πρῶτα τὰ σκεύη χαμαί ; \textit{idem} Ar.~\textit{Pax}~886 ; Ar.~\textit{Nub.}~778 εἰπὲ δή. — καὶ δὴ λέγω ; Ar.~\textit{Nub.}~1097 πλείους σκόπει. — καὶ δὴ σκοπῶ ; Ar.~\textit{Thesm.}~266 ἁνὴρ μὲν ἡμῖν οὑτοσὶ καὶ δὴ γυνὴ ; Ar.~\textit{Eccl.}~1014 καὶ δή σοι λέγω ; Pl. Com. fr.~71.8–9~K.-A. (=~69.8–9~Kock ; Λάκωνες ; Ath.~15.665b) νεοκρᾶτά τις ποιείτω. — καὶ δὴ κέκραται ; Soph.~\textit{OC}~173 πρόσθιγέ νύν μου. — ψαύω καὶ δή ; Σ Ar.~\textit{Thesm.}~214 ἀληθῶς ; Ar.~\textit{Av.}~268 καὶ δή τις ὄρνις ἔρχεται ; Ar.~\textit{Av.}~550 καὶ δὴ τοίνυν πρῶτα διδάσκω μίαν ὀρνίθων πόλιν εἶναι ; Aesch.~\textit{Supp.}~438 καὶ δὴ πέφρασμαι}
\lp{215–216 ἀποξυρεῖν ταδί]}{\textit{Test.} Suid.~α~4565 ἀποξυρεῖν ταδί, τὰ κάτω δ᾽ ἀφαύειν ; Clem.Al.~\textit{Strom.}~6.2.26.4 Ἀριστοφάνης δὲ ὁ κωμικὸς ἐν ταῖς πρώταις Θεσμοφοριαζούσαις τὰ ἐκ τῶν Κρατίνου Ἐμπιπραμένων μετήνεγκεν ἔπη ; Ar.~\textit{Thesm.}~1043–1044 ὃς ἔμ᾽ ἀπεξύρησε πρῶτον, / ὃς ἐμὲ κροκόεντ᾽ ἐνέδυσεν ; Cratin.~fr.~48~K.-A. (=~41~Kock ; Διονυσαλέξανδρος ; Poll.~7.28) νακότιλτος ὡσπερεὶ κωδάριον ἐφαινόμην ; Ath.~13.19 (565d–f) ἐρωμένους περιάγεσθαι ξυρουμένους τὴν ὑπήνην καὶ τὸν ὄρρον ; Zonar.~11.1 (= Cass. Dio~57.10.5) κείρεσθαί μου τὰ πρόβατα, ἀλλ᾽ οὐκ ἀποξυρᾶσθαι βούλομαι ; \textit{idem} Suid.~τ~552 (οὐ ξυρᾶσθαι) ; Σ Ar.~\textit{Thesm.}~215 τὰ γένεια. ταῦτα δὲ ἔλαβεν ἐκ τῶν Ἰδαίων Κρατίνου ; Mant. Prov.~3.16 Τὸν πόκον περικείρεσθαι βούλομαι τῶν προβάτων, ἀλλ᾽ οὐ ξυρᾶσθαι δυσμενῶς ἄχρις ἐπιδερμίδος: Τιβερίου βασιλέως Ῥωμαίων αὕτη ἡ παροιμία. λέγεται δὲ ἐπὶ τῶν βαρέως φορολογούντων ; Ar.~\textit{Ach.}~119 ὦ θερμόβουλον πρωκτὸν ἐξυρημένε ; Ar.~\textit{Lys.}~92 δήλη ᾽στὶν οὖσα ταυταγὶ τἀντευθενί ; Ar.~\textit{Thesm.}~92 λάθρᾳ, στολὴν γυναικὸς ἠμφιεσμένον ; Ar.~\textit{Thesm.}~191 σὺ δ᾽ εὐπρόσωπος λευκὸς ἐξυρημένος ; Eub.~fr.~31~K.-A. (=~32~Kock ; Δόλων ; Phot.~π~1377 ; Poll.~2.29) πρὸς φθεῖρα κείρασθαι ; Eur.~\textit{Alc.}~427 κουρᾷ ξυρήκει καὶ μελαμπέπλῳ στολῇ ; Eur.~\textit{El.}~241 καὶ κρᾶτα πλόκαμόν τ᾽ ἐσκυθισμένον ξυρῷ ; Eur.~\textit{El.}~335 κάρα τ᾽ ἐμὸν ξυρῆκες, ὅ τ᾽ ἐκεῖνον τεκών ; Hdt.~5.35.3 ὁ δὲ τῶν δούλων τὸν πιστότατον ἀποξυρήσας τὴν κεφαλὴν ἔστιξε ; X.~\textit{HG}~1.7.8 ἀνθρώπους μέλανα ἱμάτια ἔχοντας καὶ ἐν χρῷ κεκαρμένους ; Theopomp. Hist. ap. Ath.~6.77 (260e) οὐχ οἱ μὲν ξυρούμενοι καὶ λεαινόμενοι διετέλουν ἄνδρες ὄντες ; Theopomp. Hist. ap. Ath.~12.14 (517f) πάντες δὲ οἱ πρὸς ἑσπέραν οἰκοῦντες βάρβαροι πιττοῦνται καὶ ξυροῦνται τὰ σώματα ; Hsch.~κ~3841 κουρέων· τὸν [ἀκουντίτην] νακοτίλτην ; Hsch.~ξ~184 ξυρεῖ· ἀκονᾷ ; Synag.~ε~466 ἐν χρῷ κεκαρμένος· πρὸς αὐτῷ τῷ χρωτί, οἷον σύνεγγυς καὶ πλησίον τοῦ δέρματος, τὰς τρίχας ἐξυρημένος ; \textit{idem} Suid.~ε~1507 ; Zen.~5.38 Λέγεται δὲ καὶ Ὄνον κείρεις, ἐπὶ τῶν ἀνηνύτοις ἐπιχειρούντων ; Macar.~8.65 Ὗν κείρεις: ἐπὶ τῶν ἀνωφελῶν ; Eur.~\textit{Phoen.}~372 [κάρα ξυρῆκες καὶ πέπλους μελαγχίμους] ; LXX \textit{Is.}~7.20 ξυρήσει Κύριος τῷ ξυρῷ τῷ μεγάλῳ καὶ μεμεθυσμένῳ … τὴν κεφαλήν, καὶ τὰς τρίχας τῶν ποδῶν, καὶ τὸν πώγωνα ἀφελεῖ ; Plu.~\textit{Oth.}~2.3 ἐδεήθη δὲ ὑπομεῖναι ἕως ἂν ἀποξύρηται τὸ γένειον ; Luc.~\textit{Sacr.}~15 τίς οὕτω περὶ πολλοῦ ποιεῖται τὴν κόμην ὅστις οὐκ ἀπεξύρησε ; Hsch.~ε~6243 ἐσκνᾶσθαι· ξυρῆσαι ; Hsch.~κ~2319 κερεῖ· ξυρεῖ ; Hsch.~κ~3858 κούριξαν· ἀπέκειραν ; Suid.~ε~1871 Ἐξυρημένος: τὴν κόμην ἀποβαλών ; Diogenian.~6.99 Ὄνου πόκοι: ἐπὶ τῶν ἀνηνύτων καὶ ἀτελῶν. Οὐδὲ γὰρ κείρεται ὄνος}
\lp{215 ἀτὰρ τί μέλλεις δρᾶν μ᾽;]}{Ar.~\textit{Vesp.}~1379 τί μέλλεις δρᾶν ; Ar.~\textit{Ach.}~412 ἀτὰρ τί τὰ ῥάκι᾽ ἐκ τραγῳδίας ἔχεις ; Ar.~\textit{Vesp.}~1131 τί οὖν κελεύεις δρᾶν με; — τὸν τρίβων᾽ ἄφες ; Eur.~\textit{Hel.}~860 φεῦγ᾽· ἀτὰρ τί φευκτέον ; Eur.~\textit{Hec.}~726 Ἑκάβη, τί μέλλεις παῖδα σὴν κρύπτειν τάφῳ ; Hsch.~δ~2314 δρᾶν· πράσσειν. θύειν}
\lp{216 τὰ κάτω δ᾽ ἀφεύειν]}{Ar.~\textit{Eq.}~373 τὰς βλεφαρίδας σου παρατιλῶ ; Ar.~\textit{Lys.}~827–828 ὄντ᾽ αὐτὸν / κομήτην, ἀλλ᾽ ἀπεψιλωμένον / τῷ λύχνῳ ; Ar.~\textit{Thesm.}~236 ἀνίστασ᾽, ἵν᾽ ἀφεύσω σε, κἀγκύψας ἔχε ; Ar.~\textit{Thesm.}~590 ἀφηῦσεν αὐτὸν κἀπέτιλ᾽ Εὐριπίδης ; Ar.~\textit{Thesm.}~593 ἠλίθιος ὅστις τιλλόμενος ἠνείχετο ; Ar.~\textit{Ran.}~426–428 τὸν Κλεισθένους δ᾽ ἀκούω / ἐν ταῖς ταφαῖσι πρωκτὸν / τίλλειν ἑαυτοῦ καὶ σπαράττειν τὰς γνάθους ; Ar.~\textit{Eccl.}~12–13 μόνος δὲ μηρῶν εἰς ἀπορρήτους μυχοὺς / λάμπεις ἀφεύων τὴν ἐπανθοῦσαν τρίχα ; Ar.~\textit{Plut.}~168 ὁ δ᾽ ἁλούς γε μοιχὸς διὰ σέ που παρατίλλεται ; Aesch.~fr.~310~R. (Ath.~9.17) καὶ καλῶς ἠφευμένος / ὁ χοῖρος ; \textit{Il.}~9.467–468 πολλοὶ δὲ σύες θαλέθοντες ἀλοιφῇ / εὑόμενοι τανύοντο διὰ φλογὸς Ἡφαίστοιο ; \textit{idem} \textit{Il.}~23.32–33 (δ᾽ ἀργιόδοντες ὕες) ; \textit{Od.}~14.74–75 ἔνθεν ἑλὼν δύ᾽ ἔνεικε καὶ ἀμφοτέρους ἱέρευσεν, / εὗσέ τε μίστυλλέν τε καὶ ἀμφ᾽ ὀβελοῖσιν ἔπειρεν ; Semon.~fr.~24~W. (Ath.~14.660a) χὠς ὗν ἄφευσα χὠς ἐμίστυλα κρέα ; Plu.~\textit{Dion}~9.3 ὥστε μηδὲ τῆς κεφαλῆς τὰς τρίχας ἀφελεῖν κουρικαῖς μαχαίραις, ἀλλὰ τῶν πλαστῶν τις ἐπιφοιτῶν ἄνθρακι τὴν κόμην περιέκαιεν ; Σ Il.~23.33 εὑόμενοι ⟨τανύοντο⟩] εὕειν τὸ τὰς τρίχας ἀποκαίειν ; Hsch.~ε~7183 εὗσεν· ἐμάδισεν (ι 389) ; Hsch.~ε~7216 εὕστρα· βόθρος, ἐν ᾧ τοὺς ὗς εὕουσι. καὶ ὁ στάχυς, ὅταν μήπω πέπειρος ὢν ἀποφλογισθῇ ; Hsch.~ε~7308 εὕω· καίω, φλογίζω ; Phot.~ε~2343 εὗστραι· βόθυνοι, ἐν οἷς χοῖροι ἀφεύονται καὶ ἕτερα ἱερεῖα ; Suid.~δ~205 τοῦτο εἴρηκεν ἀφευθεὶς τὸν πρωκτόν. ἀντὶ τοῦ φλογισθείς. μετὰ γὰρ τὸ τυθῆναι τὰ δελφάκια φλογίζονται, ἵνα ψιλωθῶσι ; \textit{idem} Suid.~α~4594 (Ἀφευθεὶς τὸν πρωκτόν: Ἀριστοφάνης ἀντὶ τοῦ φλογισθείς tantum) ; Ar.~\textit{Eq.}~394 ἐν ξύλῳ δήσας ἀφαύει κἀποδόσθαι βούλεται ; Ar.~\textit{Eq.}~1236 ἐν ταῖσιν εὕστραις κονδύλοις ἡρμοττόμην ; Ar.~\textit{Nub.}~1083 τί δ᾽ ἢν ῥαφανιδωθῇ πιθόμενός σοι τέφρᾳ τε τιλθῇ ; Ar.~\textit{Pax}~1144 ἀλλ᾽ ἄφευε τῶν φασήλων ὦ γύναι τρεῖς χοίνικας ; \textit{idem} Suid.~α~4593 (om. ἀλλ᾽) ; Ar.~\textit{Lys.}~89 κομψότατα τὴν βληχώ γε παρατετιλμένη ; Ar.~\textit{Lys.}~151 γυμναὶ παρίοιμεν δέλτα παρατετιλμέναι ; Ar.~\textit{Lys.}~800–804 τὴν λόχμην πολλὴν φορεῖς. — καὶ Μυρωνίδης γὰρ ἦν / τραχὺς ἐντεῦθεν μελάμπυγός / τε τοῖς ἐχθροῖς ἅπασιν ; Ar.~\textit{Thesm.}~238 ἐνεγκάτω τις ἔνδοθεν δᾷδ᾽ ἢ λύχνον ; Ar.~\textit{Thesm.}~537–538 αὐταί τε καὶ τὰ δουλάρια τέφραν ποθὲν λαβοῦσαι / ταύτης ἀποψιλώσομεν τὸν χοῖρον ; Ar.~\textit{Ran.}~422–424 τὸν Κλεισθένους δ᾽ ἀκούω / ἐν ταῖς ταφαῖσι πρωκτὸν / τίλλειν ἑαυτοῦ ; Ar.~\textit{Ran.}~516 ἡβυλλιῶσαι κἄρτι παρατετιλμέναι ; Ar.~\textit{Eccl.}~60–61 πρῶτον μέν γ᾽ ἔχω τὰς μασχάλας / λόχμης δασυτέρας, καθάπερ ἦν ξυγκείμενον ; Ar.~\textit{Eccl.}~724 κατωνάκην τὸν χοῖρον ἀποτετιλμένας ; Ar.~fr.~425.1–2~K.-A. (=~409~Kock ; Ὁλκάδες ; Ath.~3.91b) δαρδάπτοντα, μιστύλλοντα, διαλείχοντά μου / τὸν κάτω σπατάγγην ; Pl. Com. fr.~188.13–14~K.-A. (=~174.13–14~Kock ; Φάων ; Ath.~10.442a) Κονισάλῳ δὲ καὶ παραστάταιν δυοῖν / μύρτων πινακίσκος χειρὶ παρατετιλμένων ; Pl. Com. fr.~188.15~K.-A. (=~174.15~Kock ; Φάων ; Ath.~10.441e–442a) λύχνων γὰρ ὀσμὰς οὐ φιλοῦσι δαίμονες ; Eub.~fr.~106.2~K.-A. (=~107.2~Kock ; Σφιγγοκαρίων ; Ath.~10.449e) οἰκείων ἀνέμων ταμίας, δασύς, ἄλλοτε λεῖος ; Eur.~\textit{Cycl.}~514–515 λύχνα δ᾽ ἀμμένει δαΐα σὸν / χρόα ; \textit{Od.}~2.300 αἶγας ἀνιεμένους σιάλους θ᾽ εὕοντας ἐν αὐλῇ ; \textit{Od.}~9.389 πάντα δέ οἱ βλέφαρ᾽ ἀμφὶ καὶ ὀφρύας εὗσεν ἀυτμὴ ; \textit{Od.}~14.426 τὸν δ᾽ ἔλιπε ψυχή. τοὶ δ᾽ ἔσφαξάν τε καὶ εὗσαν ; Hes.~\textit{Op.}~704–705 ἥτ᾽ ἄνδρα καὶ ἴφθιμόν περ ἐόντα / εὕει ἄτερ δαλοῖο καὶ ὠμῷ γήραϊ δῶκεν ; Σ Ar.~\textit{Eq.}~1236 εὕστραις· ἀντὶ τοῦ φλογίστραις, ὅπου τὰ ἱερεῖα θύεται. λέγει δὲ ὡς μάγειρος. ἄλλως. ἐν τοῖς μαγειρείοις. εὕστρα γὰρ τὸ μαδιστήριον, ἀπὸ τοῦ εὕειν καὶ φλογίζειν τοὺς χοίρους ; Hsch.~ε~7043 εὑόμενοι· φλογιζόμενοι (Ι 468) ; Hsch.~ε~7174 εὖσαι· καῦσαι ; Hsch.~ε~7176 εὔσανα· χύτρα. [ὄρυγμα, ἐν ᾧ τοὺς ὗς βυθίζουσι.] καὶ τὰ ἐγκαύματα ; Hsch.~μ~36 μαδίσας· τὰς τρίχας ἀποβαλών ; Hsch.~ψ~207 ψίλωθρον· σποδή. ἔνιοι μάδον (Nic. ther. 902) ; Ar.~\textit{Ach.}~31 ἀπορῶ γράφω παρατίλλομαι λογίζομαι ; Hsch.~ε~7194 εὕσοντας· ἐκδέροντας ; Hsch.~μ~33 μαδαρός· ἀραιόθριξ. ψεδνός ; Hsch.~τ~890 τιλλομένη· λοιδορουμένη (Ar. Thesm. 543) ; Hsch.~φ~629 φλογίσει· καύσει}
\lp{216 ἀλλὰ πρᾶττ᾽]}{Eur.~\textit{HF}~326 πρᾶσσ᾽· οὐ γὰρ ἀλκὴν ἔχομεν ὥστε μὴ θανεῖν ; Ar.~\textit{Vesp.}~428 ἀλλ᾽ ἀφίει τὸν ἄνδρ᾽ ; Ar.~\textit{Vesp.}~1009 ἀλλ᾽ ἴτε χαίροντες ὅποι βούλεσθ᾽ ; Trag. adesp. fr. *710.2 TrGF II (P. Aberd.~14) πορεύου· πρᾶσσε κἀπιτύγχανε}
\lp{216 εἴ σοι δοκεῖ]}{Ar.~\textit{Ach.}~338 ἀλλὰ γὰρ νῦν λέγ᾽, εἴ σοι δοκεῖ ; Eur.~\textit{Hel.}~993 κτεῖν᾽, εἰ δοκεῖ σοι ; Eur.~\textit{Heracl.}~529 καὶ στεμματοῦτε καὶ κατάρχεσθ᾽, εἰ δοκεῖ ; Ar.~\textit{Av.}~665 ἀλλ᾽ εἰ δοκεῖ σφῷν, ταῦτα χρὴ δρᾶν ; Ar.~\textit{Thesm.}~234 εἰ δοκεῖ, φέρε ; Eur.~\textit{Med.}~742 ἀλλ᾽, εἰ δοκεῖ σοι, δρᾶν τάδ᾽ οὐκ ἀφίσταμαι ; Aesch.~\textit{Ag.}~944 ἀλλ᾽ εἰ δοκεῖ σοι ταῦθ᾽ ; Soph.~\textit{Ant.}~98 ἀλλ᾽ εἰ δοκεῖ σοι, στεῖχε ; Ar.~\textit{Nub.}~11 ἀλλ᾽ εἰ δοκεῖ ῥέγκωμεν ἐγκεκαλυμμένοι ; Ar.~\textit{Nub.}~1482 εἴθ᾽ ὅ τι σοι δοκεῖ ; Ar.~\textit{Lys.}~167 εἴ τοι δοκεῖ σφῷν ταῦτα, χἠμῖν ξυνδοκεῖ ; Pherecr.~fr.~163.3~K.-A. (=~154.3~Kock ; Ψευδηρακλῆς ; Ath.~3.122e) ἀλλ᾽ εἰ δοκεῖ σοι, πρόσεχε τὸν νοῦν κἀκροῶ ; Eur.~\textit{Hipp.}~507 εἴ τοι δοκεῖ σοι, χρῆν μὲν οὔ σ᾽ ἁμαρτάνειν}
\lp{217 ἢ μὴ διδόναι γ᾽ ἑαυτόν]}{Ar.~\textit{Eq.}~739–740 σαυτὸν δὲ λυχνοπώλαισι καὶ νευρορράφοις / καὶ σκυτοτόμοις καὶ βυρσοπώλαισιν δίδως ; Soph.~\textit{Phil.}~84 δός μοι σεαυτόν ; \textit{idem} Soph.~\textit{Trach.}~1117 ; Hdt.~6.108.4 ἱκέται ἱζόμενοι ἐπὶ τὸν βωμὸν ἐδίδοσαν σφέας αὐτούς ; Th.~2.68.7 οἱ δ᾽ Ἀμφίλοχοι γενομένου τούτου διδόασιν ἑαυτοὺς Ἀκαρνᾶσι ; Ar.~\textit{Thesm.}~74 ὄντα κηδεστὴν ἐμόν ; Ar.~fr.~469~K.-A. (=~453~Kock ; Πολύιδος ; Stob.~4.22b.43) ἰδοὺ δίδωμι τήνδ᾽ ἐγὼ γυναῖκά σοι / Φαίδραν ; D.~18.219 διὰ παντὸς ἔδωκεν ἑαυτὸν εἰς οὐδὲν τῇ πόλει ; Eur.~\textit{IA}~68 δίδωσ᾽ ἑλέσθαι θυγατρὶ μνηστήρων ἕνα}
\lp{217 μὴ … ὤφελόν ποτε]}{Soph.~\textit{Phil.}~969 οἴμοι, τί δράσω; μή ποτ᾽ ὤφελον λιπεῖν ; Eur.~\textit{IT}~518 ὡς μήποτ᾽ ὤφελόν γε μηδ᾽ ἰδὼν ὄναρ ; Eur.~\textit{Or.}~879 ὃ μήποτ᾽ ὤφελον ; Ar.~\textit{Ach.}~221 μὴ γὰρ ἐγχάνῃ ποτὲ ; Ar.~\textit{Thesm.}~344 καὶ μὴ δίδωσιν ἃν ὑπόσχηταί ποτε ; D.~18.288 παθόντων ἃ μήποτ᾽ ὤφελον}
\lp{217 ὤφελον]}{Ar.~\textit{Nub.}~41 εἴθ᾽ ὤφελ᾽ ἡ προμνήστρι᾽ ἀπολέσθαι κακῶς ; Eur.~\textit{Med.}~1 Εἴθ᾽ ὤφελ᾽ Ἀργοῦς μὴ διαπτάσθαι σκάφος ; Soph.~\textit{OT}~1157 ἔδωκ᾽· ὀλέσθαι δ᾽ ὤφελον τῇδ᾽ ἡμέρᾳ ; Hsch.~ο~1955b ⟨ὄφελον·⟩ εἴθε. μακάρι (Σ 367) ; \textit{idem} Synag.~ο~282, Suid.~ο~994 (Ὄφελες καὶ Ὄφελον ; add. εὐκτικῶς) ; Phot.~ο~702 ὄφελον καὶ ὄφελες καὶ ὄφελε· τοῖς προσώποις συσχηματίζουσιν ὡς ῥῆμα· καὶ διὰ τοῦ ω· ὤφελον ἐγώ· καὶ ὤφελες σύ· καὶ ὤφελεν οὗτος ; Ar.~\textit{Ran.}~955 ὡς πρὶν διδάξαι γ᾽ ὤφελες μέσος διαρραγῆναι ; Enn.~\textit{Medea} ap. Quint.~5.10.84 utinam ne in nemore Pelio ; \textit{Il.}~9.698 μὴ ὄφελες λίσσεσθαι ; Suid.~ω~268 Ὤφελεν: εὐκτικόν. ἔστι δὲ καὶ ὤφελον}
\lp{218 Ἀγάθων σὺ μέντοι]}{Ar.~\textit{Av.}~933–934 οὗτος, σὺ μέντοι σπολάδα καὶ χιτῶν᾽ ἔχεις, / ἀπόδυθι καὶ δὸς τῷ ποιητῇ τῷ σοφῷ ; Ar.~\textit{Eq.}~168 σὺ μέντοι ; Ar.~\textit{Thesm.}~177 Ἀγάθων, σοφοῦ πρὸς ἀνδρός, ὅστις ἐν βραχεῖ}
\lp{218 ξυροφορεῖς]}{Ar.~\textit{Ach.}~848–849 οὐδ᾽ ἐντυχὼν ἐν τἀγορᾷ πρόσεισί σοι βαδίζων / Κρατῖνος ἀεὶ κεκαρμένος μοιχὸν μιᾷ μαχαίρᾳ ; Ar.~\textit{Thesm.}~191 σὺ δ᾽ εὐπρόσωπος λευκὸς ἐξυρημένος ; Ar.~\textit{Eccl.}~65–66 τὸ ξυρὸν δέ γ᾽ ἐκ τῆς οἰκίας / ἔρριψα πρῶτον, ἵνα δασυνθείην ὅλη ; Ar.~fr.~332.1~K.-A. (=~320.1~Kock ; Θεσμοφοριάζουσαι β´ ; Poll.~7.95) ξυρόν, κάτοπτρον, ψαλίδα, κηρωτήν, λίτρον ; Alex.~fr.~266.1–2,~6–7~K.-A. (=~264~Kock ; Ath.~13.565b) ἂν πιττοκοπούμενόν τιν᾽ ἢ ξυρούμενον / ὁρᾷς, … / τί γὰρ αἱ τρίχες λυποῦσιν ἡμᾶς, πρὸς θεῶν; / δι᾽ ἃς ἀνὴρ ἕκαστος ἡμῶν φαίνεται ; Aesch.~\textit{Cho.}~584 ξιφηφόρους ἀγῶνας ὀρθώσαντί μοι ; Hsch.~π~3280 πριαμωθήσομαι· ξυρήσομαι. ἐπειδὴ τὸ τραγικὸν τοῦ Πριάμου πρόσωπον ξυρίας ἐστίν (com. ad. fr. 414 K.–A.) ; Suid.~κ~1756 Κλεισθένην ὁρῶ· Ἀριστοφάνης. οὗτος ἀεὶ ἐξυρᾶτο. λέγεται οὖν ἐπὶ τῶν ἀεὶ ξυρωμένων ; Suid.~μ~1360 Μοιχός: εἶδος κουρᾶς κιναιδώδους. Ἀριστοφάνης· Κρατῖνος ἀεὶ κεκαρμένος μοιχὸν μιᾷ μαχαίρᾳ. ἀπρεπὴς κουρὰ καὶ κιναιδώδης. καὶ κῆπος ἄλλο εἶδος, καὶ σκάφη ἕτερον ; Suid.~σ~1186 Στράτων ἀεὶ ἐξυρᾶτο καὶ Κλεισθένης ἐπὶ τὸ φαίνεσθαι ἀεὶ νέοι. Ἀριστοφάνης· ὦ θερμόβουλον πρωκτὸν ἐξυρημένε ; Ar.~\textit{Lys.}~632 καὶ φορήσω τὸ ξίφος τὸ λοιπὸν ἐν μύρτου κλαδί ; Eur.~\textit{Hel.}~1072 ἀνὴρ παρ᾽ ἄνδρα στήσεται ξιφηφόρος ; Aesch.~\textit{Pers.}~56 τὸ μαχαιροφόρον τ᾽ ἔθνος ; Plu.~\textit{Art.}~29.6 καὶ καταγαγών ἀπέτεμε τῷ ξυρῷ τὸν τράχηλον ; Plu.~\textit{De\,garr.}~13 (509a) ἀδολέσχου κουρέως περιβαλόντος αὐτῷ τὸ ὠμόλινον καὶ πυθομένου πῶς σε κείρω, βασιλεῦ σιωπῶν ἔφη ; Luc.~\textit{Anach.}~34 πρόστιμόν γ᾽ ἔστιν, ὅστις ἐν ἄστει σιδηροφοροίη μηδὲν δέον ; Luc.~\textit{Merc.Cond.}~33 κίναιδόν τινα τῶν πεπιττωμένων τὰ σκέλη καὶ τὸν πώγωνα περιεξυρημένων ; Hsch.~κ~2529 κῆπος· … ἢ εἶδος κούρας, ἣν οἱ θρυπτόμενοι ἐκείροντο ὡς ἐπίπαν ἐν μιᾷ μαχαίρᾳ, καὶ τὸ ἐφήβαιον τῶν γυναικῶν ; Hsch.~σ~15 σάβυττος· εἶδος ξυρήσεως εἰς καλλωπισμόν· πότερον δὲ τοῦ πώγωνος ἢ τῆς κεφαλῆς ἄδηλον. τινὲς δὲ τὸ γυναικεῖον (Eupol. fr. 313,2 K.–A.) ; Phot.~ξ~65 ξυρήκης· ὁ ξυρήσιμος καὶ κουριῶν ; \textit{idem} Suid.~ξ~157 (Ξυρήκης ; καὶ κουριῶν) ; Ar.~\textit{Thesm.}~163 ἐμιτροφόρουν τε καὶ διεκλῶντ᾽ Ἰωνικῶς ; X.~\textit{Cyn.}~10.3 ἔχοντα τὰς λόγχας εὐπλατεῖς καὶ ξυρήκεις ; D.Chr.~64.21 τοὺς μαχαιροφόρους Θρᾷκας ; Hsch.~π~1420 πέξις· ψίλωσις, κουρά, ξύρησις ; Suid.~δ~1539 Δρῶπαξ: ὁ κεκαλλωπισμένος ἀνήρ. πολλοὺς ὁρῶ διὰ τὴν ὑγίειαν ἐπὶ τὸν ξυρὸν καὶ τὸν δρώπακα καταφεύγοντας}
\lp{218 ξυρόν (locutions proverbiales)]}{Soph.~\textit{Aj.}~786 ξυρεῖ γὰρ ἐν χρῷ τοῦτο μὴ χαίρειν τινά ; \textit{Il.}~10.173 νῦν γὰρ δὴ πάντεσσιν ἐπὶ ξυροῦ ἵσταται ἀκμῆς ; Pl.~\textit{R.}~341c ξυρεῖν ἐπιχειρεῖν λέοντα ; Diogenian.~4.41 Ἐκ τριχὸς κρέμαται: ἐπὶ τῶν σφόδρα κινδυνευόντων. Ταὐτὸ δέ ἐστι τῷ, Ἐπὶ ξυροῦ βαίνειν ; \textit{idem} Greg. Cypr.~2.50 (om. σφόδρα ; ὅμοιον τῷ Ἐπὶ ξυροῦ), Apostol.~6.100 (ὅμοιον τῷ· Ἐπὶ ξυροῦ ἀκμῇ ἕστηκεν) ; Diogenian.~6.25 Λέοντα ξυρᾷς: ἐπὶ τῶν καθ᾽ ἑαυτῶν τι ποιούντων ; \textit{idem} Suid.~λ~250 (παροιμία λεγομένη ἐπὶ τῶν τὰ ἀδύνατα ἐπιχειρούντων), Apostol.~10.54 (ξυρᾶν, ἢ κολωνὸν κυρίττειν ; ἐπὶ τῶν μὴ πειθομένων περί τι) ; Diogenian.~6.91 Ξυρὸς εἰς ἀκόνην: πρὸς τοὺς ὧν βούλονται τυγχάνοντας. Ὁμοία τῇ, Ὄνος εἰς ἄχυρα ; \textit{idem} Greg. Cypr.~3.15 (Ξυρὸς εἰς ἀκόνην tantum), Suid.~ξ~160 (Ξυρὸν ; add. παροιμία ; ὄνον), Apostol.~12.24 (ἐπὶ τῶν προσηκόντων τισίν. ἢ ἐπὶ τῶν ὧν βούλονται τυγχανόντων) ; Apostol.~12.23 Ξυρεῖ ἐν χρῷ: ἐπὶ τῶν ποιούντων τι εἴς τινα λίαν ἐπίπονον ; Erasm.~\textit{Adag.}~1.1.20 Novacula in cotem ; Eur.~\textit{HF}~630 ὧδ᾽ ἔβητ᾽ ἐπὶ ξυροῦ ; Aesch.~fr.~99.22~R. (Κᾶρες ἢ Εὐρώπη) λεπτὴ γὰρ ἐλπὶς ἥδ᾽ ἐπὶ ξυροῦ μένει ; Soph.~\textit{Ant.}~996 φρόνει βεβὼς αὖ νῦν ἐπὶ ξυροῦ τύχης ; Thgn.~557 κίνδυνός τοι ἐπὶ ξυροῦ ἵσταται ἀκμῆς ; Simon. ap. Plu. De Herodoti malign.~39 (870e) ἀκμᾶς ἑστακυῖαν ἐπὶ ξυροῦ Ἑλλάδα πᾶσαν ; Theoc.~22.6 ἀνθρώπων σωτῆρας ἐπὶ ξυροῦ ἤδη ἐόντων ; Hdt.~6.11.2 ἐπὶ ξυροῦ γὰρ ἀκμῆς ἔχεται ἡμῖν τὰ πρήγματα ; Luc.~\textit{Iupp.\,trag.}~3 τοῦτο δὴ τὸ τοῦ λόγου, ἐπὶ ξυροῦ ἕστηκεν ; Σ vet. Soph.~\textit{Aj.}~786 ξυρεῖ γὰρ ἐν χρῷ· ἅπτεται τῶν ἀναγκαίων τοῦ χρωτός· καὶ ἔστι παροιμία ἐπὶ τῶν ἐπικινδύνων πραγμάτων ξυρεῖ ἐν χρῷ· ἐνίοτε γὰρ καὶ τοῦ σώματος ἐφάπτεται ὁ σίδηρος ; Hsch.~ε~5023 ἐπὶ ξυροῦ· ἐπὶ κινδύνου. ἢ ἐπὶ τοῦ ὀξυτάτου, ἢ ἀκροτάτου (Κ 173) ; Suid.~ξ~162 Ξυροῦ ἀκμῆς: τοῦ ὀξυτάτου, κατὰ λεπτότητα τῶν σιδηρίων ; Mant. Prov.~1.56 Ἐπὶ ξυροῦ ἀκμῆς ἵσταται: ἐπὶ τῶν ἐν μεγίστῳ κινδύνῳ ὄντων. Ὅμηρος ; \textit{idem} Macar.~4.7 (ἕστηκε τὸ πρᾶγμα pro ἀκμῆς ἵσταται ; ἐπὶ τῶν ἔσχατα κινδυνευόντων ; om. Ὅμηρος) ; Aristid.~\textit{Or.}~46 (II, p.~192 Dindorf) ὥσθ᾽ ὅρα μὴ λέοντα ξυρεῖν ἐπιχειρῶμεν, οὐ Θρασύμαχον συκοφαντεῖν ἐπιχειροῦντες, ἀλλὰ κωμῳδεῖν Περικλέα ; Plaut.~\textit{Bacch.}~241–242 adibo hunc, quem quidem ego hodie faciam hic arietem / Phrixi, itaque tondebo auro usque ad vivam cutem ; Aesch.~\textit{Cho.}~884 ἔοικε νῦν αὐτῆς ἐπὶ ξυροῦ πέλας ; Suid.~ε~2498 ἐπὶ κινδύνου, ἐν αὐτῇ τῇ τοῦ πράγματος ἐπιτάσει ; Erasm.~\textit{Adag.}~2.4.13 (n°~1313) Ad vivum resecare … Mutuo sumpta metaphora a tonsoribus capillos aut ungues resecantibus, nam ii saepenumero molesti sunt, dum nimium diligentes esse student. … Plautus in Bacchidibus sub persona Chrysali : Tondebo auro usque ad vivam cutem. Et hoc ipsum tondere pro deludere Graecis in proverbio est}
\lp{218 ἑκάστοτε]}{Ar.~\textit{Eq.}~1070 οὐ τοῦτό φησιν, ἀλλὰ ναῦς ἑκάστοτε ; Ar.~\textit{Av.}~1440 ὅταν λέγωσιν οἱ πατέρες ἑκάστοτε ; Hsch.~ε~1255 ἑκάστοτε· ἀεί. παρ᾽ ἕκαστα ; Moer.~π~57 Hansen (p.~207,~27~Bekker) πάντοτε οὐδεὶς τῶν Ἀττικῶν ; Phryn.~\textit{Ecl.}~74 Fischer (82 Rutherford) Πάντοτε μὴ λέγε, ἀλλ᾽ ἑκάστοτε καὶ διαπαντός ; Suid.~ε~357 Ἑκάστοτε: ἀντὶ τοῦ ἀεί. οὐ γὰρ λέγεται πάντοτε ; Pl. Com. fr.~168.1~K.-A. (=~153.1~Kock ; Συμμαχία ; Herm. in Pl. Phaedr.~241b) εἴξασιν γὰρ τοῖς παιδαρίοις τούτοις, οἳ ἑκάστοτε γραμμὴν ; Anaxandr.~fr.~58.3~K.-A. (=~57.3~Kock ; Ath.~2.39a) διακονῶ καὶ σεμνός εἰμ᾽ ἑκάστοτε ; Hdt.~2.25.3 οὐδὲ πᾶν τὸ ὕδωρ τὸ ἐπέτειον ἑκάστοτε ἀποπέμπεσθαι τοῦ Νείλου ὁ ἥλιος ; Th.~1.68.2 οὐ περὶ ὧν ἐδιδάσκομεν ἑκάστοτε τὴν μάθησιν ἐποιεῖσθε ; Pl.~\textit{Euthphr.}~3b ὅτι δὴ σὺ τὸ δαιμόνιον φῂς σαυτῷ ἑκάστοτε γίγνεσθαι}
\lp{219 χρῆσόν τι νῦν ἡμῖν ξυρόν]}{\textit{Test.} Poll.~10.140 τὰ δὲ κουρέως σκεύη κτένες, κουρίδες, μαχαιρίδες· … καὶ ψαλὶς δὲ τῶν κουρέως σκευῶν, ἣν καὶ μίαν μάχαιραν καλοῦσιν, καὶ ξυρὸν καὶ ξυροδόχη, ὡς ἐν Θεσμοφοριαζούσαις Ἀριστοφάνης· χρῆσόν τι ἡμῖν ξυρόν. — αὐτὸς λάμβανε ἐντεῦθεν ἐκ τῆς ξυροδόχης ; Ar.~\textit{Ach.}~414–415 ἀλλ᾽ ἀντιβολῶ πρὸς τῶν γονάτων σ᾽ Εὐριπίδη, / δός μοι ῥάκιόν τι τοῦ παλαιοῦ δράματος ; Ar.~\textit{Thesm.}~250 ἀλλ᾽ ἱμάτιον γοῦν χρῆσον ἡμῖν τουτῳὶ ; Eup.~fr.~300~K.-A. (=~278~Kock ; Χρυσοῦν γένος ; Poll.~10.140) ἔπειθ᾽ ὁ κουρεὺς τὰς μαχαιρίδας λαβὼν / ὑπὸ τῆς ὑπήνης κατακερεῖ τὴν εἰσφοράν ; Ar.~\textit{Ran.}~1159 χρῆσον σὺ μάκτραν, εἰ δὲ βούλει, κάρδοπον ; Timocl. (Λήθη) ap. Ath.~9.73 (407f) ἔπειτα «χρῆσόν μοι σύ,» φησί, «τὰς χύτρας ἐν αἷσιν ἕψεις τοὺς κυάμους» ; Hsch.~κ~2828 κιχρήσει· δανείσει ; Phryn.~\textit{Ecl.}~381 Fischer (380 Rutherford) Εὐχρηστεῖν ἀπόρριψον λέγε δὲ κιχράναι ; Ar.~\textit{Thesm.}~257–258 κεκρυφάλου δεῖ καὶ μίτρας. — ἡδὶ μὲν οὖν / κεφαλὴ περίθετος, ἣν ἐγὼ νύκτωρ φορῶ ; Ev.Luc.~11.5 Φίλε, χρῆσόν μοι τρεῖς ἄρτους ; Ael.~\textit{VH}~14.10 «χρῆσόν μοι» ἔφη ‘τὴν ῥυπαρὰν χλαμύδα}
\lp{219 ξυρόν]}{Poll.~10.125–129 (Bethe II, p.~228) τῶν δὲ γυναικείων σκευῶν κτένιον, ξάνιον, κομμώτριον, ξυρόν, κάτοπτρον ; Archipp.~fr.~44~K.-A. (=~42~Kock ; Ῥίνων ; Poll.~10.177) ξυρόν Ἄρχιππος ἐν τῷ Ῥίνωνι ἀρρενικῶς εἴρηκεν ; Hsch.~α~5647 ἀξυρές· [ἄξυρον· ἄστομον] ἄτομον. ἀμβλύ ; Hsch.~ξ~187 ξυρόν· τομόν. ἰσχνόν. ὀξύ ; Phot.~μ~151 μαχαιρίδες· αἱ κουρικαὶ μάχαιραι. οὕτως Εὔπολις (fr. 300 K.-A.) ; Hsch.~ξ~186 ξυρμεύεσθαι· ἀδιακρίτως λαλεῖν. καὶ καυχᾶσθαι. Ἀττικοί}
\lp{219 ξυρόν (prosodie)]}{Moer.~ξ~5 Hansen (p.~204,~18~Bekker) ξυρόν μακρῶς Ἀττικοί, βραχέως Ἕλληνες ; Phot.~ξ~79 ξυρόν· ἐκτεταμένως τὰ πολλὰ λέγουσιν ; Phot.~ξ~78 ξυρεῖσθαι· οὐδὲ ξυρᾶσθαι λέγουσιν}
\lp{219 αὐτὸς λάμβανε]}{Ar.~\textit{Ach.}~432 ὦ παῖ δὸς αὐτῷ Τηλέφου ῥακώματα ; Ar.~\textit{Thesm.}~252 λαμβάνετε καὶ χρῆσθ᾽· οὐ φθονῶ ; Ar.~\textit{Thesm.}~262 ὑποδημάτων δεῖ. — τἀμὰ ταυτὶ λάμβανε ; Hsch.~α~8460 αὐτός· ἢ ἐγώ. ἢ σύ. ἢ ἐκεῖνος}
\lp{219 ξυρόν. — αὐτὸς (anapeste partagé)]}{Ar.~\textit{Ach.}~165 οὐ καταβαλεῖτε τὰ σκόροδ᾽; — ὦ μόχθηρε σὺ}
\lp{220 ἐντεῦθεν]}{Hsch.~ε~3363 ἐντεῦθεν· ἔνθεν, παραυτά, ἐκ τούτου}
\lp{220 ξυροδόκης]}{\textit{Test.} Poll.~2.32 ὧν τὰ ἐργαλεῖα κτένες καὶ ξυρόν, οὗ τὴν θήκην ξυροδόκην Ἀριστοφάνης καλεῖ ; Eust.~130.43 πανδοκεύτριαν γοῦν ὁ κωμικός φησι καὶ ξυροδόκην, καπνοδόκην ; Eust.~1399.38–41 κατὰ δὲ ὁμοιότητα ταύτης τῆς λέξεως, τὴν ξυροδόκην εἶπεν ὁ κωμικός. … Ἀττικὰ γοῦν, ἡ ξυροδόκη. καὶ ἡ πανδοκεύτρια ; \textit{Od.}~1.127–129 ἔγχος μέν ῥ᾽ ἔστησε φέρων πρὸς κίονα μακρὴν / δουροδόκης ἔντοσθεν ἐυξόου, ἔνθα περ ἄλλα / ἔγχε᾽ Ὀδυσσῆος ταλασίφρονος ἵστατο πολλά ; Σ Ar.~\textit{Thesm.}~220 παρὰ τὸ δουροδόκης, ὡς πολλῶν αὐτῶν ὄντων (Rutherford ; ὡς πολλὰ αυ . . Dübner) ; Hsch.~δ~2273 δουροδόκη· δόρατος θήκη (α 128) ; Mart.~11.58.9 Sed fuerit curva cum tuta novacula theca ; Hsch.~ι~733 ἰοδόκον· τὴν φαρέτραν, τῶν ἰῶν οὖσαν δεκτικήν (Ο 444) ; Hsch.~ι~1024 ἱστοδόκη· ἱστοθήκη. τὸ διὰ μέσου ⟨τῆς⟩ νεὼς φράγμα, εἰς ὃ κατακλινόμενος ⟨ὁ⟩ ἱστὸς ἐντίθεται (Α 434) ; Synag.~δ~356 δουροδόκης· τῆς τὰ δόρατα δεχομένης ; \textit{idem} Suid.~δ~1433 (Δουροδόκης) ; A.R.~2.679 ἰοδόκη τετάνυστο κατωμαδόν ; Opp.~\textit{Hal.}~2.356 ἠύτε δουροδόκῃ χαλκήλατον ἔσπασεν αἰχμήν ; Hsch.~π~329 πανδοκ⟨ε⟩ῖον· πανδοχ⟨ε⟩ῖον ; Hdn.~\textit{Pros.\,cath.}~12 ἱστοδόκη, καπνοδόκη}
\lp{220 γενναῖος εἶ]}{Ar.~\textit{Ran.}~179 χρηστὸς εἶ καὶ γεννάδας ; Hsch.~γ~354 γενναῖος· ἀνδρεῖος. εὐγενής ; Zen.~2.88 Γενναῖος εἶ ἐκ βαλαντίου: ἐπὶ τῶν διὰ τὸν πλοῦτον εὐγενῶν εἶναι δοκούντων ; \textit{idem} Diogenian.~4.6 (om. τὸν), Suid.~γ~126 (praem. Γενναῖος. ὁ εὖ γεγονώς. καὶ ; om. τὸν ; add. ὅτι χρημάτων κρείττων ὁ γενναῖος ἀνήρ) ; Ar.~\textit{Ran.}~640 οὐκ ἔσθ᾽ ὅπως οὐκ εἶ σὺ γεννάδας ἀνήρ ; Luc.~\textit{Pisc.}~51 εὖ γε, ὦ Χρύσιππε, γενναῖος εἶ ; Hsch.~γ~352 γεννάδας· γενναῖος. ἀνδρεῖος. εὐγένης. συγγένης. μέγας ; Ar.~\textit{Thesm.}~757 κακῶς ἀπόλοι᾽, ὡς φθονερὸς εἶ καὶ δυσμενής ; Eur.~\textit{El.}~253 πένης ἀνὴρ γενναῖος ἔς τ᾽ ἔμ᾽ εὐσεβής ; Hsch.~γ~351 γεννάδαι· εὐγενέστατοι, γενναιότατοι ; Synag.~γ~25 γεννάδας· γενναῖος, εὐγενέστατος}
\lp{221 κάθιζε]}{Hsch.~κ~171 κάθιξον· κάθισον ; Hor.~\textit{Sat.}~1.7.3 omnibus et lippis notum et tonsoribus esse ; Ar.~\textit{Nub.}~254 κάθιζε τοίνυν ἐπὶ τὸν ἱερὸν σκίμποδα ; Ar.~\textit{Av.}~1441 τοῖς μειρακίοις ἐν τοῖσι κουρείοις ταδί ; Ar.~\textit{Plut.}~338 ἐπὶ τοῖσι κουρείοισι τῶν καθημένων ; Plu.~\textit{Mor.}~177a ἀδολέσχου δὲ κουρέως ἐρωτήσαντος αὐτὸν πῶς σε κείρω; σιωπῶν ἔφη ; Ar.~\textit{Vesp.}~905 σίγα, κάθιζε· σὺ δ᾽ ἀναβὰς κατηγόρει ; D.L.~6.2.66 καὶ γὰρ ἐν κουρείῳ, φησί, κείρομαι}
\lp{221 φύσα τὴν γνάθον τὴν δεξιάν]}{Ar.~\textit{Nub.}~1108–1110 εὖ μοι στομώσεις αὐτόν, ἐπὶ μὲν θάτερα / οἷον δικιδίοις, τὴν δ᾽ ἑτέραν αὐτοῦ γνάθον / στόμωσον οἵαν ἐς τὰ μείζω πράγματα ; Ar.~\textit{Thesm.}~903 αἰσχύνομαί σε τὰς γνάθους ὑβρισμένη ; Ar.~\textit{Eccl.}~818 μεστὴν ἀπῆρα τὴν γνάθον χαλκῶν ἔχων ; D.~19.314 ἴσα βαίνων Πυθοκλεῖ, τὰς γνάθους φυσῶν ; \textit{idem} Demetr.~\textit{Eloc.}~269 (πορεύεται διὰ τῆς ἀγορᾶς τὰς γνάθους φυσῶν, τὰς ὀφρῦς ἐπηρκώς), Procop.~\textit{Arc.}~5.32 (καὶ τὰς γνάθους φυσῶν) ; Luc.~\textit{Merc.Cond.}~3 πρὸς τὴν γνάθον πεφυσημένην ἀποπειρώμενος ; Σ Ar.~\textit{Thesm.}~221 ἵνα ἀποτεταμένης αὐτῆς κάλλιον ξυρήσῃ (Rutherford ; καλ . . . νεν Dübner) ; Ar.~\textit{Thesm.}~575 ὅτι μὲν φίλος εἴμ᾽ ὑμῖν, ἐπίδηλος ταῖς γνάθοις ; \textit{Batr.}~17 εἰμὶ δ᾽ ἐγὼ βασιλεὺς Φυσίγναθος, ὃς κατὰ λίμνην ; Hor.~\textit{Sat.}~1.1.20–21 iratus buccas inflet ; Hsch.~γ~701 γναθμοί· γνάθοι. σιαγόνες (π 175) ; Hsch.~γ~703 γνάθους θηλείας· ἢ γυναῖκας (com. ad. fr. 302 K.–A.) ; Plaut.~\textit{Stich.}~5.6.767 Age, iam infla buccas, nunciam aliquid suaviter ; Pers.~5.13 tumidas intendis rumpere buccas ; Lib.~\textit{Decl.}~33.42 παρ᾽ οἷς οὐδὲν ἂν εὕροις ἢ γνάθους πεφυσημένας, τὰ δ᾽ ἔνδον λιμὸς πολύς ; Lib.~\textit{Or.}~2.46 καὶ τῶν μὲν τεταπεινωμένων, τῶν δὲ τὰς γνάθους φυσώντων ; Hsch.~γ~126 γαμφηλαί· γνάθοι, σιαγόνες. ⟨Ὅμηρος δὲ ἐπὶ ἀλόγων (Τ 394)⟩ ; Hsch.~γ~704 γναθοῦν· τὸ εἰς γνάθους τύπτειν (Phryn. com. fr. 29 K.–A.) ; Hsch.~κ~631 καναδοί· σιαγόνες, γνάθοι ; Hsch.~φ~15 φαγόνες· σιαγόνες, γνάθοι ; Hsch.~φ~1040 φυσᾷ· πέπρηται. ἐπήρτηται (Soph. fr. 768, 1 R.) ; Hsch.~φ~1043 φύσας· γεννήσας (Eur. Hel. 87 . .). ἢ τοὺς ἀσκοὺς πληρώσας ; Phot.~ε~2083 ἑτερόγναθος· ἀπειθής. ἢ ὁ ἄπληστος καὶ ἀμφοτέραις ταῖς γνάθοις ἐσθίων ; Diogenian.~6.100 Ὄνου γνάθος: εἰς πολυφάγον}
\lp{221 τὴν γνάθον]}{Epich.~fr.~18.2~K.-A. (=~21 Kaibel ; Βούσιρις ; Ath.~10.411a) βρέμει μὲν ὁ φάρυγξ ἔνδοθ᾽, ἀραβεῖ δ᾽ ἁ γνάθος ; Ar.~\textit{Pax}~237 ὡς αὐτίκα μάλα τὰς γνάθους ἀλγήσετε ; Eur.~\textit{Hipp.}~1223 αἳ δ᾽ ἐνδακοῦσαι στόμια πυριγενῆ γναθμοῖς ; Aesch.~\textit{Cho.}~280 σαρκῶν ἐπαμβατῆρας ἀγρίαις γνάθοις ; Soph.~fr.~507.2~R. (Ποιμένες ; Σ Nic. Ther.~382 ; Suid.~α~1695) κρυμὸν φέρων γναθμοῖσιν ἐξ ἀμφημέρου ; \textit{Il.}~13.671 τὸν βάλ᾽ ὑπὸ γναθμοῖο καὶ οὔατος ; Hippon.~fr.~79.4 Degani (=~79~W. ; P. Oxy.~2174) τὴ]ν γνάθον παρα. [ ; Hippon.~fr.~132 Degani (=~73.4–5~W. ;~62~Bergk ; Epim. Hom., Cramer An. Ox. I~287.28–288.5) † οἱ δέ μεο ὀδόντες ἐν τοῖσι γνάθοισι πάντες κεκινέαται † ; Hdt.~9.83.2 ἐφάνη δὲ καὶ γνάθος κατὰ τὸ ἄνω τῆς γνάθου ἔχουσα ὀδόντας μουνοφυέας}
\lp{222 οἴμοι]}{Ar.~\textit{Nub.}~925 ὤμοι σοφίας — ὤμοι μανίας — ; Ar.~\textit{Thesm.}~241 οἴμοι τάλας. ὕδωρ ὕδωρ ὦ γείτονες ; Hsch.~ο~297 οἴμοι· ἐπίφθεγμα θρήνου (Ioël 1, 15 . .) ; EM~822.34 Ὤμοι, ἐστὶν ἀναφώνησις λύπης δηλωτική. ἐκ τοῦ οἴμοι, κατὰ ἔκτασιν τοῦ ο εἰς ω ; Ar.~\textit{Thesm.}~625 οἴμοι τάλας, οὐδὲν λέγεις ; Eur.~\textit{Hipp.}~799 οἴμοι· τέκνων μοι μή τι συλᾶται βίος ; Aesch.~\textit{Pers.}~253 ὤμοι, κακὸν μὲν πρῶτον ἀγγέλλειν κακά ; Soph.~\textit{OC}~820 οἴμοι. — τάχ᾽ ἕξεις μᾶλλον οἰμώζειν τάδε ; Sapph.~fr.~94.4~V. ὤιμ᾽ ὠς δεῖνα πεπ[όνθ]αμεν}
\lp{222 οἴμοι (la barbe arrachée)]}{Ar.~\textit{Eccl.}~118–121 οὐκ ἂν φθάνοις τὸ γένειον ἂν περιδουμένη … / ἴθι δὴ σὺ περιδοῦ καὶ ταχέως ἀνὴρ γενοῦ ; Ar.~\textit{Ach.}~120–121 τοιόνδε γ᾽ ὦ πίθηκε τὸν πώγων᾽ ἔχων / εὐνοῦχος ἡμῖν ἦλθες ἐσκευασμένος ; Ar.~\textit{Eccl.}~493–494 ὥστ᾽ εἰκὸς ἡμᾶς μὴ βραδύνειν ἔστ᾽ ἐπαναμενούσας / πώγωνας ἐξηρτημένας ; Ar.~\textit{Eccl.}~501–502 ἀλλ᾽ ἐπείγου / ἅπασα καὶ μίσει σάκον πρὸς τοῖν γνάθοιν ἔχουσα}
\lp{222 τί κέκραγας;]}{Eup.~fr.~113~K.-A. (=~97~Kock ; Δῆμοι ; Σ Ar.~\textit{Lys.}~397) τί κέκραγας ὥσπερ Βουζύγης ἀδικούμενος ; Ar.~\textit{Eq.}~274 καὶ κέκραγας ; Ar.~\textit{Vesp.}~415 ὦγαθοὶ τὸ πρᾶγμ᾽ ἀκούσατ᾽, ἀλλὰ μὴ κεκράγετε ; Ar.~\textit{Thesm.}~692 κέκραχθι· τοῦτο δ᾽ οὐδέποτε σὺ ψωμιεῖς ; Pl. Com. fr.~27.3–4~K.-A. (=~28~Kock ; Ἑορταί ; Clem. Al. Strom.~7.6.33) καὶ οὐδὲν ἀφ᾽ ὑὸς γίγνεται πλὴν ὕστριχες καὶ πηλὸς ἡμῖν καὶ βοή ; Aesch.~\textit{PV}~743 σὺ δ᾽ αὖ κέκραγας κἀναμυχθίζῃ; τί που / δράσεις ; Soph.~\textit{Aj.}~1236 ποίου κέκραγας ἀνδρὸς ὧδ᾽ ὑπέρφρονα ; Luc.~\textit{Luct.}~16 ὦ κακόδαιμον ἄνθρωπε, τί κέκραγας ; Suid.~λ~237 Λέλακας: ἀντὶ τοῦ κέκραγας ; Clem.Al.~\textit{Strom.}~7.6.33 (Aesop.) ὅθεν καὶ ὁ Αἴσωπος οὐ κακῶς ἔφη τοὺς ὗς κεκραγέναι μέγιστον, ὅταν ἕλκωνται· συνειδέναι γὰρ αὑτοῖς εἰς οὐδὲν ἄλλο χρησίμοις πλὴν εἰς τὴν θυσίαν ; Ar.~\textit{Ach.}~335 ὡς ἀποκτενῶ, κέκραχθ᾽· ἐγὼ γὰρ οὐκ ἀκούσομαι ; Ar.~\textit{Vesp.}~198 ἔνδον κέκραχθι τῆς θύρας κεκλῃμένης ; Ar.~\textit{Ran.}~982 κέκραγε πρὸς τοὺς οἰκέτας ; Hsch.~κ~2097 κεκραγήσει· κραυγάσει ; Hsch.~κ~2105 κέκραγε· βοᾷ. φανερῶς διαμαρτύρεται (Esai. 15, 4) ; Hsch.~κ~3902 κραγόν· κραυγήν (Ar. Equ. 487) ; Moer.~κ~37 Hansen (p.~201,~13~Bekker) κεκραγμός Ἀττικοί, κραυγή Ἕλληνες ; Phot.~κ~536 κεκραγμόν· τὴν βοήν}
\lp{222–223 ἐμβαλῶ σοι πάτταλον, ἢν μὴ σιωπᾷς]}{Ar.~\textit{Eq.}~375–381 καὶ νὴ Δί᾽ ἐμβαλόντες αὐτῷ / πάτταλον μαγειρικῶς / ἐς τὸ στόμ᾽ ; Ar.~\textit{Eccl.}~284 ὑπαποτρέχειν ἔχουσι μηδὲ πάτταλον ; Aesch.~\textit{PV}~64–65 ἀδαμαντίνου νῦν σφηνὸς αὐθάδη γνάθον / στέρνων διαμπὰξ πασσάλευ᾽ ἐρρωμένως ; Herod.~3.85–86 πρός σοι βαλέω τὸν μῦν τάχ᾽, ἢν πλέω γρύξῃς. / — ἰδού, σιωπῶ· μή με, λίσσομαι, κτείνῃς ; Σ Ar.~\textit{Eq.}~376 μαγειρικῶς : Εἰώθασι γὰρ οἱ μάγειροι πασσάλοις τὰ τῶν χοίρων ἀνοίγοντες στόματα κατανοεῖν εἰ χαλαζῶσι ; Diogenian.~3.98 Γυμνότερος παττάλου: ἐπὶ τῶν σφόδρα ἀπόρων ; \textit{idem} Suid.~γ~492 (ἀπορωτάτων) ; Diogenian.~5.16 Ἥλῳ τὸν ἦλον, καὶ, Πάτταλον ἐξέκρουσας παττάλῳ: ἀντὶ τοῦ, ἁμαρτήματι τὸ ἁμάρτημα θεραπεύεις ; \textit{idem} Apostol.~8.52 (Ἧλος τὸν ἧλον: καὶ· Πάτταλος τὸν πάτταλον ἐξέκρουσεν ; ἐπὶ τῶν ἰωμένων δι᾽ ἁμαρτήματος ἁμάρτημα), Suid.~η~259 (Ἥλῳ τὸν ἧλον ἐκκρούει ; om. καὶ, Πάτταλον ἐξέκρουσας παττάλῳ ; παροιμία ἀντὶ τοῦ ἁμαρτήματι τὸ ἁμάρτημα σπεύδεις ἐξελάσαι· τὸ δὲ οὐχ οἷόν τε) ; Arist.~\textit{Pol.}~5.1314a ἥλῳ γὰρ ὁ ἧλος, ὥσπερ ἡ παροιμία ; Ar.~\textit{Eq.}~371 διαπατταλευθήσει χαμαί ; Aesch.~\textit{Ag.}~235–238 στόματός / τε καλλιπρῴρου φυλακᾷ κατασχεῖν / φθόγγον ἀραῖον οἴκοις, / βίᾳ χαλινῶν τ᾽ ἀναύδῳ μένει ; Hsch.~δ~1204 διαπατταλευθήσει· τρόπῳ βύρσης διαταθήσῃ παττάλοις (Ar. Equ. 371) ; Hsch.~π~1093 πατάλους· πασσάλους (Ar. Vesp. 129. fr. 432 K.–A.) ; Suid.~α~946 ταύτην Φασὶ καὶ σπόγγῳ πάτταλον κρούειν ; Erasm.~\textit{Adag.}~1.2.4 (n°~104) Clavum clavo pellere. Ἥλῳ τὸν ἧλον ἐκκρούειν, id est Clavum clavo pellere. Καὶ πάτταλον παττάλῳ ἐξέκρουσας, id est Paxillum paxillo pepulisti, id est malum alio malo depulisti ; Luc.~\textit{Philops.}~9 καὶ ἥλῳ, φασίν, ἐκκρούεις τὸν ἧλον ; \textit{idem} Luc.~\textit{Apol.}~9 (om. καὶ ; ἐκκρούων), Luc.~\textit{Laps.}~7 (ἄλλῳ ἥλῳ ἐκκρούειν τὸν ἧλον) ; Ar.~\textit{Vesp.}~129 ὁ δ᾽ ὡσπερεὶ κολοιὸς αὑτῷ παττάλους ; Ar.~\textit{Vesp.}~808 παρά σοι κρεμήσετ᾽ ἐγγὺς ἐπὶ τοῦ παττάλου ; Luc.~\textit{Vit.Auct.}~22 αἷς συμποδίζω τοὺς προσομιλοῦντας καὶ ἀποφράττω καὶ σιωπᾶν ποιῶ, φιμὸν ἀτεχνῶς αὐτοῖς περιτιθείς ; Hsch.~π~1066 πασσακίῳ· πασσάλῳ (Ar. Ach. 763) ; Hsch.~π~1084 παστάξ· πάσσαλος ; Suid.~δ~685 τὰς γὰρ βύρσας ἐκτείνοντες ἐπὶ τῆς γῆς, ἵνα μὴ συνάγοιντο καὶ συστέλλοιντο ἐκ τῆς τοῦ ἡλίου καύσεως, κατὰ τὰ ἄκρα παττάλοις κατακρούοντες ἐκτείνουσιν}
\lp{223 ἀτταταῖ ἰατταταῖ]}{\textit{Test.} Suid.~ε~2807 ἀττατατταταί· θρηνῶν παρατραγῳδεῖ ; Ar.~\textit{Eq.}~1 ἰατταταιὰξ τῶν κακῶν, ἰατταταῖ ; Ar.~\textit{Thesm.}~1005 ἀτταταῖ ἰατταταῖ ; Ar.~\textit{Ran.}~649 ἀτταταῖ. — τί τἀτταταῖ ; Soph.~\textit{Phil.}~743 ἀτταταῖ· διέρχεται ; Soph.~\textit{Phil.}~790 ἀτταταῖ ; Hsch.~ι~92 ἰατταταιάξ· φεῦ (Ar. Equ. 1) ; Ar.~\textit{Ach.}~1190 ἀτταταῖ ἀτταταῖ ; \textit{idem} Ar.~\textit{Ach.}~1198, Ar.~\textit{Nub.}~707 ; Hdn.~\textit{Orth.}~2.2.1.53 Τὰ σχετλιαστικὰ εἰς αι λήγοντα ἐπιρρήματα χαίρει τῇ αι διφθόγγῳ βαβαί, ἀτταταῖ, παπαῖ ; Hdn.~\textit{Pros.\,cath.}~19 καὶ τὸ ἰατταταῖ δὲ περισπᾶται ; Ar.~\textit{Vesp.}~1338 ἰαιβοῖ, αἰβοῖ ; Ar.~\textit{Thesm.}~945 ἰαππαπαιάξ· ὦ κροκώθ᾽ οἷ᾽ εἴργασαι}
\lp{223 ἢν μὴ σιωπᾷς]}{Ar.~\textit{Lys.}~364 εἰ μὴ σιωπήσει, θενών σου ᾽κκοκκιῶ τὸ γῆρας ; Aesch.~\textit{Ag.}~36–37 τὰ δ᾽ ἄλλα σιγῶ· βοῦς ἐπὶ γλώσσῃ μέγας / βέβηκεν ; Zen.~2.70 Βοῦς ἐπὶ γλώττης: παροιμία ἐπὶ τῶν μὴ δυναμένων παρρησιάζεσθαι ; \textit{idem} Diogenian.~3.48 (γλώσσης ; om. παροιμία), Hsch.~β~968 (γλώσσῃ), Greg. Cypr.~1.95 (γλώσσης ; ἐπὶ τῶν δωροδοκουμένων) ; Erasm.~\textit{Adag.}~1.7.18 Bos in lingua. In eos, qui non audent libere, quod sentiant, dicere ; Hsch.~σ~473 σιγᾷ· σιωπᾷ}
\lp{224 οὗτος σὺ ποῖ θεῖς;]}{Ar.~\textit{Ach.}~564 οὗτος σὺ ποῖ θεῖς ; \textit{idem} Ar.~\textit{Vesp.}~854, Ar.~\textit{Lys.}~728 (αὕτη) ; Ar.~\textit{Thesm.}~689 ἆ ποῖ σὺ φεύγεις; οὗτος οὗτος οὐ μενεῖς ; Pherecr.~fr.~135~K.-A. (=~129~Kock ; Πέρσαι ; Ath.~11.502b) οὗτος σύ, ποῖ τὴν ἀργυρίδα τηνδὶ φέρεις ; Phot.~π~993 ποῖ θεῖς (Ar. Vesp. 854)· ποῦ τρέχεις ; Ar.~\textit{Thesm.}~610 αὕτη σὺ ποῖ στρέφει; μέν᾽ αὐτοῦ ; Ar.~\textit{Thesm.}~1223a ὀρθὴν ἄνω δίωκε. ποῖ θεῖς; οὐ πάλιν ; Ar.~\textit{Eccl.}~703 τοιάδ᾽ ἐροῦσιν· ποῖ θεῖς οὗτος ; Eur.~\textit{Hec.}~1280 οὗτος σύ, μαίνῃ καὶ κακῶν ἐρᾷς τυχεῖν ; Eur.~\textit{Or.}~1567 οὗτος σύ, κλῄθρων τῶνδε μὴ ψαύσῃς χερί ; Soph.~\textit{OT}~532 οὗτος σύ, πῶς δεῦρ᾽ ἦλθες ; Σ Eur.~\textit{Or.}~1567 οὗτος σύ : ταῦτα ἄνωθεν Ὀρέστης ἐκ τοῦ δώματός φησι. B.M. Gu. I. προσφώνημα καταφρονητικόν ; Hsch.~π~2701 ποῖ· εἰς τίνα τόπον. τὸ δὲ ποῦ· ἐν τίνι τόπῳ ; \textit{idem} Phot.~π~999 (add. καὶ εἰς τίνα τόπον) ; Moer.~π~49 Hansen (p.~207,~18~Bekker) ποῖ Ἀττικοί, ποῦ κοινόν ; Ar.~\textit{Thesm.}~930 οὗτος τί κύπτεις ; Soph.~\textit{OT}~1121 οὗτος σύ, πρέσβυ, δεῦρό μοι φώνει βλέπων ; Hsch.~π~2702 ποῖ γῆς· εἰς ποῖον γῆς τόπον (e. g. Ar. Ran. 48. 85)}
\lp{224 εἰς τὸ τῶν σεμνῶν θεῶν]}{Ar.~\textit{Eq.}~1311–1312 καθῆσθαί μοι δοκεῖ / ἐς τὸ Θησεῖον πλεούσαις ἢ ᾽πὶ τῶν σεμνῶν θεῶν ; Eur.~\textit{El.}~1270–1272 δειναὶ μὲν οὖν θεαὶ τῷδ᾽ ἄχει πεπληγμέναι / πάγον παρ᾽ αὐτὸν χάσμα δύσονται χθονός, / σεμνὸν βροτοῖσιν εὐσεβὲς χρηστήριον ; Soph.~\textit{OC}~89–90 ὅπου θεῶν / σεμνῶν ἕδραν λάβοιμι καὶ ξενόστασιν ; Orac. ap. Paus.~7.25.1 φράζεο δ᾽ Ἄρειόν τε πάγον βωμούς τε θυώδεις / Εὐμενίδων, ὅθι χρὴ Λακεδαιμονίους ⟨σ᾽⟩ ἱκετεῦσαι / δουρὶ πιεζομένους. τοὺς μὴ σὺ κτεῖνε σιδήρῳ, / μηδ᾽ ἱκέτας ἀδικεῖν· ἱκέται δ᾽ ἱεροί τε καὶ ἁγνοί ; Th.~1.126.11 καθεζομένους δέ τινας καὶ ἐπὶ τῶν σεμνῶν θεῶν τοῖς βωμοῖς ἐν τῇ παρόδῳ ἀπεχρήσαντο ; Paus.~1.28.6 πλησίον δὲ ἱερὸν θεῶν ἐστιν ἃς καλοῦσιν Ἀθηναῖοι Σεμνάς, Ἡσίοδος δὲ Ἐρινῦς ἐν Θεογονίᾳ ; Σ Ar.~\textit{Eq.}~1312 εἰς τὸ Θησεῖον : Ἐνταῦθα οἱ καταφεύγοντες τῶν οἰκετῶν ἀσυλίαν εἶχον. ἐπὶ τῶν σεμνῶν : Εἰς τὸ τῶν Ἐρινύων ἱερόν. καὶ ἐνταῦθα δὲ οἱ οἰκέται ἔφευγον ; Σ Ar.~\textit{Thesm.}~224 ἀντὶ τοῦ εἰς ⟨τὸ⟩ τῶν Ἐρινύων. ἄσυλον γὰρ εἶχον οἱ καταλαμβάνοντες τὰ ἱερὰ τῶν Ἐρινύων ; Σ Eur.~\textit{Or.}~37 οὐκ ὀνομάζουσι δὲ, ἀλλ᾽ εὐφημιζόμενοι σεμνὰς θεὰς ἢ Εὐμενίδας καλοῦσιν ; Hsch.~σ~408 σεμναὶ θεαί· τὰς Εὐμενίδας οὕτως ἔλεγον, καὶ Ἐρινύας ἐπὶ εὐφημισμῷ (e. g. Soph. OC 89–90) ; Harp. s.v.~σεμναὶ θεαί ΣΕΜΝΑΙ ΘΕΑΙ οὕτω καλοῦσιν Ἀθηναῖοι τὰς Ἐρινύας ; \textit{idem} Suid.~σ~222 (οὕτως ; Ἐριννύας) ; Phot.~σ~140 σεμναὶ θεαί· κατ᾽ εὐφημισμὸν αἱ Ἐρινύες· ὥσπερ αἱ αὐταὶ καὶ Εὐμενίδες ἐκαλοῦντο· ἦσαν δὲ τρεῖς ; Ar.~\textit{Thesm.}~83 κἀν Θεσμοφόροιν μέλλουσι περί μου τήμερον ; Soph.~\textit{Aj.}~837 σεμνὰς Ἐρινῦς τανύποδας ; Aeschin.~1.188 εὐχὰς ὑπὲρ τῆς πόλεως ταῖς σεμναῖς θεαῖς ; Plu.~\textit{Sol.}~12.1 ὡς ἐγένοντο περὶ τὰς σεμνὰς θεὰς καταβαίνοντες ; Luc.~\textit{Bis\,Acc.}~4 σὺ δέ, ὦ θύγατερ, καθεζομένη παρὰ τὰς σεμνὰς θεὰς ἀποκλήρου τὰς δίκας ; Din.~1.47 ἐπιωρκηκὼς μὲν τὰς σεμνὰς θεὰς ἐν Ἀρείῳ πάγῳ ; Σ Luc.~\textit{Bis\,Acc.}~4 σεμνὰς θεάς: τὰς Ἐρινύας· τούτων γὰρ τὸ ἱερὸν πλησίον τοῦ Ἀρείου πάγου ; Suid.~κ~2673 Κυλώνειον ἄγος: ἀπὸ Κύλωνος. ὃν καταφυγόντα ἐπὶ τὰς σεμνὰς θεὰς ἀποσπάσαντες αὐτὸν οἱ περὶ τὸν Περικλέα τὸν Ἀθηναῖον ἀπέκτειαν ; App. Prov.~3.40 Καλλίου· πιθήκου. τὰ δυσχερῆ γὰρ τῶν ὀνομάτων εὐφημότερον εἰώθασιν οἱ Ἀττικοὶ προφέρεσθαι, καὶ τὸν πίθηκον οὖν καλλίαν προσηγόρευσαν· οὕτω καὶ τὰς Ἐριννύας Εὐμενίδας λέγουσιν ; Ar.~\textit{Lys.}~2 ἢ ᾽ς Πανὸς ἢ ᾽πὶ Κωλιάδ᾽ ἢ ᾽ς Γενετυλλίδος ; Ar.~\textit{Thesm.}~726–727 ἀλλὰ τάσδε μὲν λαβεῖν χρῆν σ᾽ ἐκφέρειν τε τῶν ξύλων, / καὶ καταίθειν τὸν πανοῦργον πυρπολεῖν θ᾽ ὅσον τάχος ; D.~21.115 ταῖς σεμναῖς θεαῖς ἱεροποιὸν αἱρεθέντ᾽ ; Arist.~\textit{Rh.}~2.23.12 ταῖς μὲν σεμναῖς θεαῖς καλῶς εἶχεν ἐν Ἀρείῳ πάγῳ δοῦναι τὰ δίκαια ; Philo \textit{Prob.}~140 τὴν ἐπὶ ταῖς Σεμναῖς Θεαῖς πομπὴν ; Hsch.~ε~5868 Ἐρινύς· δαίμων καταχθόνιος. ἢ Ἀφροδίτη. ἢ εἴδωλον ; Hsch.~τ~149 τανύποδας Ἐρινῦς· ταχύποδας, ἐν τάχει τιμωρουμένας (Soph. Ai. 837)}
\lp{225 οὐ γὰρ μὰ τὴν Δήμητρά γ᾽ ἐνταυθοῖ μενῶ]}{Ar.~\textit{Vesp.}~1442 οὔτοι μὰ τὴν Δήμητρ᾽ ἔτ᾽ ἐνταυθοῖ μενεῖς ; \textit{idem} Ar.~\textit{Nub.}~814 (Ὁμίχλην pro Δήμητρ᾽)}
\lp{225 μὰ τὴν Δήμητρα]}{Ar.~\textit{Lys.}~273 οὐ γὰρ μὰ τὴν Δήμητρ᾽ ἐμοῦ ζῶντος ἐγχανοῦνται ; Ar.~\textit{Eq.}~698 οὔτοι μὰ τὴν Δήμητρ᾽, ἐὰν μή σ᾽ ἐκφάγω ; Ar.~\textit{Ran.}~1222 οὐδ᾽ ἂν μὰ τὴν Δήμητρα φροντίσαιμί γε ; Ar.~\textit{Plut.}~64 οὔ τοι μὰ τὴν Δήμητρα χαιρήσεις ἔτι ; Antiph.~fr.~26.1–2~K.-A. (=~25.1–2~Kock ; Ἀλείπτρια ; Ath.~3.123b) ἐὰν δὲ τοὐργαστήριον / ποιῆτε περιβόητον, κατασκεδῶ, νὴ τὴν φίλην Δήμητρα ; Philippid.~fr.~5.4–5~K.-A. (=~5.4–5~Kock ; Ἀνανεοῦσα ; Ath.~9.384e) «καλοί γε» φησίν, «οἱ νεφροί, νὴ τὴν φίλην / Δήμητρα.» ; Ar.~\textit{Eq.}~1021 ταυτὶ μὰ τὴν Δήμητρ᾽ ἐγὼ οὐκ οἶδ᾽ ὅ τι λέγει ; Ar.~\textit{Av.}~11 οὐδ᾽ ἂν μὰ Δία γ᾽ ἐντεῦθεν Ἐξηκεστίδης ; Ar.~\textit{Thesm.}~20 νὴ τὸν Δί᾽ ἥδομαί γε τουτὶ προσμαθών ; Ar.~\textit{Thesm.}~623–624 νὴ Δία / ὁσέτη γε ; Ar.~\textit{Ran.}~42 οὔ τοι μὰ τὴν Δήμητρα δύναμαι μὴ γελᾶν ; Ar.~\textit{Ran.}~668 οὔ τοι μὰ τὴν Δήμητρα δύναμαί πω μαθεῖν ; Ar.~\textit{Plut.}~364 οὔ τοι μὰ τὴν Δήμητρ᾽ ὑγιαίνειν μοι δοκεῖς ; Machon ap. Ath.~13.43 (579d–581a) ναὶ τὴν φίλην Δήμητρα, δικαίως τοι δέρος, / ἄνθρωπε, φησίν, ἐξεδάρης ἀκόλαστος ὤν ; Hsch.~μ~1 μά· … συζεύγνυσι καὶ κατομόσει καὶ ἀπομόσει, τοῦτ᾽ ἔστι καὶ συγκαταθέσει καὶ ἀποφάσει, τῷ τε ναί … καὶ τῷ οὔ … ἐπιρρήμασιν}
\lp{225 ἐνταυθοῖ]}{Cratin.~fr.~39~K.-A. (=~37~Kock ; Διονυσαλέξανδρος ; Poll.~10.140) ἔνεισι δ᾽ ἐνταυθοῖ μάχαιραι κουρίδες, / αἷς κείρομεν τὰ πρόβατα καὶ τοὺς ποιμένας ; Ammon.~\textit{Diff.}~170 ἐνταυθοῖ καὶ ἐνταῦθα καὶ ἐνθάδε διαφέρει. ἐνταυθοῖ μὲν γὰρ τὴν ἐν τόπῳ σημασίαν δηλοῖ, ἐνταῦθα δὲ καὶ τὴν ἐν τόπῳ καὶ τὴν εἰς τόπον ; Eur.~\textit{IT}~1010 ἄξω δέ γ᾽, ἤνπερ καὐτὸς ἐνταυθοῖ περῶ ; \textit{Il.}~21.122 ἐνταυθοῖ νῦν κεῖσο μετ᾽ ἰχθύσιν ; \textit{Od.}~18.105 ἐνταυθοῖ νῦν ἧσο σύας τε κύνας τ᾽ ἀπερύκων ; Σ Il.~21.122 τὸ ἐνταυθοῖ περισπαστέον· ἔστι γὰρ ἀπὸ τοῦ ἐνταῦθα Ἀττικοῦ ; Hsch.~ε~3322 ἐνταυθοῖ· ἐνταῦθα (Φ 122) ; Ar.~\textit{Nub.}~843 ἀλλ᾽ ἐπανάμεινόν μ᾽ ὀλίγον ἐνταυθοῖ χρόνον ; Ar.~\textit{Lys.}~4 νῦν δ᾽ οὐδεμία πάρεστιν ἐνταυθοῖ γυνή ; Ar.~\textit{Lys.}~568 ὑπενεγκοῦσαι τοῖσιν ἀτράκτοις τὸ μὲν ἐνταυθοῖ τὸ δ᾽ ἐκεῖσε ; Ar.~\textit{Ran.}~273 τί ἔστι τἀνταυθοῖ; — σκότος καὶ βόρβορος ; Hsch.~ε~3321 ἐνταῦθα· καὶ τὸ χρονικόν. καὶ τὸ ἐπὶ τόπου. καὶ ἔνθα δή, ἢ ὧδε}
\lp{226 οὔκουν … δῆτ᾽]}{Ar.~\textit{Eq.}~878 οὔκουν σε δῆτα ταῦτα δεινόν ἐστι πρωκτοτηρεῖν ; Ar.~\textit{Ran.}~193 οὔκουν περιθρέξει δῆτα τὴν λίμνην κύκλῳ ; Ar.~\textit{Ran.}~200 οὔκουν καθεδεῖ δῆτ᾽ ἐνθαδί, γάστρων}
\lp{226 καταγέλαστος … ἔσει]}{Ar.~\textit{Lys.}~1020 γυμνὸν ὄνθ᾽ οὕτως. ὁρῶ γὰρ ὡς καταγέλαστος εἶ ; Ar.~\textit{Ran.}~480 ὦ καταγέλαστ᾽ οὔκουν ἀναστήσει ταχὺ ; Ar.~\textit{Eccl.}~125–127 ὡς καὶ καταγέλαστον τὸ πρᾶγμα φαίνεται. — πῶς καταγέλαστον; — ὥσπερ εἴ τις σηπίαις / πώγωνα περιδήσειεν ἐσταθευμέναις ; Hsch.~γ~311 γελοῖος· ἱλαρός, γελοιαστὴς καὶ ὁ καταγέλαστος ἄνθρωπος καὶ λοίδορος ; Ar.~\textit{Ach.}~1196–1197 Δικαιόπολις εἴ μ᾽ ἴδοι τετρωμένον / κᾆτ᾽ ἐγχάνοι ταῖς ἐμαῖς τύχαισιν ; Ar.~\textit{Nub.}~849 ἄμφω ταὐτό; καταγέλαστος εἶ ; Ar.~\textit{Thesm.}~941–942 ἵνα μὴ ᾽ν κροκωτοῖς καὶ μίτραις γέρων ἀνὴρ / γέλωτα παρέχω τοῖς κόραξιν ἑστιῶν ; Pl.~\textit{Ap.}~28d ἵνα μὴ ἐνθάδε μένω καταγέλαστος παρὰ νηυσὶ κορωνίσιν ; Pl.~\textit{Cri.}~53a καὶ οὐ καταγέλαστός γε ἔσῃ ἐκ τῆς πόλεως ἐξελθών ; Synag.~γ~19 γελοιος· ὁ καταγελαστότατος προπερισπωμένως· ὀξυτόνως δὲ ὁ γελωτοποιός ; \textit{idem} Suid.~γ~118 (Γέλοιος ; προπαροξυτόνως pro προπερισπωμένως ; Γελοῖος δὲ pro ὀξυτόνως δὲ) ; Suid.~γ~118 (Γελοιότερον Μελιτίδου) καὶ ἑτέρα παροιμία· Γελοιότερον Μελιτίδου, ἐπὶ τῶν ἐπὶ μωρίᾳ διαβεβλημένων ; Ar.~\textit{Ach.}~76 ἆρ᾽ αἰσθάνει τὸν κατάγελων τῶν πρέσβεων ; Ar.~\textit{Ach.}~1081 οἴμοι κακοδαίμων καταγελᾷς ἤδη σύ μου ; Ar.~\textit{Eq.}~713 ἐγὼ δ᾽ ἐκείνου καταγελῶ γ᾽ ὅσον θέλω ; Ar.~\textit{Nub.}~1238 οἴμ᾽ ὡς καταγελᾷς ; Ar.~\textit{Vesp.}~515–516 καταγελώμενος μὲν οὖν / οὐκ ἐπαΐεις ὑπ᾽ ἀνδρῶν, οὓς σὺ μόνον οὐ προσκυνεῖς ; Ar.~\textit{Vesp.}~1406 καὶ καταγελᾷς μου ; Ar.~\textit{Av.}~1407 καταγελᾷς μου, δῆλος εἶ ; Ar.~\textit{Plut.}~880 οἴμοι τάλας· μῶν καὶ σὺ μετέχων καταγελᾷς ; Hsch.~γ~310 γελοίϊον· γέλοιον, γέλωτος ἄξιον (Β 215) ; Ammon.~\textit{Diff.}~119 γέλοιος καὶ γελοῖος διαφέρει. γέλοιος μὲν ὁ καταγέλαστος, γελοῖος δὲ ὁ γελωτοποιός ; Suid.~ο~397 Ὄνου παρακύψεως: ἐπὶ τῶν καταγελάστως συκοφαντουμένων εἴρηται ἡ παροιμία}
\lp{226 τεμνόμενος]}{Aesch.~\textit{Ag.}~849 ἤτοι κέαντες ἢ τεμόντες εὐφρόνως ; Mart.~11.84.1–4 Qui nondum Stygias descendere quaerit ad umbras, / Tonsorem fugiat, si sapit, Antiochum. / Alba minus saevis lacerantur bracchia cultris, / Cum furit ad Phrygios enthea turba modos ; Pl.~\textit{Grg.}~480c παρέχειν μύσαντα εὖ καὶ ἀνδρείως ὥσπερ τέμνειν καὶ κάειν ἰατρῷ ; X.~\textit{An.}~5.8.18 καὶ γὰρ οἱ ἰατροὶ καίουσι καὶ τέμνουσιν ἐπ᾽ ἀγαθῷ ; Hp.~\textit{Aph.}~7.87 Ὁκόσα φάρμακα οὐκ ἰῆται, σίδηρος ἰῆται· ὅσα σίδηρος οὐκ ἰῆται, πῦρ ἰῆται ; Soph.~\textit{Aj.}~581–582 οὐ πρὸς ἰατροῦ σοφοῦ / θρηνεῖν ἐπῳδὰς πρὸς τομῶντι πήματι ; Mart.~11.58.5–6 Quid si me tonsor, cum stricta novacula supra est, / Tunc libertatem divitiasque roget? ; Heraclit.~fr.~58~DK οἱ γοῦν ἰατροί, φησὶν ὁ Ἡ., τέμνοντες, καίοντες πάντη, βασανίζοντες κακῶς τοὺς ἀρρωστοῦντας, ἐπαιτέονται μηδὲν ἄξιοι μισθὸν λαμβάνειν ; Charito~3.4.12 καόμενος δὲ καὶ τεμνόμενος ἀντεῖχεν ἐπὶ πλέον ; Artem.~5.94 θαρρῶν τέμνου, θεραπευθήσει τεμνόμενος}
\lp{227 ἡμίκραιραν]}{\textit{Test.} Ael.Dion.~η~9 Erbse ἡμίκραιρα (Ar. Thesm. 227)· τὸ ἥμισυ τῆς κεφαλῆς ; \textit{idem} Suid.~η~334 (Ἡμίκραιρα ; add. τὴν ἡμίκραιραν τὴν ἑτέραν ψιλὴν ἔχων), Et. Gen. (AB) p.~150 Miller (Ἡμίκραιρα καὶ ἡμίκραιρον· λέγεται οὐδετέρως ; add. Ἀριστοφάνης Θεσμοφοριαζούσαις. Οὕτως Ὦρος ὁ Μιλήσιος), Hdn.~\textit{Pass.}~266 (ἡμίκραιρον), EM~430.29 (ἡμίκραιρον) ; Hsch.~η~496 ἡμίκραιρα· ἡμικεφάλαιον. σιαγόνος τὸ ἥμισυ (Ar. Thesm. 227) ; Eust.~710.50 ὅθεν παρὰ τῷ κωμικῷ ἐψίλωταί τις τὴν ἑτέραν ἡμίκραιραν ἤτοι τὸ ἡμίκρανον ; Eust.~1127.33–35 ἐξ ἧς οἱ Ἀττικοὶ τὸ τῆς κεφαλῆς ἥμισυ ἡμίκραιραν φασίν, ὡς εὕρηται καὶ παρὰ τῷ κωμικῷ ; Ameips.~fr.~7~K.-A. =~7~Kock (Konnos, Ath.~9.6,~368e) δίδοται μάλισθ᾽ ἱερώσυνα, / κωλῆ, τὸ πλευρόν, ἡμίκραιρ᾽ ἀριστερά ; Crobyl.~fr.~6.1~K.-A. (=~6~Kock ; Ψευδυποβολιμαῖος ; Ath.~9.384c) εἰσῆλθεν ἡμίκραιρα τακερὰ δέλφακος ; Σ Il.~18.3 κραίραν γὰρ λέγουσι τὴν κεφαλὴν, καὶ οἱ Ἀττικοὶ τὸ τῆς κεφαλῆς ἥμισυ ἡμίκραιραν ; Phryn.~\textit{Ecl.}~301 Fischer (303 Rutherford ; p.~328 Lobeck) Ἡμικεφάλαιον μὴ λέγε, ἀλλὰ ἡμίκρανον ; Hsch.~κ~3931 κραῖρα· ἡ κεφαλή, καὶ ἀκροστόλιον ; Hsch.~κ~3932 κραῖραι· στόλοι νεῶν. μέτωπα, κεφαλαί ; Poll.~6.160–161 ἡμίκρανον, ἡμίκραιρα ; Hsch.~ο~1192 ὀρθοκραιράων· ἐπὶ μὲν τῶν νεῶν, τῶν ὀρθοπρύμνων (Σ 3)· ἐπὶ δὲ τῶν βοῶν, τῶν ὀρθοκεράτων (Θ 231). εἴρηται δὲ παρὰ τὸ κάρα}
\lp{227 τὴν ἡμίκραιραν τὴν ἑτέραν ψιλὴν ἔχων]}{Ar.~\textit{Thesm.}~583 ἕως ἂν οὕτως τὰς γνάθους ψιλὰς ἔχῃς ; Hermipp.~fr.~13~K.-A. (=~14~Kock ; Δημόται ; Poll.~9.70) οἴμοι, τί δράσω σύμβολον κεκαρμένος ; Hor.~\textit{Epist.}~1.1.94–95 Si curatus inaequali tonsore capillos / occurri, rides ; Hdt.~2.121δ.6 τῶν φυλάκων ἐπὶ λύμῃ πάντων ξυρῆσαι τὰς δεξιὰς παρηίδας ; Hdt.~4.191.1 οἳ τὰ ἐπὶ δεξιὰ τῶν κεφαλέων κομόωσι, τὰ δ᾽ ἐπ᾽ ἀριστερὰ κείρουσι ; Plu.~\textit{Ages.}~30.3 καὶ ξυρῶνται μέρος τῆς ὑπήνης, μέρος δὲ τρέφουσι ; Plu.~\textit{Dem.}~7.3 ξυρούμενον τῆς κεφαλῆς θάτερον μέρος ὑπὲρ τοῦ μηδὲ βουλομένῳ πάνυ προελθεῖν ἐνδέχεσθαι ; D.L.~6.33 εἰσελθών ποτε ἡμιξύρητος εἰς νέων συμπόσιον ; Σ Ar.~\textit{Thesm.}~227 τὴν μίαν κομῶν γνάθον, τὴν ἑτέραν ἐψιλωμένος ; Ar.~\textit{Thesm.}~232 οἴμοι κακοδαίμων, ψιλὸς αὖ στρατεύσομαι ; Catull.~59.5 ab semiraso tunderetur ustore ; Hdt.~4.175.1 Μάκαι, οἳ λόφους κείρονται, τὸ μὲν μέσον τῶν τριχῶν ἀνιέντες αὔξεσθαι, τὰ δὲ ἔνθεν καὶ ἔνθεν κείροντες ἐν χροΐ ; Hdt.~4.180.1 καὶ οἱ μὲν Μάχλυες τὰ ὀπίσω κομῶσι τῆς κεφαλῆς, οἱ δὲ Αὐσέες τὰ ἔμπροσθε ; Pl.~\textit{Smp.}~191d ἕκαστος οὖν ἡμῶν ἐστιν ἀνθρώπου σύμβολον, ἅτε τετμημένος ὥσπερ αἱ ψῆτται ; LXX~2 Regn.~10.4 καὶ ἐξύρησεν τοὺς πώγωνας αὐτῶν ; Luc.~\textit{Peregr.}~17 ξυρόμενος μὲν τῆς κεφαλῆς τὸ ἥμισυ ; Soph.~fr.~473~R. (Οἰνόμαος ; Ath.~9.410c) Σκυθιστὶ χειρόμακτρον ἐκκεκαρμένος ; Hdt.~2.65.4 ξυρῶντες τῶν παιδίων ἢ πᾶσαν τὴν κεφαλὴν ἢ τὸ ἥμισυ ἢ τὸ τρίτον μέρος τῆς κεφαλῆς}
\lp{227 τὴν ἑτέραν]}{Ar.~\textit{Ran.}~1415 ἐὰν δὲ κρίνω; — τὸν ἕτερον λαβὼν ἄπει}
\lp{227 ψιλήν]}{Hsch.~ψ~192 ψιλὰ τὰ γένεια· σπανοπώγων ; Hsch.~ε~3456 «ἐν χρῷ κουρά»· ἡ ψιλή, καὶ πρὸς αὐτῷ τῷ χρωτί ; Hsch.~ψ~195 ψίλακα· ψιλόν, λεῖον. πτερὸν ἢ πτενόν (Ar. fr. 922 K.–A.) ; Hsch.~ψ~202 ψιλοκόρρης· φαλακρός (com. ad. fr. 441 K.–A.) ; Hsch.~ψ~203 ψιλόν· γυμνόν. ἔρημον. εὐτελές. ἄτριχον. καὶ πᾶν ἀποδεδαρμένον (ν 437) ; Phryn.~\textit{Ecl.}~38 Fischer (41 Rutherford) Ἐν χρῷ κουρίας φαθί, καὶ μὴ ψιλόκουρος ; X.~\textit{An.}~1.8.6 Κῦρος δὲ ψιλὴν ἔχων τὴν κεφαλὴν εἰς τὴν μάχην καθίστατο ; Hsch.~ψ~206 ψιλούς· γυμνούς. σφενδονιστάς. τοξότας (Thuc. 1, 60, 1 . .) ; Moer.~ψ~4 Hansen (p.~214,~2~Bekker) ψίλαξ Ἀριστοφάνης, ψιλὸς καὶ λεῖος Ἕλληνες}
\lp{228 ὀλίγον μέλει μοι]}{Ar.~\textit{Eq.}~1195 ὀλίγον μοι μέλει ; \textit{idem} Ar.~\textit{Nub.}~1142 (ὀλίγον γάρ μοι μέλει), Ar.~\textit{Vesp.}~1411 (praem. ἔπειθ᾽ ὁ Λᾶσος εἶπεν), Ar.~\textit{Vesp.}~1446 (praem. Αἴσωπον οἱ Δελφοί ποτ᾽ —), Ar.~\textit{Av.}~1636, Ar.~\textit{Lys.}~248 (ὀλίγον αὐτῶν μοι μέλει), Ar.~\textit{Ran.}~1136 (praem. ὁρᾷς ὅτι ληρεῖς; — ἀλλ᾽ ; ὀλίγον γέ μοι μέλει) ; Ar.~\textit{Ran.}~655 οὐδέν μοι μέλει ; Hdt.~6.129.4 ὁ δὲ Ἱπποκλείδης ὑπολαβὼν εἶπε οὐ φροντὶς Ἱπποκλείδῃ ; Hsch.~μ~674 μέλει· φροντίζει ; Hsch.~ο~550 ὀλίγον· μικρὸν ἢ οὐδέν (Ε 800) ; Hsch.~ο~1920 οὐ φροντὶς Ἱπποκλείδῃ· παροιμία, ἧς μέμνηται καὶ Ἕρμιππος Δημόταις (fr. 16 K.–A.)· [μὴ πρὸς ὀργὴν ὁμιλεῖν. καὶ] ἀντὶ τοῦ ἐμοί ; \textit{idem} Phot.~ο~697 (om. καὶ ; ἐν Δημόταις ; om. μὴ πρὸς ὀργὴν … ἐμοί), Suid.~ο~978 (om. καὶ ; ἐν Δημόταις ; om. μὴ πρὸς ὀργὴν … ἐμοί), Zen.~5.31 (lemma tantum), Diogenian.~7.21 (lemma tantum) ; Erasm.~\textit{Adag.}~1.10.12 (n°~912) Allusit huc Aristophanes in Vespis, saepius repetens illud, Ὀλίγον μοι μέλει, i. e. Leviter mihi cura est ; Ar.~\textit{Lys.}~895–896 ὀλίγον αὐτῶν μοι μέλει. — ὀλίγον μέλει σοι τῆς κρόκης φορουμένης / ὑπὸ τῶν ἀλεκτρυόνων ; Eur.~\textit{Cycl.}~331 χιόνος οὐδέν μοι μέλει ; Aristid.~\textit{Or.}~51 (II, p.~573 Dindorf) ἕτεροι δ᾽ εἰσὶν οἷς λόγων ὀλίγον μέλει ; Σ Luc.~\textit{Apol.}~15 οὐ φροντὶς Ἱπποκλείδῃ] παροιμία ἐπὶ τῶν ἄγαν καταφρονούντων τινὸς καὶ ἀπεριμερίμνως διακειμένων ; Hsch.~ο~556 ὀλιγωρεῖ· ῥαθυμεῖ, ἀμελεί ; Hsch.~ο~557 ⟨ὀλιγωρήσας⟩· [ἢ] ὀλίγην ἔχων φροντίδα ; Plu.~\textit{De\,Herodoti\,malign.}~33 (867b) καθάπερ Ἱπποκλείδης ὁ τοῖς σκέλεσι χειρονομῶν ἐπὶ τῆς τραπέζης, εἰπεῖν ἂν ἐξορχούμενος τὴν ἀλήθειαν· οὐ φροντὶς Ἡροδότῳ ; Luc.~\textit{Apol.}~15 ἱκανὸν ἂν εἴη μοι τό οὐ φροντὶς Ἱπποκλείδῃ ; Luc.~\textit{Herc.}~8 ἐσιδών με ὑποπόλιον τὸ γένειον … καὶ ὁ Ἱπποκλείδης οὐ φροντιεῖ ; Julian.~\textit{Or.}~6.2 (182b) οὐ φροντὶς Ἱπποκλείδῃ· καὶ γὰρ οὐδὲ τῶν κυνιδίων ἡμῖν μέλει τὰ τοιαῦτα πλημμελούντων ; Plu.~\textit{Amat.}~6 (752d) καὶ ὁ Πεισίας ἐμοὶ μὲν εἶπεν ὀλίγον μέλει τοῦ λόγου ; Luc.~\textit{Contempl.}~12 τῷ δὲ θεῷ ὀλίγον μέλει τῶν σῶν χρυσοποιιῶν ; Luc.~\textit{Hermot.}~18 ἢ τούτων ὀλίγον σοι μέλει, ἄχρι ἂν εὐσταλὴς ἡ ἀναβολὴ καὶ ὁ πώγων βαθὺς καὶ ἐν χρῷ ἡ κουρά ; Σ Luc.~\textit{Herc.}~8 παροιμία οὐ φροντὶς Ἱπποκλείδῃ ἐπὶ τῶν μὴ πάνυ σπουδαίων ἡμῖν λεγομένη}
\lp{228 μηδαμῶς πρὸς τῶν θεῶν]}{Ar.~\textit{Pax}~322 μηδαμῶς πρὸς τῶν θεῶν ; \textit{idem} Ar.~\textit{Eccl.}~562 (add. τουτὶ ποιήσῃς) ; Eur.~\textit{Alc.}~275 μὴ πρός ⟨σε⟩ θεῶν τλῇς με προδοῦναι ; Ar.~\textit{Vesp.}~919 πρὸς τῶν θεῶν μὴ προκαταγίγνωσκ᾽ ὦ πάτερ ; Ar.~\textit{Pax}~376–377 μὴ πρὸς τῶν θεῶν / ἡμῶν κατείπῃς, ἀντιβολῶ σε δέσποτα ; Ar.~\textit{Eccl.}~1054 μηδαμῶς με περιίδῃς ; Hsch.~μ~1209 μή με· μηδαμῶς με ; Ar.~\textit{Ach.}~590 οἴμ᾽ ὡς τεθνήξεις. — μηδαμῶς ὦ Λάμαχε ; Ar.~\textit{Vesp.}~1164 φέρε καὶ τὸν ἕτερον. — μηδαμῶς τοῦτόν γ᾽ ; Ar.~\textit{Pax}~926 βοΐ; μηδαμῶς ; Ar.~\textit{Av.}~145 οἴμοι μηδαμῶς ; Ar.~\textit{Thesm.}~172 πῶς πρὸς τῶν θεῶν ; Ar.~\textit{Thesm.}~714 τοῦτο μέντοι μὴ γένοιτο μηδαμῶς, ἀπεύχομαι ; Hsch.~α~3732 δηλοῖ δὲ καὶ τὸ μηδαμῶς, κατὰ μηδένα τρόπον ; Hsch.~μ~1140 μηδαμῶς· μὴ δῆτα ; Hsch.~μ~1152 μὴ δὴ οὖν· μηδαμῶς οὖν ; Hsch.~μ~1153 μὴ δῆτα· μηδαμῶς ; Hsch.~μ~1266 μή σε· μηδαμῶς σε}
\lp{229 προδῷς με]}{Ar.~\textit{Thesm.}~210 μὴ σαυτὸν προδῷς ; Eur.~\textit{Hipp.}~1456 μή νυν προδῷς με, τέκνον, ἀλλὰ καρτέρει ; Soph.~\textit{Phil.}~757 μή με ταρβήσας προδῷς ; Eur.~\textit{IT}~706 καὶ μὴ προδῷς μου τὴν κασιγνήτην ποτέ ; Carm. pop. ap. Ath.~15.53 (697b–c) ὦ τί πάσχεις; μὴ προδῷς ἄμμ᾽, ἱκετεύω ; Hsch.~π~3432 προέσθαι· προδοῦναι}
\lp{229 χώρει δεῦρο]}{Ar.~\textit{Nub.}~509 χώρει· τί κυπτάζεις ἔχων περὶ τὴν θύραν ; Ar.~\textit{Av.}~1186 χώρει δεῦρο πᾶς ὑπηρέτης ; Ar.~\textit{Lys.}~738 χώρει πάλιν δεῦρ᾽ ; Hsch.~δ~731 δεῦρο· … ἢ παρακελευστικῶς· «δεῦρό νυν, ἢ τρίποδος περιδώμεθα» (Ψ 485) ; Ar.~\textit{Nub.}~889 χώρει δευρί, δεῖξον σαυτὸν ; Ar.~\textit{Pax}~301 δεῦρο πᾶς χώρει προθύμως εὐθὺ τῆς σωτηρίας ; Ar.~\textit{Eccl.}~730 χώρει σὺ δεῦρο κιναχύρα καλὴ καλῶς ; Theopomp. Com. fr.~33.1~K.-A. (=~32.1~Kock ; Νεμέα ; Ath.~11.470f) χώρει σὺ δεῦρο, Θηρικλέους πιστὸν τέκνον ; Hsch.~χ~869 χώρει· πορεύου ; Hsch.~χ~870 χωρεῖν· ἀπελθεῖν (Soph. Tr. 496. Eur. Suppl. 461 . .)}
\lp{229 κακοδαίμων ἐγώ]}{Ar.~\textit{Ach.}~1094 κακοδαίμων ἐγώ ; \textit{idem} Ar.~\textit{Nub.}~698, Ar.~\textit{Vesp.}~1166, Ar.~\textit{Thesm.}~604, Ar.~\textit{Thesm.}~650 ; Eur.~\textit{Hipp.}~1362 τὸν κακοδαίμονα καὶ κατάρατον ; Hsch.~κ~327 κακοδαίμων· ὁ θεῷ ἀπεχθής (Dem. 19, 115) ; Ar.~\textit{Thesm.}~1225 κακόδαιμον, ἀλλὰ τρέξι ; Pherecr.~fr.~118~K.-A. (=~117~Kock ; Μυρμηκάνθρωποι ; Phot.~α~525) οἴμοι κακοδαίμων, αἰγὶς αἰγὶς ἔρχεται ; Eur.~\textit{Phoen.}~1615 εἶἑν· τί δράσω δῆθ᾽ ὁ δυσδαίμων ἐγώ ; Antipho~5.43 οὐ γὰρ δήπου οὕτω κακοδαίμων ἐγώ ; Pl.~\textit{Men.}~78a ἔστιν οὖν ὅστις βούλεται ἄθλιος καὶ κακοδαίμων εἶναι ; Pl.~\textit{R.}~4.440a ἰδοὺ ὑμῖν, ἔφη, ὦ κακοδαίμονες, ἐμπλήσθητε τοῦ καλοῦ θεάματος ; Pl.~\textit{Smp.}~173d καὶ ἴσως αὖ ὑμεῖς ἐμὲ ἡγεῖσθε κακοδαίμονα εἶναι ; Hsch.~β~259 βαρύστονος· κακοδαίμων ; Hsch.~δ~74 δαιμόνιε· μακάριε (Β 190 . .). κακόδαιμον ; Hsch.~κ~333 κακοί· ἄνανδροι, δειλοί (Λ 408 . .), δυστυχεῖς, κακοδαίμονες}
\lp{230 ἔχ᾽ ἀτρέμα σαυτόν]}{Ar.~\textit{Nub.}~261 ἀλλ᾽ ἔχ᾽ ἀτρεμί ; Ar.~\textit{Nub.}~743 ἔχ᾽ ἀτρέμα ; \textit{idem} Ar.~\textit{Av.}~1244, A.R.~3.952 (ἔχ᾽ ἀτρέμας), Luc.~\textit{Contempl.}~7, Max.Tyr.~11.2 (Ἔχε ἀτρέμας) ; Ar.~\textit{Av.}~1200 ἔχ᾽ ἀτρέμας· αὐτοῦ στῆθ᾽· ἐπίσχες τοῦ δρόμου ; \textit{idem} Suid.~ε~2807 (praem. αὕτη σύ, ποῖ, ποῖ, ποῖ, ποῖ πέτῃ: μέν᾽ ἥσυχος ; om. αὐτοῦ στῆθ᾽) ; Ar.~\textit{Ran.}~339 οὔκουν ἀτρέμ᾽ ἕξεις, ἤν τι καὶ χορδῆς λάβῃς ; Eur.~\textit{Or.}~258 μέν᾽, ὦ ταλαίπωρ᾽, ἀτρέμα σοῖς ἐν δεμνίοις ; Hdt.~9.54.1 εἶχον ἀτρέμας σφέας αὐτοὺς ἵνα ἐτάχθησαν ; Hld.~5.2 ἀλλὰ σὺ μὲν ἀτρέμας ἔχε σαυτόν ; Synag.~α~1066 ἀτρέμα· ἠρέμα ; Ar.~\textit{Nub.}~1244 ἔχε νυν ἥσυχος ; Ar.~\textit{Av.}~1572 ἕξεις ἀτρέμας ; Ar.~\textit{Plut.}~127 ἆ μὴ λέγ᾽ ὦ πόνηρε ταῦτ᾽. — ἔχ᾽ ἥσυχος ; Pherecr.~fr.~6.3~K.-A. (=~6.3~Kock ; Ἄγριοι ; Σ Ar.~\textit{Av.}~858) Β. ἔχ᾽ ἀτρέμ᾽, ἐγᾦδα, Χαῖρις ; Strattis fr.~55~K.-A. (=~52~Kock ; Χρύσιππος ; Poll.~10.55) πρόσαγε τὸν πῶλον ἀτρέμα, προσλαβὼν / τὸν ἀγωγέα βραχύτερον ; Eur.~\textit{Ba.}~1072 ἀτρέμα, φυλάσσων μὴ ἀναχαιτίσειέ νιν ; Eur.~\textit{Hipp.}~1358 φεῦ φεῦ· πρὸς θεῶν, ἀτρέμας, δμῶες ; Eur.~\textit{IT}~1226 ἐκποδὼν δ᾽ αὐδῶ πολίταις τοῦδ᾽ ἔχειν μιάσματος ; Eur.~\textit{Or.}~149 κάταγε κάταγε, πρόσιθ᾽ ἀτρέμας, ἀτρέμας ἴθι ; \textit{Il.}~15.318 ὄφρα μὲν αἰγίδα χερσὶν ἔχ᾽ ἀτρέμα Φοῖβος Ἀπόλλων ; Hdt.~3.79.3 ἀλλὰ κατ᾽ οἴκους ἑωυτοὺς οἱ Μάγοι ἔχουσι ; Pl.~\textit{Ep.}~7.348c ἐκποδὼν αὑτὸν ἔσχεν ἀφανῆ ; Σ Eur.~\textit{Or.}~258 μένε ἀτρέμας καὶ ἡσύχως ἔχε ἐν τοῖς δεμνίοις ; Hsch.~α~8144 ἀτρέμας· ἡσύχως, ἡσυχῇ (Β 200 . .) ; Hsch.~α~8146 ἀτρεμῆσαι· ἡσυχάσαι ; Suid.~α~4384 Ἀτρέμας: ἠρέμα. καὶ Ἀτρεμία, ἡσυχία. καὶ Ἀτρεμίζων, ἑδράζων ; Ar.~\textit{Ach.}~861 κατάθου τὺ τὰν γλάχων᾽ ἀτρέμας Ἰσμηνία ; Ar.~fr.~592.15~K.-A. (P. Oxy.~212 ; CGFP~62) (Α.) τί οὖν γένοιτ᾽ ἄν; (Β.) ἔχ᾽ ἀ.[ ; Hsch.~α~8147 ἀτρέμητον· ἀσάλευτον ; Hsch.~α~8148 ἀτρεμία· ἡσυχία ; Hsch.~η~750 ἠρέμα· ἀτρέμα. ἡσύχως. σχολαίως ; Hsch.~η~751 ἠρεμεῖ· ἡσυχάζει ; Synag.~η~105 ἠρέμα· ἀτρέμα ; Phot.~η~235 (Naber I, p.~267) ἠρέμα· ἀτρέμα, πρᾴως, ἡσυχῆ· ἡ γὰρ ἠρεμία ἡσυχία ; Suid.~η~491 Ἠρέμα: ἀτρέμας. καὶ Ἠρεμαίως καὶ Ἠρεμεῖ, ἡσύχως}
\lp{230 ἀτρέμα]}{Ar.~\textit{Eq.}~24 ὥσπερ δεφόμενος νῦν ἀτρέμα πρῶτον λέγε ; Ap.Dysc.~\textit{Adv.}~160.23–161.2 Schneider ὡς γὰρ τῷ ἠρεμῶ τὸ ἠρέμα παράκειται, τῷ σιγῶ τὸ σῖγα, ἀντῶ ἄντα, ⟨οὕτως⟩ ἀτρεμῶ ἀτρέμα καὶ ἀτρέμας ; Synag. (Σᵇ) α~2364 ἀτρέμας = Σ α 1066. — καὶ ἀτρέμα λέγουσιν καὶ ἀτρεμεῖν ; Ar.~\textit{Nub.}~390 ἀτρέμας πρῶτον παππὰξ παππάξ, κἄπειτ᾽ ἐπάγει παπαπαππάξ ; \textit{Od.}~13.92 δὴ τότε γ᾽ ἀτρέμας εὗδε, λελασμένος ὅσσ᾽ ἐπεπόνθει ; Bacchyl.~5.7 ἀτρέμ᾽ ἀμπαύσας μεριμνᾶν ; Antipho~3.3.10 διὰ τὸ ὑπὸ τὴν φορὰν τοῦ ἀκοντίου ὑπελθεῖν καὶ μὴ ἀτρέμας ἑστάναι … ἀλλ᾽ εἴπερ τούτου μὴ ἀκοντίζοντος ἀλλ᾽ ἀτρέμα ἑστῶτος ἀπέθανεν ὁ παῖς ; Antipho~3.4.7 ἐξῆν γὰρ αὐτῷ ἀτρέμα ἑστάναι ; Hsch.~α~308 ἀγανηδά· ἀτρέμας ; Hsch.~α~2413 ἀκιρῶς· εὐλαβῶς. ἀτρέμας ; Hsch.~α~8149 ἀτρεμέων· ἡσυχάζων (Greg. Naz. c. 1, 1, 36, 26 [37, 520, 2 M.]) ; Hsch.~ε~6910 εὐκαλεῖ· ἀτρεμίζει ; Hsch.~ε~7201 εὐσταλεῖ· ἀτρεμεῖ ; Ap.Dysc.~\textit{Adv.}~152.1–5 Schneider ἠρεμῶ ἠρέμα, μίγω μίγα, ἀντῶ ἄντα … σιγῶ σῖγα, ἀτρεμῶ ἀτρέμα ; Suid.~η~493 Ἠρεμία: ἀτρεμία τις οὖσα, κατὰ στέρησιν τοῦ τρέμειν, οἷον ἡ ἀκινησία ; Zonar. col.~332~Tittmann (s.v.~ἀτρέμας) Ἀτρέμας. ἥσυχος, ἄφοβος ; Diogenian.~3.9 Ἀτρέμας βοῦς: ἐπὶ τῶν βραδέως συντελούντων ; \textit{idem} Apostol.~4.24 ; Erasm.~\textit{Adag.}~2.1.3 (n°~1003) Placide bos. Ἀτρέμας βοῦς, id est Lente bos, subaudiendum incedit aut simile quippiam. Congruet in illos, qui placide ac paulatim citraque tumultum, sed tamen assiduitate rem conficiunt}
\lp{230 κἀνάκυπτε]}{Ar.~\textit{Eq.}~1354 οὗτος τί κύπτεις; οὐχὶ κατὰ χώραν μενεῖς ; Ar.~\textit{Thesm.}~236 κἀγκύψας ἔχε ; Ar.~\textit{Thesm.}~930 οὗτος τί κύπτεις; δῆσον αὐτὸν εἰσάγων ; Ar.~\textit{Thesm.}~1187a ἀνακύπτι καὶ παρακύπτι ; Ar.~fr.~697~K.-A. (=~679~Kock ; Poll.~7.143) ὥστ᾽ ἀνακύπτων καὶ κατακύπτων τοῦ σχήματος εἵνεκα τοῦδε / κηλωνείου τοῖς κηπουροῖς ; Eur.~\textit{Cycl.}~211–212 βλέπετ᾽ ἄνω καὶ μὴ κάτω. — ἰδού, πρὸς αὐτὸν τὸν Δί᾽ ἀνακεκύφαμεν ; Pl.~\textit{R.}~7.529b ἐν ὀροφῇ ποικίλματα θεώμενος ἀνακύπτων ; Ar.~\textit{Av.}~146–147 ἵν᾽ ἀνακύψεται / κλητῆρ᾽ ἄγουσ᾽ ἕωθεν ἡ Σαλαμινία ; Ar.~\textit{Ran.}~1068 παρὰ τοὺς ἰχθῦς ἀνέκυψεν ; X.~\textit{Eq.}~7.10 ἢν δὲ μᾶλλον ἀνακεκυφώς, κατωτέρω ; Ev.Luc.~21.28 ἀνακύψατε καὶ ἐπάρατε τὰς κεφαλὰς ὑμῶν ; Thphr.~\textit{Char.}~11.3 καὶ ὅταν σιωπήσῃ τὸ θέατρον, ἀνακύψας ἐρυγεῖν ; Hsch.~α~4480 †ἀναπαντοῦσιν· ἀνακύπτουσιν ; Hsch.~κ~4656 κυπτάζειν· … ἐπικύπτειν. καὶ χρονίζειν ; Hsch.~κ~4657 κύπτει· κάμπτει, ἑαυτὸν κλίνει ; Suid.~α~1926 Ἀνακύπτω· γενικῇ}
\lp{230 ποῖ στρέφει;]}{Ar.~\textit{Thesm.}~610 αὕτη σὺ ποῖ στρέφει; μέν᾽ αὐτοῦ ; Luc.~\textit{Catapl.}~25 ποῖ στρέφῃ; πρόσιθι ; Pl.~\textit{Phdr.}~236e τί δῆτα ἔχων στρέφῃ ; Hsch.~σ~1992 στρέφει· κάμπτει (Ψ 323). λυγίζει. ἀποσείεται. διώκει. ψιλοῖ}
\lp{231 μυμῦ]}{Ar.~\textit{Eq.}~10 μυμῦ μυμῦ μυμῦ μυμῦ μυμῦ μυμῦ ; Erasm.~\textit{Adag.}~1.8.2 (n°~702) Est autem Graecis μύζω, tum sugo, tum naribus vocem reddo. Aristophanes ἐν Θεσμοφοριαζούσαις, Τί μύζεις, πάντα πεποίηται καλῶς. Quum enim Mnesilochus dixisset, μῦ μῦ, id est, mu mu, Euripides respondet, Quid mussitas? peracta sunt cuncta optimè ; Varro L.~7.101 (= Enn. inc.~10~V.) Mussare dictum, quod muti non amplius quam μῦ dicunt; a quo idem dicit id quod minimum est: Neque, ut aiunt, μῦ facere audent ; \textit{idem} Tosi n°~2123 (Neque… μυ facere audent) ; \textit{Test.} Suid.~ε~2807 ἀττατατατατταῖ, ἐπὶ τῶν ἀλγούντων. μῦ μῦ, ἐπὶ τῶν διὰ μυκτήρων ἡχούντω[ν] ; Ar.~\textit{Plut.}~895 ὓ ὗ ὓ ὗ ὓ ὗ ὓ ὗ ὓ ὗ ὓ ὗ ; Hippon.~fr.~123 Deg. (ap. S. E. M.~1.275) μηδέ μοι μῦ / λαλεῖν ; Synag. Σᵇ~α~304~Cunningham ἄγρυκτα καὶ ἄλεκτα πέπονθα· τὸ μὲν ἄγρυκτα ἐστὶν ὥστε μηδὲ γρύξαι διὰ τὴν ὑπερβολὴν τῶν κακῶν· γρύξαι δέ ἐστιν τὸ βραχύτατον φθέγξασθαι, ὃ καὶ ἄναρθρόν ἐστιν, μυγμῷ ἢ στεναγμῷ παραπλήσιον ; \textit{idem} Phot.~α~273 (γρύξαι δέ ἐστι … παραπλήσιον tantum) ; Suid.~γ~461 Γρύ: βραχύ. ἔνιοι δὲ τὸν ἀπὸ τῶν ὀνύχων ῥύπον. Ἀριστοφάνης Πλούτῳ· καὶ ταῦτ᾽ ἀποκρινομένου τοπαράπαν οὐδὲ γρύ. τινὲς παρὰ τὸν γρυλισμὸν, τουτέστι τὴν φωνὴν τῶν χοίρων ; Suid.~μ~1416 Μὺ μῦ, μὺ μῦ, μὺ μῦ, μὺ μῦ, μὺ μῦ, μὺ μῦ: τοῦτο ὡς θρηνοῦντές φασιν ; EM~572.48 (s.v.~Μῦ) Μῦ τὸ στοιχεῖον, ὅτι μυγμόν τινα ἔχει ἡ τούτου ἐκφώνησις. μυγμὸς δέ ἐστιν, ὁ τοῦ μῦ ἦχος διὰ τοῦ μυκτῆρος ἐξερχόμενος ; Et. Gud.~400,~35–37 et~56 Sturz Μῦνδος, ὁ μὴ αὐδῶν· Σοφοκλῆς μύνδος ὁ μὴ αὐδῶν· … τὸ δὲ μυ, ὅτι μυγμόν τινα ἔχει, ἢ τούτου ἐκφώνησις. — Μῦ, ὅτι μεμιγμένον τινὰ ἔχει ἢ τούτου ἐκφώνησις ; Erasm.~\textit{Adag.}~1.8.3 (n°~703) Μηδὲ γρῦ φθέγγεται, id est, Ne gry quidem loquitur: hyperbole item proverbialis apud Graecos, pro eo quod est, ne tantulum quidem ; Lucil.~426 M. non laudare hominem quemquam, neque mu facere, inquam ; Petr.~57.8 Tu lacticulosus, nec mu nec ma argutas ; Ar.~\textit{Plut.}~17 καὶ ταῦτ᾽ ἀποκρινομένῳ τὸ παράπαν οὐδὲ γρῦ ; Hsch.~μ~1871 μύνδος· ἄφωνος. καὶ πόλις Ἀσίας. ἢ ἐνεός (Soph. fr. 1072 R.) ; Phot.~γ~216–218 γρῦ· βραχύ. ἔνιοι δὲ τὸν ἀπὸ τῶν ὀνύχων ῥύπον. — γρῦ· ἀντὶ τοῦ οὐδέν … — γρύξαι· φθέγξασθαι, διακράξαι ; Zen.~5.54 Οὐδὲ τὸ Δίωνος γρῦ ; \textit{idem} App. Prov.~4.91 (Τὸ Δίωνος γρῦ ; add. … αὐτὸς μὲν οὐδὲν ἐφθέγξατο. πρὶν ἂν ἐντὸς τοῦ ἰδίου πυλῶνος γενόμενος καὶ μηδὲν ὀργῆς ὑπομείνας, εἶπεν· οὐδὲ γρῦ) ; Hsch.~γ~938 γρῦ· ὁ ὑπὸ τῷ ὄνυχι ῥύπος. ἤδη δὲ καὶ τὸ ἐλάχιστον. καὶ ἡ γρύτη. καὶ τὸ βραχύ ; Ter.~\textit{An.}~505 Tibi verba censes falso: itaque hercle nihil iam muttire audeo}
\lp{231 τί μύζεις;]}{\textit{Test.} Suid.~μ~1383 Μῦ, μῦ: ἀπὸ τοῦ μῦ παρῆκται τὸ μύζειν, πολλοῖς ἄλλοις ὁμοίως. μύζειν δέ ἐστι τὸ τοῖς μυκτῆρσι ποιὸν ἦχον ἀποτελεῖν. Ἀριστοφάνης· τὶ μύζεις; πάντα πεποίηται καλῶς ; Et. Gen. B ap. Miller, Mélanges, p.~218 Μύζεις· παρὰ τὸ μῦ· Ἀριστοφάνης· μῦ μῦ. — Τί μύζεις; μύζω, ὡς τρύω τρύζω· τί μύζεις; πάντα πεποίηται καλῶς· Ἀριστοφάνης Θεσμοφοριαζούσαις ; EM~593.38 Μύζεις] τί μύζεις; Ἀριστοφάνης. παρὰ τὸ μύω, μύζω, ὡς τρύω, τρύζω. σημαίνει δὲ τὸ θηλάζω ; Eust.~440.23 ἔστι δὲ μύζειν, ὡς καὶ ὁ κωμικὸς δηλοῖ, τὸ τὴν φωνὴν τοῦ μ̄ στοιχείου ἐκφωνεῖν ; Aesch.~\textit{Eum.}~117–118 (μ̔υγμός.) / μύζοιτ᾽ ἄν, ἁνὴρ δ᾽ οἴχεται φεύγων πρόσω ; Aesch.~\textit{PV}~743 σὺ δ᾽ αὖ κέκραγας κἀναμυχθίζῃ ; D.L.~10.118 ὅτε μέντοι στρεβλοῦται, ἔνθα καὶ μύζει καὶ οἰμώζει ; Σ Ar.~\textit{Thesm.}~231 ὅτι τὸ μύζειν παρῆκται ἀπὸ τοῦ μῦ πολλοῖς ἄλλοις ὁμοίως. μύζειν δέ ἐστι τὸ τοῖς μυκτῆρσι ποιὸν ἦχον ἀποτελεῖν ; \textit{Test.} Eust.~723.9 ἔστι δὲ μύζειν τὸ μεμυκότων χειλέων ἦχον ἀποτελεῖν διὰ τῆς ῥινός ; Ar.~\textit{Thesm.}~173 παῦσαι βαΰζων· καὶ γὰρ ἐγὼ τοιοῦτος ἦν ; Ar.~\textit{Thesm.}~1095 ἔτι γὰρ γρύζεις ; Ar.~fr.~475~K.-A. (=~459~Kock ; EM~451,~52) θησομύζειν ἐν τῷ αὐτῷ (Πολυΐδῳ) λέγεται ; Men.~fr.~81~Kock (Erot.~67,~12 Klein) ἔμυξεν· ἐστέναξεν. μέμνηται τῆς λέξεως καὶ Μένανδρος ἐν Ἀσπίδι ; Aesch.~\textit{Ag.}~1308 τί τοῦτ᾽ ἔφευξας; εἴ τι μὴ φρενῶν στύγος ; Aesch.~\textit{Eum.}~189 λευσμός τε, καὶ μύζουσιν οἰκτισμὸν πολὺν ; \textit{Il.}~4.20 ὣς ἔφαθ᾽, αἳ δ᾽ ἐπέμυξαν Ἀθηναίη τε καὶ Ἥρη ; \textit{idem} \textit{Il.}~8.457 ; D.~\textit{H.\,Comp.}~14 ῥοῖζον ἢ σιγμὸν ἢ μυγμὸν ἢ τοιούτων τινῶν ἄλλων ἤχων δηλωτικούς ; Σ Il.~4.20 ἐπέμυξαν] μεμυκόσι τοῖς χείλεσιν ἐπεμύχθησαν. ἠθικῶς δὲ καὶ τὸ γυναικεῖον τῷ σχήματι ἐμφαίνει ; Hsch.~ε~2524 ἔμυζεν· ἔστενεν (Hippocr. epid. 5, 6 [5, 206, 12 L.]). ἀπεθήλαζεν (Archil. fr. 42, 2 W. e ci.) ; Hsch.~ε~4401 ἐπέμυξαν· ἐπεμυκτήρισαν, ἐπεμύχθισαν. ἐξεφαύλισαν (Δ 20) ; Hsch.~μ~1789 μύζει· θηλάζει, λείχει. πιέζει ; \textit{idem} Phot.~μ~574–575 (om. πιέζει ; add. μύζοντες· ἐκπιέζοντες), Suid.~μ~1381 (Μυζεῖ καὶ Μύζει ; om. πιέζει ; add. Ξενοφῶν· ἔδει δέ, εἴ τις διψῴη, λαβόντα εἰς τὸ στόμα μυζεῖν τοὺς καλάμους) ; Hsch.~μ~1791 μύζοντος· μυκτηριζομένου ; Hsch.~μ~1841 μυκτηρίζει· χλευάζει. καταγελᾷ. ἀπὸ τοῦ μύζειν τοῖς μυκτῆρσι (Prov. 11, 12) ; Suid.~μ~1372 Μυγμός: ὁ στεναγμός. καὶ ὁ τοῦ μ ἦχος ; Zonar. col.~1379~Tittmann Μυζῶ. τὸ ῥοφῶ. καὶ μυζεῖς ἀντὶ τοῦ θρηνεῖς. … [Μύζεις. θρηνήσεις.] ; Ar.~\textit{Lys.}~509 οὐ γὰρ γρύζειν εἰᾶθ᾽ ἡμᾶς ; Ar.~\textit{Thesm.}~895 βάυζε τοὐμὸν σῶμα βάλλουσα ψόγῳ ; Ar.~\textit{Thesm.}~1081b οἴμωζ᾽. — οἴμωζ᾽. — ὀτότυζ᾽. — ὀτότυζ᾽ ; Aesch.~\textit{Pers.}~13 νέον δ᾽ ἄνδρα βαΰζει ; Aesch.~fr.~461~R. (dubium) μυχθίζειν ; \textit{Od.}~24.416 μυχμῷ τε στοναχῇ τε δόμων προπάροιθ᾽ Ὀδυσῆος ; Theoc.~20.13 χείλεσι μυχθίζοισα καὶ ὄμμασι λοξὰ βλέποισα ; Hp.~\textit{Carn.}~6 Τὸ δὲ παιδίον ἐν τῇ γαστρὶ συνέχον τὰ χείλεα μύζει ; Hp.~\textit{Morb.}~2.55 καὶ τὰ σπλάγχνα μύζει, … καὶ ὅταν τι φάγῃ, ἐπὶ τοῖσι σπλάγχνοισι μύζει ; Arist.~\textit{HA}~535b32 Ἀφίησι δὲ καὶ ὁ δελφὶς τριγμὸν καὶ μύζει, ὅταν ἐξέλθῃ, ἐν τῷ ἀέρι ; Plu.~\textit{Cor.}~38.2 καὶ μυγμῷ καὶ στεναγμῷ ψόφον ὅμοιον ; Plb.~15.26.8 μυχθίζοντες δὲ καὶ διαψιθυρίζοντες ; Aret.~\textit{SA}~1.5 μυγμὸς δὲ καὶ στεναγμὸς ἡ φωνή ; Σ Il.~8.457 ἐπέμυξαν〉] μεμυκόσι τοῖς χείλεσι ἐγόγγυσαν ; Hsch.~γ~940 γρύζειν· φθέγγεσθαι, λέγειν (Ar. Plut. 454) ; Hsch.~γ~952 γρύξαι· ἠρέμα κράξαι, ἢ ἠρέμα φθέγξασθαι (Ios. 10, 21), ἢ ὑλακτῆσαι (Exod. 11, 7) ; Hsch.~ε~4402 ἐπεμύξατο· ἐπεστέναξεν, ἐπεγόγγυσεν ; Hsch.~ε~4990 ἐπίμυξις· στεναγμός ; Hsch.~μ~1792 μύζουσι· θηλάζουσιν ; Hsch.~μ~2002 μυχθίζουσι· μυκτηρίζουσι, χλευάζουσι ; Hsch.~μ~2003 μυχθισμός· στεναγμός (Eur. Rhes. 789) ; Hsch.~μ~2007 μυχμός· ἰδίωμα ἤχου (ω 416?) ; Hdn.~\textit{Pros.\,cath.} (Lentz, I, μυζῶ) = Antim. fr.~90 καὶ τὸ μυζῶ, ἐξ οὗ τὸ μεμυζότε μυδαλέω τε ; Hdn.~\textit{Περὶ\,παθῶν}~360a μύω μύζω, τρύω τρύζω}
\lp{231 πάντα πεποίηται καλῶς]}{\textit{Test.} Apostol.~16.60c Τί μύζεις; πάντα πεποίηται καλῶς. Ἀριστοφάνους ; Ar.~\textit{Lys.}~1273 ἄγε νυν ἐπειδὴ τἄλλα πεποίηται καλῶς ; Eur.~\textit{Hec.}~875 τὰ δ᾽ ἄλλα — θάρσει — πάντ᾽ ἐγὼ θήσω καλῶς}
\lp{232 οἴμοι κακοδαίμων]}{Ar.~\textit{Thesm.}~237 οἴμοι κακοδαίμων δελφάκιον γενήσομαι ; Ar.~\textit{Eq.}~1206 οἴμοι κακοδαίμων, ὑπεραναιδευθήσομαι ; Ar.~\textit{Eq.}~1243 οἴμοι κακοδαίμων· οὐκέτ᾽ οὐδέν εἰμ᾽ ἐγώ ; Ar.~\textit{Nub.}~504 οἴμοι κακοδαίμων ἡμιθνὴς γενήσομαι ; Ar.~\textit{Nub.}~1324 οἴμοι κακοδαίμων τῆς κεφαλῆς καὶ τῆς γνάθου ; Ar.~\textit{Ran.}~33 οἴμοι κακοδαίμων· τί γὰρ ἐγὼ οὐκ ἐναυμάχουν ; Ar.~fr.~339,~1~K.-A. (=~324~Kock) οἴμοι κακοδαίμων τῆς τόθ᾽ ἡμέρας, ὅτε / εἶπέν μ᾽ ὁ κῆρυξ· οὗτος ἀλφάνει ; Ar.~\textit{Ach.}~105 οἴμοι κακοδαίμων ὡς σαφῶς ; Ar.~\textit{Eq.}~234 οἴμοι κακοδαίμων ὁ Παφλαγὼν ἐξέρχεται ; Ar.~\textit{Lys.}~845 οἴμοι κακοδαίμων, οἷος ὁ σπασμός μ᾽ ἔχει ; Ar.~\textit{Thesm.}~1004 οἴμοι κακοδαίμων, μᾶλλον ἐπικρούεις σύ γε ; Pherecr.~fr.~118~K.-A. (=~117~Kock) οἴμοι κακοδαίμων, αἰγὶς αἰγὶς ἔρχεται}
\lp{232 κακοδαίμων]}{Ar.~\textit{Ach.}~1094 κακοδαίμων ἐγώ ; \textit{idem} Ar.~\textit{Nub.}~698, Ar.~\textit{Vesp.}~1166, Ar.~\textit{Thesm.}~229, Ar.~\textit{Thesm.}~604, Ar.~\textit{Thesm.}~650 ; Eur.~\textit{Hipp.}~1361–1362 σύντονα δ᾽ ἕλκετε / τὸν κακοδαίμονα καὶ κατάρατον ; Phot.~κ~90 κακοδαίμων· ὁ θεῷ †ἐπαχθής†· καὶ ὁ τὴν ψυχὴν ἔχων ἐμπαθῆ}
\lp{232 ψιλὸς αὖ στρατεύσομαι]}{\textit{Test.} Suid.~ψ~105 Ψιλός: τάξις ἐστὶν ἀγεννεστάτη, ἥ τις προβάλλεται ἔμπροσθεν τοῦ στρατοπέδου γυμνῶν ἀόπλων· καὶ ἐὰν ταραχὴ γένηται, ἐκείνη πρώτη ἀναιρεῖται. οἴμοι κακοδαίμων· ψιλὸς αὖ στρατεύσομαι. τοῦτο δὲ εἴρηκε, διότι ἐψίλωτο ; Ar.~\textit{Thesm.}~227 τὴν ἡμίκραιραν τὴν ἑτέραν ψιλὴν ἔχων ; Soph.~\textit{Aj.}~1122–1123 μέγ᾽ ἄν τι κομπάσειας, ἀσπίδ᾽ εἰ λάβοις. / κἂν ψιλὸς ἀρκέσαιμι σοί γ᾽ ὡπλισμένῳ ; Soph.~\textit{Phil.}~952–953 ὦ σχῆμα πέτρας δίπυλον, αὖθις αὖ πάλιν / εἴσειμι πρὸς σὲ ψιλός, οὐκ ἔχων τροφήν ; Th.~4.94.1 ψιλοὶ δὲ ἐκ παρασκευῆς μὲν ὡπλισμένοι οὔτε τότε παρῆσαν οὔτε ἐγένοντο τῇ πόλει ; Arist.~\textit{Ath.}~33.2 πολέμου τε καθεστῶτος καὶ ἐκ τῶν ὅπλων τῆς πολιτείας οὔσης ; Σ Ar.~\textit{Thesm.}~232 ψιλὴ τάξις ἐστὶν ἀγεννεστάτη, ἥτις βάλλεται ἔμπροσθεν τοῦ στρατοπέδου γυμνὴ οὖσα ὅπλων, καὶ ἐὰν ταραχὴ γένηται, ἐκείνη πρώτη ἀναιρεῖται. τοῦτο δὲ λέγει, ἐπεὶ ἐψίλωται ; Hsch.~π~3280 πριαμωθήσομαι· ξυρήσομαι. ἐπειδὴ τὸ τραγικὸν τοῦ Πριάμου πρόσωπον ξυρίας ἐστίν (com. ad. fr. 414 K.–A.) ; Hsch.~ψ~192 ψιλὰ τὰ γένεια· σπανοπώγων ; Hsch.~ψ~203 ψιλόν· γυμνόν. ἔρημον. εὐτελές. ἄτριχον. καὶ πᾶν ἀποδεδαρμένον (ν 437) ; Ar.~\textit{Thesm.}~583 ἕως ἂν οὕτως τὰς γνάθους ψιλὰς ἔχῃς ; Eur.~\textit{HF}~159–161 ὃς οὔποτ᾽ ἀσπίδ᾽ ἔσχε πρὸς λαιᾷ χερὶ / οὐδ᾽ ἦλθε λόγχης ἐγγύς, ἀλλὰ τόξ᾽ ἔχων, / κάκιστον ὅπλον, τῇ φυγῇ πρόχειρος ἦν ; Soph.~\textit{OC}~1029 οὐ ψιλὸν οὐδ᾽ ἄσκευον ἐς τοσήνδ᾽ ὕβριν ; Th.~4.32.4 οἱ πολέμιοι ἔσεσθαι ψιλοὶ καὶ οἱ ἀπορώτατοι, τοξεύμασι καὶ ἀκοντίοις καὶ λίθοις καὶ σφενδόναις ἐκ πολλοῦ ἔχοντες ἀλκήν ; Th.~8.97.1 εἶναι δὲ αὐτῶν ὁπόσοι καὶ ὅπλα παρέχονται ; D.~23.148 ὅσα μὲν δὴ στρατιώτης ὢν ἐν σφενδονήτου καὶ ψιλοῦ μέρει τὸ ἀπ᾽ ἀρχῆς ἐναντί᾽ ἐστράτευται τῇ πόλει ; Arist.~\textit{Pol.}~1297b1 δεῖ δὲ τὴν πολιτείαν εἶναι μὲν ἐκ τῶν τὰ ὅπλα ἐχόντων μόνον ; Arist.~\textit{Pol.}~1321a13 ἡ δὲ ψιλὴ δύναμις καὶ ναυτικὴ δημοκρατικὴ πάμπαν ; Lys.~14.14 καὶ ἡδέως ἂν οἱ μὲν ταῖς πόλεσι καταμείναντες ἐθεραπεύοντο, οἱ δὲ οἴκαδ᾽ ἀπελθόντες τῶν οἰκείων ἐπεμέλοντο, οἱ δὲ ψιλοὶ ἐστρατεύοντο ; Σ Soph.~\textit{Aj.}~1123 κἂν ψιλὸς ἀρκέσαιμι. τὰ τοιαῦτα σοφίσματα οὐκ οἰκεῖα τραγῳδίας ; Hsch.~π~1420 πέξις· ψίλωσις, κουρά, ξύρησις ; Hsch.~σ~15 σάβυττος· εἶδος ξυρήσεως εἰς καλλωπισμόν ; Hsch.~ψ~195 ψίλακα· ψιλόν, λεῖον. πτερὸν ἢ πτενόν (Ar. fr. 922 K.–A.) ; Hsch.~ψ~202 ψιλοκόρρης· φαλακρός (com. ad. fr. 441 K.–A.) ; Hsch.~ψ~206 ψιλούς· γυμνούς. σφενδονιστάς. τοξότας (Thuc. 1, 60, 1) ; Hsch.~ψ~207 ψίλωθρον· σποδή. ἔνιοι μάδον (Nic. ther. 902) ; Hsch.~ψ~209 ψιλῶς· κούφως. εὐτελῶς ; Hsch.~ψ~244 ψολῶσαι· ψιλῶσαι. ῥινῆσαι ; Suid.~ο~466 ψιλοὶ δὲ οἳ κουφοτάτῃ κέχρηνται, τόξῳ καὶ ἀκοντίῳ καὶ λίθοις ἐκ σφενδόνης ἢ ἐκ χειρός ; Suid.~ψ~106 Ψιλοί: οἱ μὴ καθωπλισμένοι μέν, ἐν καιρῷ δὲ τῆς μάχης τῷ παρατυχόντι ἢ λίθῳ ἢ ξύλῳ ἢ ἄλλῳ τινὶ ὀργάνῳ χρώμενοι· ἢ τοξόται ; Et. Gud.~576,~24–25 Sturz ψῶ τὸ λεπτύνω, ἐξ οὗ καὶ ψιλὸν στρατιώτην λέγομεν τὸν ἄοπλον καὶ γυμνόν ; Et. Parv. s.v.~Δαψιλῶς (Miller, Mélanges, p.~324) Δαψιλῶς· ψιλὸν λέγεται τὸ μικρὸν καὶ ταπεινὸν, ὡς καὶ ψιλὸς στρατιώτης ὁ μὴ ἔχων ὅπλα ; \textit{idem} Et. Gud.~134,~57–135,~2 Sturz (Δαψιλός ; μόνον pro μικρὸν ; ὁ μεμονωμένος καὶ ὅπλα μὴ ἔχων) ; Eup.~fr.~384,~7–8~K.-A. (=~117~Kock) νυνὶ δ᾽ ὅπηι τύχοιμεν / στρατευόμεσθ᾽ αἱρούμενοι καθάρματα στρατηγούς ; Eur.~\textit{Phoen.}~1147 γυμνῆτες ἱππῆς ἁρμάτων τ᾽ ἐπιστάται ; Soph.~\textit{Ant.}~426 οὕτω δὲ χαὔτη, ψιλὸν ὡς ὁρᾷ νέκυν ; Soph.~\textit{OC}~866 ὅς μ᾽, ὦ κάκιστε, ψιλὸν ὄμμ᾽ ἀποσπάσας ; Hdt.~9.29.1 ὁπλῖται μὲν οἱ πάντες συλλεγέντες ἐπὶ τὸν βάρβαρον ἦσαν τοσοῦτοι, ψιλῶν δὲ πλῆθος ἦν τόδε ; Th.~1.60.1 ὁπλίτας καὶ ψιλοὺς τετρακοσίους ; Th.~2.13.8 ἑξακοσίους δὲ καὶ χιλίους τοξότας ; Th.~2.79.3 οἱ δὲ ἱππῆς τῶν Χαλκιδέων καὶ ψιλοὶ νικῶσι τοὺς τῶν Ἀθηναίων ἱππέας καὶ ψιλούς ; Th.~3.22.3 ἔπειτα ψιλοὶ δώδεκα ξὺν ξιφιδίῳ καὶ θώρακι ἀνέβαινον ; Th.~4.28.4 Λημνίους δὲ καὶ Ἰμβρίους τοὺς παρόντας καὶ πελταστὰς ; Th.~4.129.2 καὶ Θρᾳξὶ μισθωτοῖς χιλίοις καὶ ἄλλοις τῶν αὐτόθεν ξυμμάχων πελτασταῖς ; Th.~6.43 καὶ σφενδονήταις Ῥοδίων ἑπτακοσίοις, καὶ Μεγαρεῦσι ψιλοῖς φυγάσιν εἴκοσι καὶ ἑκατόν ; Th.~7.60.4 τοξότας τε ἐπ᾽ αὐτὰς πολλοὺς καὶ ἀκοντιστὰς τῶν τε Ἀκαρνάνων καὶ τῶν ἄλλων ξένων ; Ar.~\textit{Byz.} Epit.~2.139 ἔχει δὲ ὁ μὲν ἄρρην χαίτας, ἡ δὲ θήλεια ψιλή ἐστι ; Hsch.~ε~6243 ἐσκνᾶσθαι· ξυρῆσαι ; Hsch.~κ~3471 κομμώτρια· ἡ κοσμοῦσα τὰς τρίχας καὶ ψιλοῦσα, κούρισσα ; Hsch.~μ~36 μαδίσας· τὰς τρίχας ἀποβαλών ; Hsch.~σ~1963 στρατεύομαι· διατρίβω (Ar. Nub. 131)}
\lp{232 ψιλός (rasé, épilé)]}{Hdt.~5.20.3 αὐτὸς δὲ ὁ Ἀλέξανδρος ἴσους τῇσι γυναιξὶ ἀριθμὸν ἄνδρας λειογενείους τῇ τῶν γυναικῶν ἐσθῆτι σκευάσας ; Ar.~\textit{Ach.}~158 τίς τῶν Ὀδομάντων τὸ πέος ἀποτεθρίακεν ; Suid.~α~3586 Ἀποτεθρίακεν. ἀντὶ τοῦ ἔτιλεν, ψιλὸν εἰργάσατο ; Hdt.~4.26.2 τὴν δὲ κεφαλὴν αὐτοῦ ψιλώσαντες ; Ael.~\textit{VH}~11.4 ψιλουμένης γὰρ αὐτῆς, κατὰ μικρὰ ὑπορρεουσῶν αὐτῷ τῶν τριχῶν}
\lp{233 μὴ φροντίσῃ]}{Ar.~\textit{Vesp.}~998 μὴ φροντίσῃς ὦ δαιμόνι᾽· ἀλλ᾽ ἀνίστασο ; Ar.~\textit{Eq.}~1356 ἀλλ᾽ οὐ σὺ τούτων αἴτιος, μὴ φροντίσῃς ; Ar.~\textit{Vesp.}~25 ἰδόντι τοιοῦτον ἐνύπνιον; / μὴ φροντίσῃς ; Ar.~\textit{Thesm.}~247 μὴ φροντίσῃς· ἕτερος γὰρ αὐτὰ σφογγιεῖ ; Ar.~\textit{Eccl.}~549 μὴ φροντίσῃς· ἄρρεν γὰρ ἔτεκε παιδίον ; Alex.~fr.~124.3~Kock (=~129~K.-A. ; Ath.~9.383c) ἔπειτα προσκέκαυκε. γλ. μηδὲν φροντίσῃς· ἰάσιμον γὰρ τὸ πάθος ἐστί ; Soph.~\textit{Phil.}~1404 μὴ φροντίσῃς ; Hdt.~6.129.4 ὁ δὲ Ἱπποκλείδης ὑπολαβὼν εἶπε οὐ φροντὶς Ἱπποκλείδῃ. ἀπὸ τούτου μὲν τοῦτο ὀνομάζεται ; Ar.~\textit{Vesp.}~228 μὴ φροντίσῃς· ἐὰν ἐγὼ λίθους ἔχω ; Ar.~\textit{Ran.}~493 οὐ μὰ Δί᾽ οὐδ᾽ ἐφρόντισα ; Ar.~\textit{Plut.}~215 μὴ φρόντιζε μηδὲν ὦγαθέ ; Eur.~\textit{Ion}~257 ἐγώ τε σιγῶ, καὶ σὺ μὴ φρόντιζ᾽ ἔτι ; X.~\textit{Cyr.}~2.2.4 ἀλλὰ μὴ φρόντιζε· αὐτίκα γὰρ ἀφ᾽ ἡμῶν ἄρξεται ; Plu.~\textit{Mor.}~867b (De Hdt. malign.~33) εἰπεῖν ἂν ἐξορχούμενος τὴν ἀλήθειαν οὐ φροντὶς Ἡροδότῳ ; Luc.~\textit{Apol.}~15 ἱκανὸν ἂν εἴη μοι τό οὐ φροντὶς Ἱπποκλείδῃ ; Hsch.~ε~7318 εὐωρεῖν· τὸ μηδενὸς ἔχειν λόγον, ἢ μηδενὸς φροντίζειν ; Hsch.~ε~7320 εὐωριάζειν· ὀλιγωρεῖν, μὴ ἔχειν φροντίδα. παρακούειν (Soph. fr. 561 R.) ; Hsch.~φ~909 φροντίζει· μεριμνᾷ ; Hsch.~ο~1920 οὐ φροντὶς Ἱπποκλείδῃ· παροιμία, ἧς μέμνηται καὶ Ἕρμιππος Δημόταις (fr. 16 K.–A.) ; \textit{idem} Suid.~ο~978 (om. καὶ ; ἐν Δημόταις), Zen.~5.31 (lemma tantum), Diogenian.~7.21 (lemma tantum)}
\lp{233 εὐπρεπὴς φανεῖ]}{Ar.~\textit{Thesm.}~192 γυναικόφωνος ἁπαλὸς εὐπρεπὴς ἰδεῖν ; Eur.~\textit{El.}~948–951 ἀλλ᾽ ἔμοιγ᾽ εἴη πόσις / μὴ παρθενωπός, ἀλλὰ τἀνδρείου τρόπου. / τὰ γὰρ τέκν᾽ αὐτῶν Ἄρεος ἐκκρεμάννυται, / τὰ δ᾽ εὐπρεπῆ δὴ κόσμος ἐν χοροῖς μόνον ; Ar.~\textit{Eccl.}~427 μετὰ τοῦτο τοίνυν εὐπρεπὴς νεανίας ; Eur.~\textit{IA}~822 γυναῖκα, μορφὴν εὐπρεπῆ κεκτημένην ; Hsch.~ε~7085 εὐπρέπεια· εὐμορφία (2. Regn. 15, 25) ; Hsch.~κ~483 καλλιπάρηος· καλὰς παρειὰς ἔχουσα, εὐπρεπής (Α 143) ; Moer.~ε~39 Hansen εὔμορφον παρ᾽ οὐδενὶ εὗρον, ἀλλὰ εὐπρεπῆ. εὐμορφίαν παρὰ Πλάτωνι ἐν Νόμοις (716a; 744c) καὶ παρὰ Ξενοφῶντι (Symp. 8, 6), σπάνιον δέ ; Ar.~\textit{Lys.}~1315 ἁγνὰ χοραγὸς εὐπρεπής ; Ar.~\textit{Eccl.}~701 τοῖς εὐπρεπέσιν δ᾽ ἀκολουθοῦντες ; Eur.~\textit{Alc.}~161 ἐσθῆτα κόσμον τ᾽ εὐπρεπῶς ἠσκήσατο ; Eur.~\textit{IA}~386 εὐπρεπῆ γυναῖκα χρῄζεις ; Hsch.~ε~737 εἰδαλίμας· καλάς. ἐπὶ τοῦ εἴδους εὐμόρφους, εὐπρεπεῖς (ω 279) ; Hsch.~ε~7086 εὐπρεπῶς· εὐκόσμως (1. Esdr. 1, 10) ; Hsch.~ι~165 ἰδανή· τρυφερή, εὐπρεπή⟨ς⟩ (Callim. fr. 114, 9 Pf.)}
\lp{233 πάνυ]}{Ar.~\textit{Thesm.}~749 πάνυ γ᾽ ἐμπίμπρατε ; Ar.~\textit{Ran.}~137 εὐθὺς γὰρ ἐπὶ λίμνην μεγάλην ἥξεις πάνυ ; Eur.~\textit{Cyc.}~646 ἀλλ᾽ οἶδ᾽ ἐπῳδὴν Ὀρφέως ἀγαθὴν πάνυ ; Aesch.~\textit{Ag.}~1456 μία τὰς πολλάς, τάς πάνυ πολλὰς ; Aesch.~\textit{Pers.}~926 τοξοδάμαντες, πάνυ ταρφύς τις ; Theodect. ap. Ath.~10.454 ἔπειτα δισσοὶ κανόνες ἰσόμετροι πάνυ ; Xenoph. B~1.18~D.-K. (Ath.~11.462c–463a) οἴκαδ᾽ ἄνευ προπόλου, μὴ πάνυ γηραλέος ; Th.~6.18.6 καὶ τὸ πάνυ ἀκριβὲς ; Pl.~\textit{Smp.}~218c πάνυ ἀνόητον ἡγοῦμαι εἶναι ; D.~56.37 πάνυ ἁπλοῦν ἐστιν διαγνῶναι ὑμῖν}
\lp{233 μὴ φροντίσῃ· ὡς εὐπρεπὴς φανεῖ πάνυ (propos de barbier)]}{Ar.~\textit{Plut.}~337–339 καίτοι λόγος γ᾽ ἦν νὴ τὸν Ἡρακλέα πολὺς / ἐπὶ τοῖσι κουρείοισι τῶν καθημένων, / ὡς ἐξαπίνης ἀνὴρ γεγένηται πλούσιος ; Plu.~\textit{De\,garr.}~13 (Mor.~508f–509c) χαριέντως γοῦν ὁ βασιλεὺς Ἀρχέλαος, ἀδολέσχου κουρέως περιβαλόντος αὐτῷ τὸ ὠμόλινον καὶ πυθομένου πῶς σε κείρω, βασιλεῦ σιωπῶν ἔφη ; \textit{idem} Plu.~\textit{Mor.}~177a (Reg. et imp. apophth., Archel.~2) (ἀδολέσχου δὲ κουρέως ἐρωτήσαντος αὐτὸν πῶς σε κείρω; σιωπῶν ἔφη) ; Hor.~\textit{Sat.}~1.7.3 omnibus et lippis notum et tonsoribus esse ; \textit{idem} Tosi n°~1967 (« Omnibus et lippis notum et tonsoribus ») ; Luc.~\textit{Hist.Conscr.}~24 οὔτε τὸ λεγόμενον δὴ τοῦτο τῶν ἐπὶ κουρείῳ τὰ τοιαῦτα μυθολογούντων ἀκούσας ; Thphr.~fr.~76 Wimmer (ap. Plu.~\textit{Quaest.\,conv.}~5.5.2, Mor.~678e–679e) Διὸ καὶ Θ. ἄοινα συμπόσια παίζων ἐκάλει τὰ κουρεῖα διὰ τὴν λαλιὰν τῶν προσκαθιζόντων ; Plb.~3.20.5 οὐ γὰρ ἱστορίας, ἀλλὰ κουρεακῆς καὶ πανδήμου λαλιᾶς ἔμοιγε δοκοῦσι τάξιν ἔχειν καὶ δύναμιν}
\lp{234 βούλει θεάσασθ᾽ αὑτόν;]}{Ar.~\textit{Thesm.}~140 τίς δαὶ κατόπτρου καὶ ξίφους κοινωνία ; Ar.~fr.~332.1~K.-A. (=~320.1~Kock ; Θεσμοφοριάζουσαι β´ ; Poll.~7.95) ξυρόν, κάτοπτρον, ψαλίδα, κηρωτήν, λίτρον ; Eur.~\textit{El.}~1071 ξανθὸν κατόπτρῳ πλόκαμον ἐξήσκεις κόμης ; Eur.~\textit{Hipp.}~428–430 προθεὶς κάτοπτρον ὥστε παρθένῳ νέᾳ ; Eur.~\textit{Med.}~1161 λαμπρῷ κατόπτρῳ σχηματίζεται κόμην ; Eur.~\textit{Or.}~1112 οἵους ἐνόπτρων καὶ μύρων ἐπιστάτας ; Anacreont.~6 [11]~Bergk,~1–5 Λέγουσιν αἱ γυναῖκες· / Ἀνακρέων, γέρων εἶ, / λαβὼν ἔσοπτρον ἄθρει / κόμας μὲν οὐκέτ᾽ οὔσας, / ψιλὸν δέ σευ μέτωπον ; Plu.~\textit{Mor.}~42b (De rect. aud.~8) οὐ γὰρ ἐκ κουρείου μὲν ἀναστάντα δεῖ τῷ κατόπτρῳ παραστῆναι καὶ τῆς κεφαλῆς ἅψασθαι, τὴν περικοπὴν τῶν τριχῶν ἐπισκοποῦντα καὶ τῆς κουρᾶς τὴν διαφοράν ; Luc.~\textit{Ind.}~29 τοὺς κουρέας τούτους ἐπίσκεψαι, καὶ ὄψει τοὺς μὲν τεχνίτας αὐτῶν ξυρὸν καὶ μαχαιρίδας καὶ κάτοπτρον σύμμετρον ἔχοντας ; Poll.~1.217 αἱ δὲ ὥσπερ ἐν κατόπτρῳ θεασάμεναι τὴν αἰσχύνην τοῦ σώματος ἀπηγλαϊσμένην τῆς κόμης, ἀνέχονται τότε τὴν πρὸς τὰ χείρω ὁμιλίαν ; Poll.~10.126–127 (Bethe II, p.~228) τῶν δὲ γυναικείων σκευῶν κτένιον, ξάνιον, κομμώτριον, ξυρόν, κάτοπτρον, οὗ τὴν θήκην λοφεῖον καλοῦσι, ψαλίς ; Tosi n°~1900 (com. adesp. fr.~903~K.-A. = Stob.~4.30.6a) τίς γὰρ κατόπτρῳ καὶ τυφλῷ κοινωνία ; Ar.~\textit{Nub.}~1454–1455 αὐτὸς μὲν οὖν σαυτῷ σὺ τούτων αἴτιος, / στρέψας σεαυτὸν ἐς πονηρὰ πράγματα ; Ar.~\textit{Pax}~907 θέασ᾽ ὡς προθύμως ὁ πρύτανις παρεδέξατο ; Ar.~\textit{Thesm.}~1178 ὀρχησομένη γὰρ ἔρχεθ᾽ ὡς ἄνδρας τινάς ; Alexand.~fr.~1~K.-A. (Διόνυσος ; Σ A Il.~1.216) ἰδοὺ κάτοπτρον· εἰπέ μοι, τούτῳ τί χρῇ ; Eur.~\textit{Hec.}~925–926 χρυσέων ἐνόπτρων λεύσ- ; Eur.~fr.~322,~3 Kn. φιλεῖ κάτοπτρα καὶ κόμης ξανθίσματα ; Soph.~\textit{El.}~344 κείνης διδακτά, κοὐδὲν ἐκ σαυτῆς λέγεις ; Soph.~\textit{OT}~707 σύ νυν ἀφεὶς σεαυτὸν ὧν λέγεις πέρι ; Pl.~\textit{Alc.} I~133a τὸ πρόσωπον ἐμφαίνεται ἐν τῇ τοῦ καταντικρὺ ὄψει ὥσπερ ἐν κατόπτρῳ ; Plu.~\textit{Mor.}~456a (De cohib. ira~6) ἔσοπτρον, ὥσπερ ἐνίοις προσφέρουσι λουσαμένοις ἐπ᾽ οὐδενὶ χρησίμῳ. τὸ δ᾽ αὑτὸν ἰδεῖν παρὰ φύσιν ἔχοντα καὶ συντεταραγμένον οὐ μικρόν ἐστιν εἰς διαβολὴν τοῦ πάθους ; Luc.~\textit{Pisc.}~45 χρυσίον τουτὶ καὶ μύρον καὶ μαχαίριον κουρευτικὸν καὶ κάτοπτρον καὶ κύβους ; D.L.~2.33 ἠξίου δὲ καὶ τοὺς νέους συνεχὲς κατοπτρίζεσθαι ; Hsch.~ε~3248 ἔνοπτρα· κάτοπτρα, ἔσοπτρον (Eur. Tr. 1107) ; Hsch.~κ~1843 κάτοπτρα· δίοπτρον ; Plu.~\textit{De\,imposs.}~26–27 Φαλακρῷ κτένας δανείζεις. Τυφλῷ κάτοπτρον χαρίζῃ ; Tosi n°~1425 Κάτοπτρον εἴδους χαλκός ἐστ᾽, οἶνος δὲ νοῦ ; Luc.~\textit{Ind.}~19 οὐ μᾶλλον ἢ φαλακρὸς ἄν τις πρίαιτο κτένας ἢ κάτοπτρον ὁ τυφλὸς ἢ ὁ κωφὸς αὐλητὴν ; Ar.~\textit{Ach.}~275 καταγιγαρτίσ᾽ ὦ ; Ar.~\textit{Nub.}~752 ὥσπερ κάτοπτρον, κᾆτα τηροίην ἔχων — ; Ar.~\textit{Pax}~392 ἀλλὰ χάρισ᾽ ὦ φιλανθρωπότατε ; Ar.~\textit{Thesm.}~230 ἔχ᾽ ἀτρέμα σαυτὸν κἀνάκυπτε· ποῖ στρέφει ; Eur.~\textit{Tro.}~1107 χρύσεα δ᾽ ἔνοπτρα, παρθένων ; Aesch.~\textit{Ag.}~839 ὁμιλίας κάτοπτρον, εἴδωλον σκιᾶς ; Aesch.~fr.~61~R. (Ἠδωνοί) ποδαπὸς ὁ γύννις; τίς πάτρα; τίς ἡ στολή ; [Luc.]~\textit{Am.}~44 οὐδὲ ἔσοπτρα τῶν ἀντιμόρφων χαρακτήρων ἀγράφους εἰκόνας ; Artem.~5.67 Ἔδοξέ τις κάτοπτρον ἔχειν κουρέως ἑστὸς ἐν τῇ πλατείᾳ πλησίον ἀγορᾶς ; Hsch.~ε~6284 ἔσοπτρον· κάτοπτρον (Sir. 12, 11) ; Hsch.~ε~6473 ἔσοπτρον· κάτοπτρον. καθρέπτην ; Hsch.~κ~1844 κατοπτρίζεται· ἐν αἰσθήσει φαντάζεται, ἢ ὁρᾷ (Cyr. in Ep. 2. Cor. 340, 13) ; Antiatt. p.~102,~10~Bekker Κάτοπτρον οὔ φασι δεῖν λέγειν, ἀλλ᾽ ἔσοπτρον ; Suid.~ε~1391 Ἔνοπτρον: τὸ κάτοπτρον ; Suid.~λ~697 Λοφεῖον: ἡ τοῦ κατόπτρου θήκη, ἢ ἡ θήκη τοῦ λόφου τῆς περικεφαλαίας. περιφερὴς δὲ αὕτη. Ἀριστοφάνης Νεφέλαις· καθέλοιμι νύκτωρ τὴν σελήνην ; Apostol.~7.16i Ἐν μὲν τοῖς ἐσόπτροις ὁ τῆς ὄψεως, ἐν δὲ ταῖς ὁμιλίαις ὁ τῆς ψυχῆς χαρακτὴρ βλέπεται ; Eust.~1587.47 Ὅτι οὐδὲν μὲν θεάτρου μετὸν τυφλῷ ἀνδρὶ ὥσπερ οὐδὲ κατόπτρου κατὰ τὴν παροιμίαν}
\lp{234 εἰ δοκεῖ]}{Ar.~\textit{Nub.}~11 ἀλλ᾽ εἰ δοκεῖ ῥέγκωμεν ἐγκεκαλυμμένοι ; Ar.~\textit{Av.}~56 πάνυ γ᾽, εἰ δοκεῖ ; Ar.~\textit{Thesm.}~216 εἴ σοι δοκεῖ ; Pl.~\textit{Cra.}~383a βούλει οὖν καὶ Σωκράτει τῷδε ἀνακοινωσώμεθα τὸν λόγον; εἴ σοι δοκεῖ ; Pl.~\textit{Thg.}~121a βούλει οὖν δεῦρο εἰς τὴν τοῦ Διὸς τοῦ ἐλευθερίου στοὰν ἐκποδὼν ἀποχωρήσωμεν; εἰ σοὶ δοκεῖ ; Ar.~\textit{Lys.}~1176 ἀλλ᾽ εἰ δοκεῖ δρᾶν ταῦτα}
\lp{234 φέρε]}{Ar.~\textit{Thesm.}~261 φέρ᾽ ἔγκυκλον ; Ar.~\textit{Ran.}~498 φέρε δὴ ταχέως αὔτ᾽· οὐ γὰρ ἀλλὰ πειστέον ; Ar.~\textit{Ach.}~584 φέρε νυν ἀπὸ τοῦ κράνους μοι τὸ πτερόν ; Ar.~\textit{Eq.}~118 φέρ᾽ αὐτὸν ἵν᾽ ἀναγνῶ ; Ar.~\textit{Nub.}~1297 φέρε μοι τὸ κέντρον ; Ar.~\textit{Thesm.}~1115 φέρε δεῦρό μοι τὴν χεῖρ᾽, ἵν᾽ ἅψωμαι κόρης ; Ar.~\textit{Thesm.}~1196 τἀργύριον τοίνυν φέρε}
\lp{235 ὁρᾷς σεαυτόν;]}{Hermipp.~fr.~7~Kock (Ἀθηνᾶς γοναί ; Σ Pl. p.~353~Bekker) τὸ ἀκκίζεσθαι ἐκ γυναικὸς εἰρῆσθαί φασιν Ἀκκοῦς … εἴς τε τὸ κάτοπτρον βλέπουσαν πρὸς τὴν παρ᾽ αὑτῆς ἔμφασιν … ὡς ἑτέρᾳ προσλαλεῖν γυναικί. μέμνηται ταύτης Ἕρμιππος ἐν Ἀθηνᾶς γοναῖς ; Zen.~1.53 Ἀκκώ: ἐπὶ τῶν μωραινόντων. Ἡ γὰρ Ἀκκὼ γυνὴ γέγονεν ἐπὶ μωρίᾳ διαβεβοημένη, ἥν φασιν ἐσοπτριζομένην τῇ εἰκόνι ὡς ἑτέρᾳ διαλέγεσθαι ; \textit{idem} Diogenian.~2.4 (Ἀκκίζεσθαι ; Ἀκκὼ γυνή τις ἐπὶ μωρίᾳ διαβαλλομένη, ἥτις ἐσοπτριζομένη τῇ οἰκείᾳ σκιᾷ ὡς ἑτέρᾳ διελέγετο), Suid.~α~946 (γυνὴ ἐπὶ μωρίᾳ διαβαλλομένη ; ἐνοπτριζομένην τῇ ἰδίᾳ εἰκόνι ; add. ὅθεν καὶ τὸ ἀκκίζεσθαι. ταύτην φασὶ καὶ σπόγγῳ πάτταλον κρούειν), Apostol.~1.51 (Ἀηκίζεσθαι μοι δοκεῖς: Ἀκὼ γυνή τις ἦν ἐπὶ μωρίᾳ διαβαλλομένη, ἥτις ἐσοπτριζομένη τῇ οἰκείᾳ σκιᾷ ὡς ἑτέρᾳ διελέγετο) ; Eur.~\textit{Ba.}~925–927 τί φαίνομαι δῆτ᾽; οὐχὶ τὴν Ἰνοῦς στάσιν / ἢ τὴν Ἀγαύης ἑστάναι, μητρός γ᾽ ἐμῆς; / αὐτὰς ἐκείνας εἰσορᾶν δοκῶ σ᾽ ὁρῶν ; Hsch.~α~2422 ἀκκίζειν· γυναικίζεσθαι. προσποιεῖσθαι μωραίνειν ; Hsch.~α~2423 ἀκκίζεται· θηλύνεται ; Suid.~α~878 Ἀκκιζόμενος: προσποιούμενος, γυναικιζόμενος, ἢ μωραίνων, ἀπό τινος γυναικὸς μωρᾶς, ἥτις ἐκαλεῖτο Ἀκκώ}
\lp{235 οὐ μὰ Δί᾽ ἀλλὰ]}{Ar.~\textit{Pax}~1045–1046 τίς ἄρα ποτ᾽ ἐστίν; … μάντις τίς ἐστιν. / οὐ μὰ Δί᾽ ἀλλ᾽ Ἱεροκλέης ; Ar.~\textit{Vesp.}~193 ἐγὼ πονηρός; οὐ μὰ Δί᾽ ἀλλ᾽ οὐκ οἶσθα σὺ ; Ar.~\textit{Eccl.}~556 τί δρᾶν; ὑφαίνειν; — οὐ μὰ Δί᾽ ἀλλ᾽ ἄρχειν ; Ar.~\textit{Plut.}~704 αὐτὸς δ᾽ ἐκεῖνος; — οὐ μὰ Δί᾽ οὐδ᾽ ἐφρόντισεν}
\lp{235 Κλεισθένη]}{\textit{Test.} Suid.~κ~1756 Κλεισθένην ὁρῶ· Ἀριστοφάνης. οὗτος ἀεὶ ἐξυρᾶτο. λέγεται οὖν ἐπὶ τῶν ἀεὶ ξυρωμένων. ἦν δὲ Σιβυρτίου· καὶ τοῦτο ἔπραττε διὰ τὸ νέος φαίνεσθαι ἀεί ; Ar.~\textit{Ach.}~117–121 καὶ τοῖν μὲν εὐνούχοιν τὸν ἕτερον τουτονὶ / ἐγᾦδ᾽ ὅς ἐστι, Κλεισθένης ὁ Σιβυρτίου. / ὦ θερμόβουλον πρωκτὸν ἐξυρημένε ; Ar.~\textit{Nub.}~355 καὶ νῦν γ᾽ ὅτι Κλεισθένη εἶδον, ὁρᾷς, διὰ τοῦτ᾽ ἐγένοντο γυναῖκες ; Ar.~\textit{Av.}~829–831 ὅπου θεὸς γυνὴ γεγονυῖα πανοπλίαν / ἕστηκ᾽ ἔχουσα, Κλεισθένης δὲ κερκίδα ; Ar.~\textit{Thesm.}~97–98 ἀλλ᾽ ἢ τυφλὸς μέν εἰμ᾽· ἐγὼ γὰρ οὐχ ὁρῶ / ἄνδρ᾽ οὐδέν᾽ ἐνθάδ᾽ ὄντα, Κυρήνην δ᾽ ὁρῶ ; Ar.~\textit{Thesm.}~574–575 φίλαι γυναῖκες ξυγγενεῖς τοὐμοῦ τρόπου· / ὅτι μὲν φίλος εἴμ᾽ ὑμῖν, ἐπίδηλος ταῖς γνάθοις ; Ar.~\textit{Ran.}~422–424 τὸν Κλεισθένους δ᾽ ἀκούω / ἐν ταῖς ταφαῖσι πρωκτὸν / τίλλειν ἑαυτοῦ καὶ σπαράττειν τὰς γνάθους ; Ar.~fr.~422~K.-A. (=~407~Kock ; Ὁλκάδες ; Σ Ar.~\textit{Ach.}~122) παῖδες ἀγένειοι, Στράτων. Schol. Στράτων· καὶ οὗτος κωμῳδεῖται ὡς λωβώμενος τὸ γένειον καὶ λειαίνων τὸ σῶμα, ὡς Κλεισθένης, ὥς φησιν αὐτὸς Ἀριστοφάνης ἐν Ὁλκάσι ; Σ Ar.~\textit{Nub.}~355 τοῦτον ἰδοῦσαι, φησὶν, αἱ Νεφέλαι, εἰς γυναῖκας μετεμόρφωσαν ἑαυτάς ; Σ Ar.~\textit{Thesm.}~235 διαβάλλει τοῦτον ὡς ξυρούμενον ; Suid.~σ~1186 Στράτων ἀεὶ ἐξυρᾶτο καὶ Κλεισθένης ἐπὶ τὸ φαίνεσθαι ἀεὶ νέοι. Ἀριστοφάνης· ὦ θερμόβουλον πρωκτὸν ἐξυρημένε ; Suid.~κ~1758 Κλεισθένους ἀκρατέστερος: οὗτος ἐπὶ κιναιδίᾳ διεβέβλητο καὶ ἐγυναικίζετο ; Ar.~\textit{Eq.}~1373–1374 οὐδ᾽ ἀγοράσει γ᾽ ἀγένειος οὐδεὶς ἐν ἀγορᾷ. / ποῦ δῆτα Κλεισθένης ἀγοράσει καὶ Στράτων ; Ar.~\textit{Lys.}~621 δεῦρο συνεληλυθότες ἄνδρες ἐς Κλεισθένους ; Ar.~\textit{Lys.}~1092 οὐκ ἔσθ᾽ ὅπως οὐ Κλεισθένη βινήσομεν ; Ar.~\textit{Ran.}~48 ἐπεβάτευον Κλεισθένει— ; Ar.~\textit{Ran.}~57 ξυνεγένου τῷ Κλεισθένει ; Cratin.~fr.~208~K.-A. (=~195~Kock ; Πυτίνη) ληρεῖς ἔχων· γράφ᾽ αὐτὸν / ἐν ἐπεισοδίωι. γελοῖος ἔσται Κλεισθένης κυβεύων / † ἐν τῆι τοῦ κάλλους ἀκμῆι ; \textit{idem} Suid.~κ~1758 (γράφων αὐτὸν ἐν ἐπεισοδίῳ ; om. † ἐν τῆι τοῦ κάλλους ἀκμῆι) ; Pherecr.~fr.~143~K.-A. (=~135~Kock ; Πετάλη ; Ath.~9.395b–c) ἀλλ᾽ ὦ περιστέριον ὅμοιον Κλεισθένει, / πέτου, κόμισον δέ μ᾽ ἐς Κύθηρα καὶ Κύπρον ; Lys.~25.25 ἴστε γὰρ Ἐπιγένη καὶ Δημοφάνη καὶ Κλεισθένη ἰδίᾳ μὲν καρπωσαμένους τὰς τῆς πόλεως συμφοράς, δημοσίᾳ δὲ ὄντας μεγίστων κακῶν αἰτίους ; Σ Ar.~\textit{Ach.}~118 οὗτος δὲ ὁ Κλεισθένης ἀεὶ τὸ γένειον ἐξυρᾶτο πρὸς τὸ ἀεὶ φαίνεσθαι νέος. διὸ εὐνούχῳ αὐτὸν εἰκάζει ; Σ Ar.~\textit{Eq.}~1374 Κλεισθένης καὶ Στράτων: Γύννιδες οὗτοι καὶ πιττούμενοι τὰ γένεια, λεῖοι δὲ καὶ μαλακοί ; Σ Ar.~\textit{Ran.}~48 Κλεισθένης ὁ Σιβυρτίου ἐπ᾽ αἰσχρότητι κωμῳδεῖται. ἦν δὲ λεῖος τὸ γένειον· διὸ καὶ εὐνούχῳ εἰκάζει αὐτόν ; \textit{idem} Suid.~ε~2021 (ὁ δὲ Κλεισθένης ; om. διὸ ; εἰκάζειν) ; Ar.~\textit{Vesp.}~1187 ὡς ξυνεθεώρεις Ἀνδροκλεῖ καὶ Κλεισθένει ; Ar.~\textit{Thesm.}~763 φύλαξον αὐτόν, ἵνα λαβοῦσα Κλεισθένη ; Timocl.~fr.~5~K.-A. (Δημοσάτυροι ; Ath.~4.165f) οὐδ᾽ ὁ Χαβρίου Κτήσιππος ἔτι τρὶς κείρεται, / ἐν ταῖς γυναιξὶ λαμπρός, οὐκ ἐν ἀνδράσιν ; Σ Ar.~\textit{Lys.}~1092 Οὗτος Σιβυρτίου ὁ παῖς ἐπὶ θηλύτητι κωμῳδούμενος}
\lp{236 ἀνίστασ᾽]}{Ar.~\textit{Vesp.}~998 μὴ φροντίσῃς ὦ δαιμόνι᾽· ἀλλ᾽ ἀνίστασο ; Ar.~\textit{Thesm.}~643 ἀνίστασ᾽ ὀρθός. ποῖ τὸ πέος ὠθεῖς κάτω ; Ar.~\textit{Vesp.}~286 ἀλλ᾽ ὦγάθ᾽ ἀνίστασο μηδ᾽ οὕτω σεαυτὸν ; Ar.~\textit{Lys.}~929 ἀνίστασ᾽, ἀναπήδησον ; Eur.~\textit{Hec.}~499 ἀνίστασ᾽, ὦ δύστηνε, καὶ μετάρσιον}
\lp{236 ἵν᾽ ἀφεύσω σε]}{\textit{Test.} Suid.~α~4594 Ἀφευθεὶς τὸν πρωκτόν: Ἀριστοφάνης ἀντὶ τοῦ φλογισθείς ; Ar.~\textit{Eq.}~1236 ἐν ταῖσιν εὕστραις κονδύλοις ἡρμοττόμην ; Ar.~\textit{Lys.}~825–828 ἀλλ᾽ ὅμως ἂν οὐκ ἴδοις / καίπερ οὔσης γραὸς ὄντ᾽ αὐτὸν / κομήτην, ἀλλ᾽ ἀπεψιλωμένον / τῷ λύχνῳ ; Ar.~\textit{Thesm.}~216 τὰ κάτω δ᾽ ἀφεύειν ; Ar.~\textit{Thesm.}~590–591 ἀφηῦσεν αὐτὸν κἀπέτιλ᾽ Εὐριπίδης / καὶ τἄλλ᾽ ἅπανθ᾽ ὥσπερ γυναῖκ᾽ ἐσκεύασεν ; Ar.~\textit{Eccl.}~12–13 μόνος δὲ μηρῶν εἰς ἀπορρήτους μυχοὺς / λάμπεις ἀφεύων τὴν ἐπανθοῦσαν τρίχα ; Aesch.~fr.~310~R. (Ath.~9.375e) λευκός· τί δ᾽ οὐχί; καὶ καλῶς ἠφευμένος / ὁ χοῖρος. ἕψου, μηδὲ λυπηθῇς πυρί ; \textit{Il.}~9.467–468 πολλοὶ δὲ σύες θαλέθοντες ἀλοιφῇ / εὑόμενοι τανύοντο διὰ φλογὸς Ἡφαίστοιο ; \textit{idem} \textit{Il.}~23.32–33 (δ᾽ ἀργιόδοντες ὕες) ; \textit{Od.}~14.74–75 ἔνθεν ἑλὼν δύ᾽ ἔνεικε καὶ ἀμφοτέρους ἱέρευσεν, / εὗσέ τε μίστυλλέν τε καὶ ἀμφ᾽ ὀβελοῖσιν ἔπειρεν ; Semon.~fr.~24,~1~W. (Ath.~14.659e) χὠς ὗν ἄφευσα χὠς ἐμίστυλα κρέα ; Σ Ar.~\textit{Eq.}~1236 ἐν τοῖς μαγειρείοις. εὖστρα γὰρ τὸ μαδιστήριον, ἀπὸ τοῦ εὕειν καὶ φλογίζειν τοὺς χοίρους ; Hsch.~ε~7183 εὗσεν· ἐμάδισεν (ι 389) ; Hsch.~ε~7216 εὕστρα· βόθρος, ἐν ᾧ τοὺς ὗς εὕουσι. καὶ ὁ στάχυς, ὅταν μήπω πέπειρος ὢν ἀποφλογισθῇ ; \textit{idem} EM~398.31 (s.v.~Εὔστρα) (σῦς ; om. ὢν ; add. Εὖστραι, βόθυνοι ἐν οἷς οἱ χοῖροι ἤγουν οἱ σῦς εὕονται, καὶ καίεται ἱερεῖα) ; Hsch.~ε~7308 εὕω· καίω, φλογίζω ; Hsch.~σ~1056 σκόλεφραι· κατακεκαυμέναι τὰς τρίχας ; Poll.~6.91 τὰ δὲ τῶν μαγείρων ἔργα ἀφεῦσαι, εὗσαι, καθῆραι, κόψαι, τεμεῖν … εὗστραι δ᾽ οἱ βόθροι ἐκαλοῦντο, ἐν οἷς εὕεται τὰ χοιρίδια ; \textit{Od.}~2.300 αἶγας ἀνιεμένους σιάλους θ᾽ εὕοντας ἐν αὐλῇ ; \textit{Od.}~14.426 τὸν δ᾽ ἔλιπε ψυχή. τοὶ δ᾽ ἔσφαξάν τε καὶ εὗσαν ; Hsch.~ε~7174 εὖσαι· καῦσαι ; Hsch.~ε~7176 εὔσανα· χύτρα. [ὄρυγμα, ἐν ᾧ τοὺς ὗς βυθίζουσι.] καὶ τὰ ἐγκαύματα ; Hsch.~π~4471 †πυρσολεῖφοι ἢ πυρσόλειφθοι†· ὑπὸ πυρὸς μεμαδισμένοι καὶ †πεφλασμένοι† ; Phot.~ε~2342–2343,~2347 εὗστρα, οὐχὶ φλογίστρα. … — εὗστραι· βόθυνοι, ἐν οἷς χοῖροι ἀφεύονται καὶ ἕτερα ἱερεῖα. … — εὗστραι· βόθροι, ἐν οἷς οἱ σῦς εὕονται ; Suid.~ε~3756 Εὔστρα: τὸ μαδιστήριον, ἢ ἡ φλογίστρα. Ἀριστοφάνης· ἐν ταῖσιν εὔστραις κονδύλοις ἡρμοττόμην ; Ar.~\textit{Pax}~1144 ἀλλ᾽ ἄφευε τῶν φασήλων ὦ γύναι τρεῖς χοίνικας ; Hsch.~ε~7194 εὕσοντας· ἐκδέροντας ; Hsch.~σ~1068 σκολοφρή· κατακεκαυμένη ; EM~398 (s.v.~Εὐσῶ) Εὐσῶ τὸ χοιρίδιον. οἱ οὕτω λέγοντες ἁμαρτάνουσι. δεῖ γὰρ θῦσαι λέγειν· καὶ παρ᾽ Ὁμήρῳ δὲ κακῶς γράφεται, σιάλους τ᾽ εὕοντας ἐν αὐλῇ}
\lp{236 κἀγκύψας ἔχε]}{Ar.~\textit{Nub.}~191–193 τί γὰρ οἵδε δρῶσιν οἱ σφόδρ᾽ ἐγκεκυφότες; / οὗτοι δ᾽ ἐρεβοδιφῶσιν ὑπὸ τὸν Τάρταρον. / τί δῆθ᾽ ὁ πρωκτὸς ἐς τὸν οὐρανὸν βλέπει ; Moer.~ε~73 Hansen εἴσκυψον, οὐκ ἔγκυψον· Ἀριστοφάνης ; Ar.~\textit{Av.}~1082 τὰς περιστεράς θ᾽ ὁμοίως ξυλλαβὼν εἵρξας ἔχει ; Ar.~\textit{Thesm.}~706 δεινὰ δῆθ᾽, ὅστις γ᾽ ἔχει μου ᾽ξαρπάσας τὸ παιδίον ; Ar.~\textit{Ran.}~425 κἀκόπτετ᾽ ἐγκεκυφώς ; Ar.~\textit{Ran.}~804 ἔβλεψε γοῦν ταυρηδὸν ἐγκύψας κάτω ; Ar.~\textit{Eccl.}~355 ἀχράς τις ἐγκλῄσασ᾽ ἔχει τὰ σιτία ; Pl.~\textit{Phdr.}~254d ἐγκύψας καὶ ἐκτείνας τὴν κέρκον, ἐνδακὼν τὸν χαλινόν, μετ᾽ ἀναιδείας ἕλκει ; Ar.~\textit{Thesm.}~489 παρὰ τὸν Ἀγυιᾶ κύβδ᾽ ἐχομένη τῆς δάφνης ; Ar.~\textit{Eccl.}~957 πόθος, ὅς με διακναίσας ἔχει ; Hes.~\textit{Op.}~42 κρύψαντες γὰρ ἔχουσι θεοὶ βίον ἀνθρώποισιν ; Pl.~\textit{R.}~359d ἵππον χαλκοῦν, κοῖλον, θυρίδας ἔχοντα, καθ᾽ ἃς ἐγκύψαντα ἰδεῖν ἐνόντα νεκρόν ; Pl.~\textit{R.}~555e οἱ δὲ δὴ χρηματισταὶ ἐγκύψαντες, οὐδὲ δοκοῦντες τούτους ὁρᾶν ; Teucer ap. Ath.~10.455e ὄμματά τ᾽ ἐκκύπτοντα προμήκεα κεἰσκύπτοντα ; Hsch.~κ~4657 κύπτει· κάμπτει, ἑαυτὸν κλίνει}
\lp{237 οἴμοι κακοδαίμων]}{Ar.~\textit{Vesp.}~207 οἴμοι κακοδαίμων, στροῦθος ἁνὴρ γίγνεται ; Ar.~\textit{Thesm.}~232 οἴμοι κακοδαίμων, ψιλὸς αὖ στρατεύσομαι}
\lp{237 δελφάκιον γενήσομαι]}{\textit{Test.} Suid.~δ~205 Δελφάκιον: ὁ μικρὸς χοῖρος. Ἀριστοφάνης· οἴμοι κακοδαίμων, δελφάκιον γέγονα. τοῦτο εἴρηκεν ἀφευθεὶς τὸν πρωκτόν. ἀντὶ τοῦ φλογισθείς. μετὰ γὰρ τὸ τυθῆναι τὰ δελφάκια φλογίζονται, ἵνα ψιλωθῶσι ; Ar.~\textit{Eccl.}~1021 οἴμοι Προκρούστης τήμερον γενήσομαι ; Ar.~fr.~608~K.-A. (=~594~Kock ; Eust.~1467.35) καὶ κύων ἀκράχολος / Ἑκάτης ἄγαλμα φωσφόρου γενήσομαι ; Ar.~fr.~957~K.-A. (dub. ; = adesp.~397~Kock ; Σ Pi.~\textit{P.}~2.125a) εὐθὺς ⟨δὲ⟩ Φοῖνιξ γίγνομαι· / τῆι μὲν δίδωμι χειρί, τῆι δὲ λαμβάνω ; Men.~fr.~264,~4–6~K.-A. (=~363~Kock ; Ὀργή ; Ath.~4.166a) καὶ βάψομαι / καὶ παρατιλοῦμαι, νὴ Δία, καὶ γενήσομαι / Κτήσιππος, οὐκ ἄνθρωπος, ἐν ὀλίγωι χρόνωι ; Eub.~fr.~122,~5~Kock (=~119~K.-A.) πεντηκοστολόγος γενήσομαι}
\lp{237 δελφάκιον]}{Ar.~\textit{Ach.}~785–787 κέρκον οὐκ ἔχει. / νεαρὰ γάρ ἐστιν· ἀλλὰ δελφακουμένα / ἑξεῖ μεγάλαν τε καὶ παχεῖαν κἠρυθράν ; Ar.~\textit{Lys.}~1061–1062 καὶ δελφάκιον ἦν τί μοι, / καὶ τοῦτο τέθυχ᾽, ὡς τὰ κρέ᾽ ἔδεσθ᾽ ἁπαλὰ καὶ καλά ; Σ Ar.~\textit{Thesm.}~236 [237] μετὰ γὰρ τὸ τυθῆναι τὰ δελφάκια φλογίζονται, ἵνα ψιλωθῶσιν ; Hsch.~δ~599 δελφάκιον· χοιρίδιον. οὕτως ἔλεγον καὶ τὸ γυναικεῖον (Ar. Thesm. 237) ; Anaxil.~fr.~12,~1~K.-A. (Κίρκη ; Ath.~9.374e) τοὺς μὲν ὀρειονόμους ὑμῶν ποιήσει δέλφακας ὑλιβάτας ; Alex.~fr.~124,~1–4~Kock (=~129~K.-A. ; Ath.~9.383c) ἧψε, μοὐδόκει, πνικτόν τι ὄψον δελφάκειον. γλ. ἡδύ γε. α. ἔπειτα προσκέκαυκε. γλ. μηδὲν φροντίσῃς· ἰάσιμον γὰρ τὸ πάθος ἐστί ; Ath.~9.374e–375a περὶ δὲ τῆς ἡλικίας τοῦ ζῴου Κρατῖνός φησιν ἐν Ἀρχιλόχοις· ἤδη δέλφακες, χοῖροι δὲ τοῖσιν ἄλλοις ; Alciphr.~2.29 [III~73] Ἡ ὗς ἡ πρῴην ἐπίτεξ εἶναι νομιζομένη ἀρτίως τέτοκε, καὶ ἔχω δελφάκων ἀφθονίαν ; Σ Ar.~\textit{Ach.}~739 Τοὺς παρ᾽ ἡμῖν νῦν λεγομένους δέλφακας, λέγω δὴ τοὺς μικροὺς σῦς, οὓς οἱ ἀρχαῖοι χοίρους ἐκάλουν. … ἐκάλουν δὲ τοὺς μεγάλους χοίρους δέλφακας ; Hsch.~δ~600 δελφακοῦσθαι· τελειοῦσθαι τὰς ὗς (Ar. Ach. 786) ; Antiatt. p.~88,~31~Bekker Δέλφακα: τὴν θήλειάν φασι δεῖν λέγειν, ἀδιαφόρως τῶν παλαιῶν τιθέντων ; Suid.~δ~204 Δέλφακας καὶ Δελφάκιον: τοὺς μεγάλους χοίρους, οἱ δὲ καὶ τοὺς μικρούς ; Plu.~\textit{Mor.}~674b (Quaest. conv.~5.1.2) τὴν ὗν τὴν Παρμένοντος οὕτως ἐθαύμασαν, ὥστε παροιμιώδη γενέσθαι ; cf. App. Prov.~2.87 Εὖ μὲν, ἀλλ᾽ οὐδὲν πρὸς τὴν Παρμένοντος ὗν ; Ar.~\textit{Thesm.}~222 τί κέκραγας; ἐμβαλῶ σοι πάτταλον ; Ar.~\textit{Thesm.}~239 ἐπίκυπτε· τὴν κέρκον φυλάττου νυν ἄκραν ; Ar.~fr.~236~K.-A. (=~224~Kock ; Δαιταλῆς ; Ath.~9.368d) καὶ δελφακίων ἁπαλῶν κωλαῖ καὶ χναυμάτια πτερόεντα ; Epicr.~fr.~6,~4–5~K.-A. (Ἔμπορος ; Ath.~14.655f) οὐκ Ἦλις, ἔνθα δελφάκων ἐγὼ κρέα / κάλλιστ᾽ ὄπωπα πυρὸς ἀκμαῖς ἠνθισμένα ; Luc.~\textit{Sat.}~23 τὸν δὲ ὗν καὶ τὸν ἔλαφον καὶ τὰ δελφάκια μεταξὺ ὀπτώμενα ; Ath.~14.656e–657a δελφάκια δὲ Αἰσχίνης εἴρηκεν ἐν Ἀλκιβιάδῃ οὕτως· ὥσπερ αἱ καπηλίδες τὰ δελφάκια τρέφουσιν … δέλφακα δὲ ἀρσενικῶς εἴρηκεν Πλάτων ; Hsch.~α~5756 †ἀπάλιον· θῦμα, δελφάκιον ; Hsch.~φ~625 φλογίδες· κρέα ὀπτά. διὰ τὸ φλογίζεσθαι (Stratt. fr. 12, 2. Archipp. fr. 10, 3 K.–A.) ; Poll.~7.187 καὶ ὕες σύες, χοῖροι χοιρίδια, δελφάκια, καὶ τὰ ὀψίγονα καὶ σμικρὰ μετάχοιρα ; Phot.~δ~171 δελφάκιον· χοῖρος}
\lp{238 ἐνεγκάτω τις ἔνδοθεν]}{Ar.~\textit{Ach.}~805 ἐνεγκάτω τις ἔνδοθεν τῶν ἰσχάδων ; Ar.~\textit{Nub.}~1490 ἐμοὶ δὲ δᾷδ᾽ ἐνεγκάτω τις ἡμμένην ; Ar.~\textit{Vesp.}~860 ἀλλ᾽ ὡς τάχιστα πῦρ τις ἐξενεγκάτω ; Ar.~\textit{Plut.}~1194 ἀλλ᾽ ἐκδότω τις δεῦρο δᾷδας ἡμμένας ; Cratin.~fr.~271~K.-A. (=~251~Kock ; Ὧραι ; Poll.~10.76) μῶν βδελυγμία σ᾽ ἔχει; / πτερὸν ταχέως τις καὶ λεκάνην ἐνεγκάτω ; Ar.~\textit{Vesp.}~529 ἐνεγκάτω μοι δεῦρο τὴν κίστην τις ὡς τάχιστα ; Ar.~\textit{Pax}~1149 κἀξ ἐμοῦ δ᾽ ἐνεγκάτω τις τὴν κίχλην καὶ τὼ σπίνω ; Ar.~\textit{Ran.}~1304 ἐνεγκάτω τις τὸ λύριον ; Ar.~\textit{Plut.}~228 τῶν ἔνδοθέν τις εἰσενεγκάτω λαβών ; Pl.~\textit{Phd.}~116d καὶ ἐνεγκάτω τις τὸ φάρμακον, εἰ τέτριπται}
\lp{238 δᾷδ᾽ ἢ λύχνον]}{Ar.~\textit{Nub.}~543 οὐδ᾽ εἰσῇξε δᾷδας ἔχουσ᾽, οὐδ᾽ ἰοὺ ἰοὺ βοᾷ ; Ar.~\textit{Vesp.}~1329–1331 ὑμᾶς / ὦ πόνηροι ταυτῃὶ τῇ / δᾳδὶ φρυκτοὺς σκευάσω ; Ar.~\textit{Lys.}~828 ἀλλ᾽ ἀπεψιλωμένον / τῷ λύχνῳ ; Ar.~\textit{Nub.}~18 ἅπτε παῖ λύχνον ; Ar.~\textit{Lys.}~1217–1218 μῶν ἐγὼ τῇ λαμπάδι / ὑμᾶς κατακαύσω ; Ar.~\textit{Plut.}~1052 τὴν δᾷδα μή μοι πρόσφερ᾽ ; Hsch.~δ~53 δαΐδες· λαμπάδες. λύχνοι ; Hsch.~λ~1477 λύχνος· ἡ λαμπάς, ἡ δαΐς, παρ᾽ Ὁμήρῳ (τ 34). καὶ ἰχθῦς κάλλιστος ; Ar.~\textit{Pax}~690–692 ὅτι τυγχάνει λυχνοποιὸς ὤν. πρὸ τοῦ μὲν οὖν / ἐψηλαφῶμεν ἐν σκότῳ τὰ πράγματα, / νυνὶ δ᾽ ἅπαντα πρὸς λύχνον βουλεύσομεν ; Ar.~\textit{Pax}~1317 δᾷδάς τε φέρειν, καὶ πάντα λεὼν συγχαίρειν κἀπιχορεύειν ; Ar.~\textit{Thesm.}~101–103 ἱερὰν … δεξάμεναι / λαμπάδα κοῦραι ; Ar.~\textit{Eccl.}~978 τοῦ δαὶ δεόμενος δᾷδ᾽ ἔχων ἐλήλυθας ; Ar.~\textit{Eccl.}~1150 ἔχω δέ τοι καὶ δᾷδα ταυτηνὶ καλῶς ; Ar.~\textit{Plut.}~423–425 ἴσως Ἐρινύς ἐστιν ἐκ τραγῳδίας· / βλέπει γέ τοι μανικόν τι καὶ τραγῳδικόν. / ἀλλ᾽ οὐκ ἔχει γὰρ δᾷδας ; Ar.~\textit{Plut.}~1041 στέφανόν γέ τοι καὶ δᾷδ᾽ ἔχων πορεύεται ; Hsch.~δ~9 δᾷδα· λαμπάδα ; Hsch.~λ~1478 λυχνοῦχος· ὁ φανός, λαμπτήρ. οἱ δέ, ἐφ᾽ οὗ ὁ λύχνος ὀχεῖται ; Hsch.~σ~1860 στίλβη· λύχνος, ἐλλύχνιον (frr. com. Hermipp. 26, 2. Plat. 206, 2. Ar. 573 K.–A.) ; Diogenian.~6.27 Λύχνον ἐν μεσημβρίᾳ ἅπτειν: ἐπὶ τῶν παρὰ καιρόν τι ποιούντων}
\lp{239 ἐπίκυπτε]}{Ar.~\textit{Lys.}~1003 ᾇπερ λυχνοφορίοντες ἐπικεκύφαμες ; Anaxandr.~fr.~37~Kock (Πάνδαρος ; Ath.~3.105) οὐκ ἐπικεκυφὼς ὀρθός, ὦ βέλτιστ᾽, ἔσει· / αὕτη δὲ καριδοῖ τὸ σῶμα καμπύλη ; Hsch.~κ~2850 Κλαζομένιοι· οὕτω κωμῳδοῦνται. δοκοῦσι γὰρ ἐπικύπτοντες προστιλᾶν (com. ad. fr. 368 K.–A.) ; Suid.~α~3520 Ἀποσιμῶσαι. τὸ ἐπικύψαι καὶ τὴν πυγὴν προθεῖναι γυμνήν. Φίλιππος ; X.~\textit{Smp.}~2.22 ἐκεῖνος ταὐτὰ εἰς τὸ ἔμπροσθεν ἐπικύπτων μιμεῖσθαι τροχοὺς ἐπειρᾶτο ; Arist.~\textit{HA}~522b18 ὁ δὲ βδάλλων ὀρθὸς ἕστηκεν, μικρὸν ἐπικύπτων ; Hsch.~ε~4912 ἐπικύψ· ἐπικύψας}
\lp{239 τὴν κέρκον]}{Ar.~\textit{Ach.}~785–787 κέρκον οὐκ ἔχει. / νεαρὰ γάρ ἐστιν· ἀλλὰ δελφακουμένα / ἑξεῖ μεγάλαν τε καὶ παχεῖαν κἠρυθράν ; Ar.~\textit{Pax}~1054–1055 ἡ κέρκος ποιεῖ / καλῶς; / καλῶς δῆτ᾽ ὦ πότνι᾽ Εἰρήνη φίλη ; Eub.~fr.~127~K.-A. (=~130~Kock ; Clem.Al.~\textit{Strom.}~7.6.30) αὐτοῖς δὲ τοῖς θεοῖσι τὴν κέρκον μόνην / καὶ μηρὸν ὥσπερ παιδερασταῖς θύετε ; Aesch.~\textit{PV}~496–497 κνίσῃ τε κῶλα συγκαλυπτὰ καὶ μακρὰν / ὀσφῦν πυρώσας δυστέκμαρτον ἐς τέχνην ; Herod.~5.45 ῥάκος καλύψαι τὴν ἀνώνυμον κέρκον ; Hsch.~κ~2333 κέρκος· θηρίδιον τὰς ἀμπέλους βλάπτον. καὶ οὐρά. καὶ ἀνδρεῖον αἰδοῖον. ἢ ἀλεκτρυών. ἢ ἀρουραῖος μῦς ; Hsch.~σ~172 σάννιον· τὸ αἰδοῖον, ἀντὶ τοῦ κέρκιον· παρὰ τὸ τῇ κέρκῳ σαίνειν (Cratin. fr. 490 K.–A.?). τὸ γὰρ αἰδοῖον ἔσθ᾽ ὅτε οὐρὰν ἔλεγον, ὡς Εὔπολις (fr. 471 K.–A.) ; \textit{idem} Phot.~σ~67 (σαννίον{ον} ; om. τὸ γὰρ αἰδοῖον … ὡς Εὔπολις) ; Soph.~\textit{Ant.}~1009–1011 κἄτυφε κἀνέπτυε, καὶ μετάρσιοι / χολαὶ διεσπείροντο, καὶ καταρρυεῖς / μηροὶ καλυπτῆς ἐξέκειντο πιμελῆς ; Hor.~\textit{Sat.}~1.2.45–46 accidit, ut cuidam testis caudamque salacem / demeterent ferro ; Hsch.~ο~1820 οὐρά· ἡ κέρκος. καὶ τὸ αἰδοῖον (Eupol. fr. 471 K.–A. vel Soph. fr. 1078 R.) ; Suid.~κ~1952 Κόλουρα: τὰ μὴ ἔχοντα κέρκον, ἢ τὰ κολοβὴν ἔχοντα οὐράν· ἅπερ ἐν ταῖς ἱερουργίαις οὐ θύεται ; Diogenian.~5.15 Ἡ κέρκος τῇ ἀλώπεκι μαρτυρεῖ: ἐπὶ τῶν δεικνυόντων ἀπὸ μικρᾶς πράξεως τὸ ἦθος ; \textit{idem} Erasm.~\textit{Adag.}~1.9.35 (n°~835) (id est Cauda de vulpe testificatur) ; Ar.~\textit{Eq.}~909 ἰδοὺ δέχου κέρκον λαγῶ ; Ar.~\textit{Eq.}~1031 ὃς κέρκῴ σαίνων σ᾽ ὁπόταν δειπνῇς ἐπιτηρῶν ; Hsch.~θ~151 θαυλέα· οὐρά, κέρκος ; Phot.~τ~335 τιτίς· βραχὺ ὀρνίθιον· σημαίνει καὶ τὸ γυναικεῖον αἰδοῖον· τιτὶς καὶ ἡ κέρκος ; Suid.~κ~1402 Κέρκος: ἡ τῶν ἑρπετῶν οὐρά, καὶ παντὸς ζώου ; Zen.~4.50 Κερκωπίζειν: ἡ παροιμία ἀπὸ τῶν προσσαινόντων τῇ κέρκῳ ζώων μετενήνεκται}
\lp{239 τὴν κέρκον φυλάττου νυν ἄκραν]}{Eur.~\textit{Phoen.}~1255–1257 μάντεις δὲ μῆλ᾽ ἔσφαζον, ἐμπύρους τ᾽ ἀκμὰς / ῥήξεις τ᾽ ἐνώμων ὑγρότητ᾽ ἐναντίαν / ἄκραν τε λαμπάδ᾽, ἣ δυοῖν ὅρους ἔχει ; Σ Eur.~\textit{Phoen.}~1257 τουτέστι τὸ ἄκρον τοῦ πυρὸς ἐσκόπουν, ὃ καλεῖται κέρκος … τὴν οὐρὰν τοῦ ἱερείου· ἦν γὰρ καὶ δι᾽ αὐτῆς μαντεία· καμπτομένη γὰρ ὑπὸ τοῦ πυρὸς δυσχέρειαν ἐδήλου τοῦ πολέμου}
\lp{239 φυλάττου]}{Ar.~\textit{Vesp.}~1343 ἔχου· φυλάττου δ᾽, ὡς σαπρὸν τὸ σχοινίον ; Orac. ap. Hdt.~7.148.3 εἴσω τὸν προβόλαιον ἔχων πεφυλαγμένος ἧσο / καὶ κεφαλὴν πεφύλαξο· κάρη δὲ τὸ σῶμα σαώσει}
\lp{239 ἄκραν]}{Hdt.~2.40.2 τὴν ὀσφὺν ἄκρην ; Ar.~\textit{Ach.}~638 εὐθὺς διὰ τοὺς στεφάνους ἐπ᾽ ἄκρων τῶν πυγιδίων ἐκάθησθε ; Men. ap. Clem.Al.~\textit{Strom.}~7.6.31 Μένανδρός τε τὴν ὀσφὺν ἄκραν πεποίηκεν ⟨καὶ⟩ τὴν χολὴν ὀστέα τὰ ἄβρωτα (φησί) τοῖς θεοῖς ἐπιθέντες αὐτοὶ τὰ ἄλλα ἀναλίσκουσιν ; Ar.~\textit{Lys.}~443 εἴ τἄρα νὴ τὴν Φωσφόρον τὴν χεῖρ᾽ ἄκραν}
\lp{240 ἐμοὶ μελήσει]}{Ar.~\textit{Thesm.}~1207 ἐμοὶ μελήσει ταῦτά γ᾽ ἢν ἅπαξ λυθῶ ; Ar.~\textit{Ach.}~932 ἐμοὶ μελήσει ταῦτ᾽ ; Ar.~\textit{Pax}~149 ἐμοὶ μελήσει ταῦτά γ᾽. ἀλλὰ χαίρετε ; Ar.~\textit{Pax}~1041 ἐμοὶ μελήσει ταῦτά γ᾽· ἀλλ᾽ ἥκειν ἐχρῆν ; Ar.~\textit{Thesm.}~1064 ἐμοὶ μελήσει ταῦτά γ᾽. ἀλλ᾽ ἄρχου λόγων ; Ar.~\textit{Plut.}~229 ἐμοὶ μελήσει τοῦτό γ᾽· ἀλλ᾽ ἀνύσας τρέχε ; Ar.~\textit{Pax}~1311 ἡμῖν μελήσει ταῦτά γ᾽· εὖ ποιεῖς δὲ καὶ σὺ φράζων ; Eur.~\textit{Ion}~35–36 τὰ δ᾽ ἄλλ᾽ — ἐμὸς γάρ ἐστιν, ὡς εἰδῇς, ὁ παῖς — / ἡμῖν μελήσει ; Orac. ap. Aristid.~\textit{Or.}~26.339 ἐμοὶ μελήσει ταῦτα καὶ λευκαῖς κόραις ; \textit{idem} App. Prov.~2.55 (add. βαρβάρων τινῶν στρατευομένων ἐπὶ Δελφοὺς ἐρωτώμενος ὁ θεὸς ταῦτα ἔφη), Suid.~ε~1060 (add. λέγουσιν, ὅτι βαρβάρων τινῶν στρατευομένων ἐπὶ Δελφούς, ἐρωτώμενος ὁ θεὸς ἔφησεν) ; X.~\textit{An.}~1.4.16 ἐμοὶ μελήσει ; Hsch.~μ~674 μέλει· φροντίζει ; Hsch.~μ~768 μέλει τάδε· ἐννοεῖ}
\lp{240 πλήν γ᾽ ὅτι]}{Ar.~\textit{Eq.}~27 νὴ Δία· πλήν γε περὶ τῷ δέρματι ; Ar.~\textit{Pax}~675–676 ψυχήν γ᾽ ἄριστος πλήν γ᾽ ὅτι / οὐκ ἦν ἄρ᾽ οὗπέρ φησιν εἶναι τοῦ πατρός ; Ar.~\textit{Nub.}~1429 ἡμῶν ἐκεῖνοι, πλήν γ᾽ ὅτι ψηφίσματ᾽ οὐ γράφουσιν ; Ar.~\textit{Ran.}~11 μὴ δῆθ᾽, ἱκετεύω, πλήν γ᾽ ὅταν μέλλω ᾽ξεμεῖν}
\lp{240 κάομαι]}{Ar.~\textit{Lys.}~9 ἀλλ᾽ ὦ Καλονίκη κάομαι τὴν καρδίαν ; Ar.~\textit{Lys.}~133–134 ἄλλ᾽ ἄλλ᾽ ὅ τι βούλει· κἄν με χρῇ διὰ τοῦ πυρὸς / ἐθέλω βαδίζειν· τοῦτο μᾶλλον τοῦ πέους ; Posidipp.~fr.~1,~8–10~K.-A. (Ἀναβλέπων ; Ath.~14.661f) ἐκ τοῦ πυρὸς / εἰς τὰς μαχαίρας ἦλθον· ὧν εἷς οὑτοσὶ / διὰ τῶν μαχαιρῶν τοῦ πυρός τ᾽ ἐλήλυθεν ; Hsch.~η~386 ἧλον· κατακεκαυμένον. λάγνον ; Moer.~κ~46 Hansen κλάειν καὶ κάειν σὺν τῷ α Ἀττικοί· μετὰ δὲ τοῦ ι Ἕλληνες ; Zen.~3.19 Διὰ μαχαιρῶν καὶ πυρὸς ῥίπτειν δεῖ; ἐπὶ τῶν παραβαλλομένων καὶ ῥιψοκίνδυνα ποιούντων ; \textit{idem} Suid.~δ~648 (ῥιπτεῖν χρή ; ἐπὶ τῶν παραβόλως τι ποιούντων) ; Diogenian.~8.45 Τὸν καπνὸν φεύγων, εἰς τὸ πῦρ ἐνέπεσον: ἐπὶ τῶν τὰ μικρὰ τῶν δεινῶν φευγόντων, καὶ εἰς μείζονα δεινὰ ἐμπιπτόντων ; \textit{idem} Tosi n°~671 (ἐνέπεσεν ; om. ἐπὶ τῶν …), Erasm.~\textit{Adag.}~1.5.5 (n°~405) (Καπνὸν γε φεύγων εἰς τὸ πῦρ περιέπεσον) ; Pl.~\textit{R.}~569b καί, τὸ λεγόμενον, ὁ δῆμος φεύγων ἂν καπνὸν δουλείας ἐλευθέρων εἰς πῦρ δούλων δεσποτείας ἂν ἐμπεπτωκὼς εἴη ; X.~\textit{Mem.}~1.3.9 οὗτος κἂν εἰς μαχαίρας κυβιστήσειε κἂν εἰς πῦρ ἅλοιτο ; Longus~1.14.1 κάομαι, καὶ ἐν σκιᾷ τοσαύτῃ κάθημαι ; Hsch.~κ~1141 κατακαίειν· καταφλέγειν ; Hsch.~φ~598 φλέξαι· κατακαῦσαι (Ι 653 v.l.) ; Diogenian.~6.68 Μὴ τὴν τέφραν φεύγων εἰς τὴν ἀνθρακιὰν ἐμπέσῃς}
\lp{241 οἴμοι τάλας]}{Ar.~\textit{Ach.}~174 οἴμοι τάλας μυττωτὸν ὅσον ἀπώλεσα ; Ar.~\textit{Av.}~62 οἴμοι τάλας ὀρνιθοθήρα τουτωί ; Ar.~\textit{Thesm.}~1038 τάλας ἐγὼ τάλας ; Ar.~fr.~477,~1~K.-A. (=~462~Kock) οἴμοι τάλας, τί μου στρέφει τὴν γαστέρα ; Pherecr.~fr.~56,~3~K.-A. (=~51~Kock) “οἴμοι τάλας”, φησίν, “χαράδρα κατελήλυθεν” ; Hermipp.~fr.~51~K.-A. (=~52~Kock ; Ath.~14.649c) οἴμοι τάλας, δάκνει, δάκνει, ἀπεσθίει μου τὴν ἀκοήν ; Antiph.~fr.~257,~1~K.-A. οἴμοι τάλας, ὡς ἄδικον ὅτ⌞αν⌟ ἡ μὲν τύχη / λίπῃ τινὸς τὴν τάξ⌞ιν⌟, ὁ δὲ τρόπος μένῃ ; Soph.~\textit{OT}~744 οἴμοι τάλας· ἔοικ᾽ ἐμαυτὸν εἰς ἀρὰς ; Soph.~\textit{Phil.}~416 οἴμοι τάλας. ἀλλ᾽ οὐχ ὁ Τυδέως γόνος ; Soph.~\textit{Phil.}~622 οἴμοι τάλας· ἦ κεῖνος, ἡ πᾶσα βλάβη ; Hsch.~ο~297 οἴμοι· ἐπίφθεγμα θρήνου ; Hsch.~τ~67 τάλας· ἄθλιος … ταλαίπωρος}
\lp{241 ὕδωρ ὕδωρ]}{Ar.~\textit{Ach.}~1174–1175 ὦ δμῶες οἳ κατ᾽ οἶκόν ἐστε Λαμάχου, / ὕδωρ ὕδωρ ἐν χυτριδίῳ θερμαίνετε ; Ar.~\textit{Lys.}~333–334 δημότισιν καομέναις / φέρουσ᾽ ὕδωρ βοηθῶ ; Xenarch.~fr.~7,~12–14~K.-A. (Πορφύρα ; Ath.~6.225c) καιρίαν δ᾽ εἰληφέναι / δόξας καταπίπτει καὶ λιποψυχεῖν δοκῶν / ἔκειτο μετὰ τῶν ἰχθύων. βοᾶι δέ τις / “ὕδωρ ⟨ὕδωρ⟩” ; Prop.~4.8.58 territa uicinas Teia clamat aquas ; Plu.~\textit{Rom.}~20.6 οἱ δ᾽ ὥσπερ ἐμπρησμῷ βοηθοῦντες ἐβόων ὕδωρ, ὕδωρ, καὶ συνέτρεχον πανταχόθεν ἀγγεῖα πλήρη κομίζοντες ; Sen.~\textit{Ep.}~17.3 cum aqua conclamata est, quomodo exeat, non quid efferat, quaerit ; [Quint.] Decl.~12.6 ut arma bello, ut aqua incendio inclamari publice solent ; Ar.~\textit{Thesm.}~508 ἄπελθ᾽ ἄπελθ᾽, ἤδη γὰρ ὦνέρ μοι δοκῶ ; Ar.~\textit{Thesm.}~514 λέων λέων σοι γέγονεν ; Ar.~\textit{Thesm.}~913–915 λαβέ με λαβέ με πόσι, περίβαλε δὲ χέρας. / φέρε σὲ κύσω. ἄπαγέ μ᾽ ἄπαγ᾽ ἄπαγέ με ; D.~23.22 δικάζειν δὲ τὴν βουλὴν τὴν ἐν Ἀρείῳ πάγῳ φόνου καὶ τραύματος ἐκ προνοίας καὶ πυρκαϊᾶς καὶ φαρμάκων ; Hsch.~π~4432 πυρκαϊά· ἐμπυρισμός (lex ap. Dem. 23, 22)}
\lp{241 ὦ γείτονες]}{Ar.~\textit{Nub.}~1322 ὦ γείτονες καὶ ξυγγενεῖς καὶ δημόται ; Ar.~\textit{Pax}~79 οἴμοι τάλας· ἴτε δεῦρο δεῦρ᾽ ὦ γείτονες ; Hes.~\textit{Op.}~343–345 τὸν δὲ μάλιστα καλεῖν, ὅς τις σέθεν ἐγγύθι ναίει· / εἰ γάρ τοι καὶ χρῆμ᾽ ἐγχώριον ἄλλο γένηται, / γείτονες ἄζωστοι ἔκιον, ζώσαντο δὲ πηοί ; Hes.~\textit{Op.}~346 πῆμα κακὸς γείτων, ὅσσον τ᾽ ἀγαθὸς μέγ᾽ ὄνειαρ ; Hor.~\textit{Epist.}~1.18.84 nam tua res agitur, paries cum proximus ardet ; \textit{idem} Tosi n°~523 (om. nam), Erasm.~\textit{Adag.}~3.6.71 (Tua res agitur quum proximus ardet paries ; add. … tum nobis esse cauendum quum vicini domus ardet) ; Juv.~3.198–199 iam poscit aquam, iam frivola transfert / Vcalegon, tabulata tibi iam tertia fumant ; Ov.~\textit{Rem.}~625–626 Proximus a tectis ignis defenditur aegre; / Utile finitimis abstinuisse locis ; X.~\textit{Mem.}~2.2.12 καὶ τῷ γείτονι βούλῃ σὺ ἀρέσκειν, ἵνα σοι καὶ πῦρ ἐναύῃ, ὅταν τούτου δέῃ, καὶ ἀγαθοῦ τέ σοι γίγνηται συλλήπτωρ καί, ἄν τι σφαλλόμενος τύχῃς, εὐνοϊκῶς ἐγγύθεν βοηθῇ σοι ; Verg.~\textit{Aen.}~2.311–312 iam proxumus ardet / Ucalegon ; \textit{idem} Tosi n°~670 (Iam) ; Hsch.~α~890 ἀγχιτέρμων· γείτων (Eur. Rhes. 426). καὶ ἀγχίθυρον· ⟨γείτονα⟩ ; Hsch.~γ~276 γειτνιῶντες· πλησιόχωροι ; Hsch.~γ~280 γείτον᾽· ἐγγύς. καὶ γεῖτος γράφεται αὐτό ; Hsch.~ο~788 ὅμορος· πλησιόχωρος … γειτνιῶν ; Hsch.~π~1741 περικτίονες· περίοικοι, γείτονες, πέριξ οἰκοῦντες ; Zen.~5.21 Μυκώνιος γείτων ; App. Prov.~4.31 Ὀξύτερον οἱ γείτονες βλέπουσι τῶν ἀλωπέκων: ἐπὶ τῶν λανθόνειν μὴ δυναμένων ; Alciphr.~1.18 (I~15) Δυσμενὴς καὶ βάσκανος ὁ τῶν γειτόνων ὀφθαλμός, φησιν ἡ παροιμία}
\lp{242 πρὶν ἀντιλαβέσθαι πρωκτὸν ἕτερον τῆς φλογός]}{Ar.~\textit{Ach.}~923 κεἴπερ λάβοιτο τῶν νεῶν τὸ πῦρ ἅπαξ ; Ar.~\textit{Nub.}~1497 οἴμοι τίς ἡμῶν πυρπολεῖ τὴν οἰκίαν ; Hor.~\textit{Epist.}~1.18.85 et neglecta solent incendia sumere vires ; Σ Ar.~\textit{Thesm.}~242 ὡς εἰ ἔλεγεν οἰκίαν ; Ar.~\textit{Plut.}~1053–1054 ἐὰν γὰρ αὐτὴν εἷς μόνος σπινθὴρ λάβῃ / ὥσπερ παλαιὰν εἰρεσιώνην καύσεται ; Plu.~\textit{Crass.}~2.4 ἐξηγόραζε τὰ καιόμενα καὶ γειτνιῶντα τοῖς καιομένοις, διὰ φόβον καὶ ἀδηλότητα τῶν δεσποτῶν ἀπ᾽ ὀλίγης τιμῆς προϊεμένων ; Plu.~\textit{Mor.}~454e (De cohib. ira~4) καθάπερ οὖν τὴν φλόγα θριξὶ λαγῴαις ἀναπτομένην καὶ θρυαλλίσι καὶ συρφετῷ ῥᾴδιόν ἐστιν ἐπισχεῖν· ἐὰν δ᾽ ἐπιλάβηται τῶν στερεῶν καὶ βάθος ἐχόντων, ταχὺ διέφθειρε ; Liv.~7.30.12 cum uideremus finitimum populum nefario latrocinio Samnitium peti et, ubi conflagrassent Sidicini, ad nos traiecturum illud incendium esse ; Eur.~\textit{Hel.}~1162 τείχεα δὲ φλογμὸς ὥστε Διός ἐπέσυτο φλόξ ; Hsch.~φ~644 φλόξ· ἡ τοῦ πυρός (Θ 135). καὶ ἄνθος τι (Mel. AP 4, 1, 51 = HE 3976)}
\lp{242 πρὶν ἀντιλαβέσθαι]}{Th.~7.77.6 καὶ ἢν ἀντιλαβώμεθά του φιλίου χωρίου τῶν Σικελῶν ; Ar.~\textit{Pax}~485 ἅπασιν ἡμῖν αὖθις ἀντιληπτέον ; Ar.~\textit{Lys.}~596 κἂν τούτου μὴ ᾽πιλάβηται ; Hsch.~α~5428 ἀντιλάζυσθαι· ἀντιλαβέσθαι (Eur. Or. 753) ; Hsch.~α~5429 ἀντιλαμβάνεται· βοηθεῖ. ἢ κωλύει}
\lp{242 πρωκτόν]}{Ar.~\textit{Ach.}~119 ὦ θερμόβουλον πρωκτὸν ἐξυρημένε ; Σ Ar.~\textit{Ach.}~119 οὗτος οὖν σκώπτων Εὐριπίδην προσέθηκε πρωκτὸν παρὰ προσδοκίαν ; Hsch.~ε~6767 «εὕδοντι δ᾽ αἱρεῖ πρωκτός»· παροιμία ἀπὸ τῆς εὕδοντι κύρτος αἱρεῖ. λέγουσι δὲ ὅτι, ὅταν τοὺς κύρτους καθῶσιν, αὐτοὶ μὲν καθεύδουσιν, ὁ δὲ κύρτος αἱρεῖ τοὺς ἰχθῦς. Κρατῖνος Ἀρχιλόχοις (fr. 3 K.–A.) ; Hsch.~α~7248 Ἀριστόδημος· Ἀριστόδημον οἱ κωμικοὶ τὸν πρωκτόν (Cratin. fr. 160 K.–A. Ar. fr. 242 K.–A.) ; Hsch.~ε~3839 Ἐξήκεστος· ἡταιρηκώς. ὅθεν καὶ τοὺς πρώκτους ὁμωνύμως Ἐξηκέστους ἔλεγον (com. ad. fr. 337 K.–A.) ; Hsch.~θ~267 Θεόδωρος· Θεοδώρους ἔλεγον οἱ κωμικοὶ τοὺς πρώκτους, ἀπὸ Θεοδώρου τινὸς οὐκ εὖ τῇ ἑαυτοῦ ὥρᾳ χρησαμένου ; Hsch.~ι~426 ἰθών· πυγή. λαγαρός. καὶ πρωκτός. ἄλλοι δὲ τὰ γλουτά ; Hsch.~π~4371 πύννος· ὁ πρωκτός ; Ar.~\textit{Thesm.}~1119–1120 ἀτὰρ εἰ τὸ πρωκτὸ δεῦρο περιεστραμμένον, / οὐκ ἐπτόνησά σ᾽ αὐτὸ πυγίζεις ἄγων ; Ar.~\textit{Thesm.}~1124 τὴ σανίδο τρήσας ἐξόπιστο πρώκτισον ; Hsch.~π~1833 περιπρωκτιῶσα· τρυφερευομένη ἐπὶ τῇ πυγῇ (com. ad. fr. 410 K.–A.) ; Hsch.~τ~899 †Τιμάναξ†· ὁ πρωκτός. ὁ δὲ αὐτὸς καὶ Θεόδωρος}
\lp{243 θάρρει. — τί θαρρῶ]}{Plaut.~\textit{Curc.}~553–554 bellator, vale. — Quid valeam? ; Plaut.~\textit{Curc.}~599 Phaedrome, propera. — Quid properem? ; Aesch.~fr.~310,~2~R. (Ath.~9.375e) ἕψου, μηδὲ λυπηθῇς πυρί ; Ar.~\textit{Ran.}~302 θάρρει· πάντ᾽ ἀγαθὰ πεπράγαμεν ; Plaut.~\textit{Curc.}~182–183 Tace. — Quid, taceam? ; Ar.~\textit{Ach.}~830 θάρρει Μεγαρίκ᾽· ἀλλ᾽ ἧς τὰ χοιρίδι᾽ ἀπέδου ; Ar.~\textit{Thesm.}~144–145 τί φῄς; τί σιγᾷς; ἀλλὰ δῆτ᾽ ἐκ τοῦ μέλους / ζητῶ σ᾽ ; Ar.~\textit{Eccl.}~796 θάρρει, καταθήσεις, κἂν ἔνης ἔλθῃς ; Ar.~\textit{Plut.}~328 θάρρει· βλέπειν γὰρ ἄντικρυς δόξεις μ᾽ Ἄρη ; Ar.~\textit{Plut.}~452 θάρρει· μόνος γὰρ ὁ θεὸς οὗτος οἶδ᾽ ὅτι ; Ar.~\textit{Plut.}~1091 θάρρει, μὴ φοβοῦ ; Ar.~\textit{Plut.}~1188 θάρρει· καλῶς ἔσται γάρ, ἢν θεὸς θέλῃ ; Hsch.~θ~117 θάρσει· θάρρει. θράσει (Ζ 126) ; Hsch.~θ~118 θαρρῶ· θρασύνομαι}
\lp{243 καταπεπυρπολημένος]}{Ar.~\textit{Thesm.}~727 καὶ καταίθειν τὸν πανοῦργον πυρπολεῖν θ᾽ ; Ar.~\textit{Vesp.}~1079 τῷ καπνῷ τύφων ἅπασαν τὴν πόλιν καὶ πυρπολῶν ; Ar.~\textit{Av.}~1580 τυρὸν φερέτω τις· πυρπόλει ; Eur.~\textit{Supp.}~640 ὃν Ζεὺς κεραυνῷ πυρπόλῳ καταιθαλοῖ ; Aesch.~\textit{Pers.}~1056 καί μοι γενείου πέρθε λευκήρη τρίχα ; Luc.~\textit{Podagr.}~21–22 σπλάγχνων δ᾽ ἐπ᾽ αὐτῶν διάπυρον τρέχει κακόν, / δίναισι φλογμῶν σάρκα πυρπολούμενον ; Hdt.~8.50.1 ἀγγέλλων ἥκειν τὸν βάρβαρον ἐς τὴν Ἀττικὴν καὶ πᾶσαν αὐτὴν πυρπολέεσθαι ; Plb.~5.19.8 αὐτόν τε τὸν τόπον τοῦτον πάντα κατεπυρπόλει καὶ διέφθειρε τοὺς ἐν αὐτῷ καρπούς ; Hsch.~κ~1290 κατανθρακῶσαι· κατακαῦσαι. κατοπτῆσαι (Aesch. fr. 281(d), 4 R.) ; Hsch.~π~4452 πυρπολεῖ· ἅπτει, καίει, πυρσεύει (Ar. Nub. 1497) ; Hsch.~π~4453 πυρπολέοντας· πυρπολουμένους, ἢ περὶ τὴν πυρὰν ἀναστρεφομένους (κ 30) ; Hsch.~π~4454 πυρπολούμενος· ἐμπυριζόμενος, παρὰ τὸ πολ[εμ]εῖν †ἐν αὐτῷ† τὸ πῦρ (4. Macc ; Synag.~π~812~Cunningham πυρπολεῖται· καίεται, ἐμπυρίζεται ; Ar.~\textit{Thesm.}~1215 ὀρτῶς δὲ συβήνη ᾽στί· καταβηνῆσι γάρ ; Eur.~\textit{Hel.}~767 τὰ Ναυπλίου τ᾽ Εὐβοικὰ πυρπολήματα ; Palaeph.~52 πλειόνων τῆς περιοικίδος καταπυρποληθέντων μερῶν ; Hsch.~κ~1101 κάταιθε· κατάκαιε (Eur. Or. 1620) ; Hsch.~π~4468 πυρσεύει· δᾳδουχεῖ, λάμπει, ἐξάπτει, ἅπτει πῦρ (Prov. 16, 28?)}
\lp{244 ἀλλ᾽ οὐκ ἔτ᾽ οὐδὲν πρᾶγμά σοι]}{Ar.~\textit{Eccl.}~462 οὐδὲ στένειν τὸν ὄρθρον ἔτι πρᾶγμ᾽ ἆρά μοι ; Ar.~\textit{Pax}~244 τουτὶ μὲν ἄνδρες οὐδὲν ἡμῖν πρᾶγμά πω ; Ar.~\textit{Ran.}~1215 ἀλλ᾽ οὐδὲν ἔσται πρᾶγμα ; Ar.~\textit{Eccl.}~670 τί γὰρ αὐτῷ πρᾶγμα μάχεσθαι ; Eur.~\textit{Med.}~451 κἀμοὶ μὲν οὐδὲν πρᾶγμα ; Pl.~\textit{Euthphr.}~3c ἀλλὰ τὸ μὲν καταγελασθῆναι ἴσως οὐδὲν πρᾶγμα ; Ar.~\textit{Eq.}~1243 οἴμοι κακοδαίμων· οὐκέτ᾽ οὐδέν εἰμ᾽ ἐγώ ; Ar.~\textit{Vesp.}~1392–1393 πράγματ᾽ αὖ δεῖ καὶ δίκας / ἔχειν ; Ar.~\textit{Pax}~685 ἀλλ᾽ οὐκέτ᾽ αὐτῷ χρησόμεθ᾽ οὐδέν, ἀλλὰ νῦν ; Nicom. Com. fr.~1,~24~K.-A. (Εἰλείθυια ; Ath.~7.290e–291d) γεωμετρικῆι δὲ καὶ σοὶ πρᾶγμα τί}
\lp{244–245 τὰ πλεῖστα γὰρ ἀποπεπόνηκας. — φεῦ ἰοὺ τῆς ἀσβόλου]}{\textit{Test.} Ael.Dion.~α~186 Erbse ἄσβολος· θηλυκῶς ἡ ἄσβολος, οὐκ ἀρσενικῶς ὁ ἄσβολος. Ἀριστοφάνης ἐν Θεσμοφοριαζούσαις· «τὰ πλεῖστα γὰρ / ἀποπεπόνηκας. — φεῦ, ἰοῦ τῆς ἀσβόλου» ; \textit{idem} Phot.~α~2946 (οὐχὶ ἡ ἀσβόλη, οὐδὲ pro οὐκ ; om. ἐν ; φεῦ ἰοὺ Bekker : σφουΐου cod.), Synag. Σᵇ~α~2213~Cunningham (=~Bekker, Anecd. I~451) (οὐχὶ ἡ ἀσβόλη οὐδὲ pro οὐκ ; om. ἐν ; σφουΐου)}
\lp{245 ἀποπεπόνηκας]}{\textit{Test.} Suid.~α~3473 Ἀποπεπόνηκας: τῶν πόνων ἔληξας}
\lp{245 φεῦ ἰοὺ τῆς ἀσβόλου]}{\textit{Test.} Suid.~α~4132 Ἀσβόλη καὶ Ἄσβολος: ἡ αἰθάλη τοῦ πυρός, ἡ στακτή. φεῦ, ἰοὺ τῆς ἀσβόλου. καὶ ἐπὶ δοτικῆς τῇ ἀσβόλῳ ; Ar.~\textit{Lys.}~295b ἰοὺ ἰοὺ τοῦ καπνοῦ ; \textit{idem} Ar.~\textit{Lys.}~305b ; Ar.~\textit{Lys.}~312 φεῦ τοῦ καπνοῦ ; Hsch.~ι~757 ἰού· σχετλιαστικὸν ἐπίρρημα, ὡς τὸ φεῦ, καὶ τὸ παπαί ; Hsch.~φ~328 φεῦ· σχετλιασμοῦ ἐπίρρημα ; Ar.~\textit{Vesp.}~161 Ἄπολλον ἀποτρόπαιε τοῦ μαντεύματος ; Ar.~\textit{Av.}~1723 ὦ φεῦ φεῦ τῆς ὥρας τοῦ κάλλους ; Ar.~\textit{Lys.}~293–295a καὶ τὸ πῦρ φυσητέον, / μή μ᾽ ἀποσβεσθὲν λάθῃ πρὸς τῇ τελευτῇ τῆς ὁδοῦ. / φῦ φῦ ; Phot.~φ~127 φεῦ· σχετλιαστικὸν ἐπίρρημα ἢ ὀδυνηρὸν οἰμώγημα ; Suid.~ε~2807 βαβαί, βαβαιάξ· σχετλιαστικά, ἀντὶ τοῦ φεῦ, φεῦ. … φῦ, φῦ· ἐπὶ τοῦ φυσῶντος τὸ πῦρ ; Ar.~\textit{Ran.}~653 ἰοὺ ἰού. — τί ἔστιν; — ἱππέας ὁρῶ ; Phot.~ι~148 ἰού· σχετλιαστικὸν ἐπίρρημα· ἀντὶ τοῦ οἴμοι}
\lp{245 τῆς ἀσβόλου]}{Σ Ar.~\textit{Thesm.}~245 ἄσβολος καλεῖται ἡ αἰθάλη τοῦ πυρός ; Hsch.~ψ~243 ψόλος· καπνός, αἰθάλη, φλόξ, ἀσβόλη (Aesch. fr. 24 R.) ; \textit{Test.} Phryn.~\textit{PS}~28,~1 de Borries ἄσβολος (Aristoph. Thesm. 246)· θηλυκῶς λέγουσιν. Ἱππῶναξ (fr. 105) δ᾽ ἀρσενικῶς. τινὲς δὲ καὶ τὴν ἀσβόλην ; Moer.~α~17 Hansen (= p.~187,~17~Bekker) ἄσβολος Ἀττικοί· ἀσβόλη Ἕλληνες ; Luc.~\textit{DDeor.}~8.4 ὑπὸ τῆς ἀσβόλου ; Luc.~\textit{Tim.}~2 ἀναπλησθήσονται τῆς ἀσβόλου ; Hsch.~α~988 ἀδαλός· ἄσβολος ; Hsch.~α~1431 ἄζα· ἄσβολος. κόνις. παλαιότης. κόπρος ἐν ἀγγείῳ ὑπομείνασα ; Hsch.~α~1842 αἰθάλη· σποδός, τὸ ἐκ καμίνου μέλαν (Exod. 9, 8) ; Hsch.~α~1846 αἰθαλόεν· κεκαυμένον. λαμπρόν. ἐμπεπρησμένον (Β 415) ; Hsch.~α~1847 αἴθαλος· ἡ τοῦ πυρὸς ἀποφορά. καὶ ἡ ἐν τῷ ὀρόφῳ αἰθάλη ; Hsch.~α~7633 ἀσβολόεν· μέγα ὑψηλόν. μέλαν ; Hsch.~α~7634 ἄσβολος· ἀσβόλη ; Phryn.~\textit{Ecl.}~82 F. (=~90 Rutherford) Ἀσβόλη μὴ λέγε, ἀλλὰ ἄσβολος ; Eust.~1390.55 Αἴλιος δὲ Διονύσιος παρασημειούμενός τινα ὅπως κατὰ γένη προφέρονται, φησὶν οὕτω. … ἔτι θηλυκῶς, καὶ ἡ τύλη τὸ φύμα τὸ ἐπὶ τοῦ αὐχένος. καὶ ἡ ὄμφαξ. καὶ ἡ ἄσβολος. καὶ ἡ ψύλλα ; Alex.~fr.~103,~16~K.-A. (=~98~Kock ; Ἰσοστάσιον ; Ath.~13.568a) τὰς ὀφρῦς πυρρὰς ἔχει τις· ζωγραφοῦσιν ἀσβόλωι ; Semon.~fr.~7,~61 (Ael.~\textit{NA}~16.25) οὔτε πρὸς ἰπνὸν ἀσβόλην ἀλευμένη ἵζοιτ᾽ ; LXX \textit{La.}~4.8 Ἐσκότασεν ὑπὲρ ἀσβόλην τὸ εἶδος αὐτῶν ; Epict.~3.16.3 ἀμήχανον τὸν συνανατριβόμενον τῷ ἠσβολωμένῳ μὴ καὶ αὐτὸν ἀπολαῦσαι τῆς ἀσβόλης ; Luc.~\textit{Bis\,acc.}~8 ἐπιχρισάμενοι τῆς ἀσβόλου ; Luc.~\textit{DDeor.}~8.5 ὁ Ἥφαιστος καὶ οἱ δάκτυλοι αὐτοῦ ἀνάξιοι τῆς σῆς κύλικος καὶ ἀσβόλου μεστός ἐστι ; Ael.~\textit{VH}~2.15 τοὺς θρόνους ἀσβόλῳ ; Σ Od.~24.539 ψολόεντα] καυστικὸν, ἀσβόλην ποιοῦντα. ψόλον γὰρ τὴν ἀσβόλην λέγει ; Hsch.~α~1843 αἰθάλη· τέφρα, σποδός (Hippocr. mul. 1, 91 [8, 220, 4 L.]) ; Hsch.~θ~862 θυμάλωψ· ἡ λιγνυώδης τοῦ πυρὸς ἀναφορά. τινὲς δὲ ξύλον κεκαυμένον, δαλόν ; Hsch.~ψ~116 ψελός· αἰθαλός ; Hsch.~ψ~236 ψοθίον· αἰθαλῶδες}
\lp{246 αἰθὸς γεγένημαι]}{\textit{Test.} Suid.~αι~135 Αἰθός: ὁ καιόμενος ; Eust.~83.6 ἡ δ᾽ αὐτὴ κόμη ῥηθείη ἂν καὶ αἰθὴ ὡς ἐκ τοῦ αἰθός, οὗ χρῆσις ἐν τῷ, αἰθὸς γεγένημαι ; Eust.~1314.15 αἰθὸς γὰρ ὀξυτόνως ὁ οἷον ἀπηνθρακωμένος κατὰ τό, αἰθὸς γεγένημαι παρὰ τῷ κωμικῷ ; Eust.~1385.59 ὅτι ἀπὸ τοῦ αἴθω ἐξ οὗ καὶ ὁ παρὰ τῷ κωμικῷ αἰθὸς καὶ λέων αἴθων καὶ οἶνος αἶθοψ ; Eust.~992.58–60 οὕτως αἰθὸς πάναιθος, οὗ θηλυκὸν ἡ παναίθη ; Ar.~fr.~469,~2~K.-A. (=~453~Kock ; Πολύιδος ; Stob.~4.22b.43) Φαίδραν· ἐπὶ πῦρ δὲ πῦρ ἔοιχ᾽ ἥκειν ἄγων ; [Eur.]~\textit{Rh.}~990 τείχη τ᾽ Ἀχαιῶν ναυσὶν αἶθον ἐμβαλεῖν ; Hsch.~α~1851 αἴθειν· καίειν (Eur. Rhes. 78) ; Hsch.~α~1878 αἶθος· καῦμα (Eur. Rhes. 990) ; Hsch.~κ~1901 καυαλέον ἢ καυαλές· ὑπὸ Αἰολέων τὸ αἶθος. ἢ κατακεκαυμένον, ⟨κα⟩πυρόν, ξηρόν, θερμόν ; Suid.~αι~134 Αἶθος: ὁ ἐκκαύστης ; EM~32.30 (s.v.~Αἶθος) Ἰστέον ὅτι αἶθός ἐστιν ὁ αἴθων, παρὰ τὸ αἴθω, τὸ καίω· αἰθὸς δὲ, ὁ αἰθόμενος, τουτέστιν ὁ καιόμενος. οὕτω Φιλόξενος εἰς τὰ ῥηματικὰ αὐτοῦ ; Diogenian.~6.71 Μὴ πῦρ ἐπὶ πῦρ: ἀνθρακεὺς ἐμπιπράμενος ἐν τῇ καμίνῳ τοῦτο λέγει. Ἐπὶ δὲ τῶν κακοῖς συνεχέστερον περιπιπτόντων ; \textit{idem} Erasm.~\textit{Adag.}~1.2.8 (n°~108) (add. id est Ne incendium incendio addas. Usurpatur a Platone. … Diogenianus putat a carbonario quodam esse natum, qui invalescente incendio conflagrarit clamans interim : μὴ πῦρ ἐπὶ πῦρ) ; Ter.~\textit{Eun.}~85 Accede ad ignem hunc; iam calesces plus satis ; A.R.~3.1303–1304 ὧς ἄρα τώγε θοὴν φλόγα φυσιόωντες / ἐκ στομάτων ὁμάδευν, τὸν δ᾽ ἄμφεπε δήιον αἶθος ; Pl.~\textit{Lg.}~666a διδάσκοντες ὡς οὐ χρὴ πῦρ ἐπὶ πῦρ ὀχετεύειν εἴς τε τὸ σῶμα καὶ τὴν ψυχήν ; Hsch.~α~1863 αἴθινος καπνός· ὁ δριμύς ; Hsch.~α~1872 αἰθόλικας· φλυκταίνας (Hippocr. liqu. 6 [6, 130, 6 L.]) ; Hsch.~α~1880 αἴθουσαι· καίουσαι ; Zen.~5.69 Πῦρ ἐπὶ πῦρ: παροιμία, ἧς μέμνηται Πλάτων. Καὶ κακὸν ἐπὶ κακῷ ; \textit{idem} Hsch.~π~4415}
\lp{246 αἰθός]}{Hsch.~α~1861 αἰθὴς πέπλος· Κρατῖνος Κλεοβουλίναις (fr. 95 K.–A.)· «†ἐκβάλλοντες τοὺς αἰθεῖς πέπλους» ; Hsch.~α~1873 αἰθόν· λαμπρόν, πυρρόν, μέλαν ; Soph.~fr.~269a,~54~R. (Inachos) κάρβανος αἰθὸς ; Call.~\textit{Dian.}~68–69 ὃ δὲ δώματος ἐκ μυχάτοιο / ἔρχεται Ἑρμείης σποδιῆι κεχριμένος αἰθῆι ; Nic.~\textit{Ther.}~288 ἄλλοτε μὲν χροιῇ ψολόεις, ὁτὲ δ᾽ ἔμπαλιν αἰθός ; Bacchyl.~fr.~13,~6–7~Bergk (= fr.~4.69–70 Snell–Maehler) Ἐν δὲ σιδαροδέτοις πόρπαξιν αἰθᾶν / ἀραχνᾶν ἱστοὶ πέλονται ; Hsch.~α~1841 αἰθαί· μέλαιναι || καὶ πομφόλυγες ; Hsch.~α~1870 Αἰθίοψ· ὁ Λέσβιος. μέλας. ἢ τὸ κεράμιον ; Hsch.~α~1877 αἴθοπος· διαπύρου. μέλανος ; Phryn.~\textit{Ecl.}~85 F. (=~91 Rutherford) Αἴθαλος λέγε ἀρσενικῶς, ἀλλὰ μὴ αἰθάλη θηλυκῶς ; EM~33.5 (s.v.~Αἰθίοψ) Αἰθίοψ, παρὰ τὸ αἴθω τὸ καίω, καὶ τὸ ὢψ ὠπὸς, ὃ σημαίνει τὸν ὀφθαλμὸν καὶ τὸ πρόσωπον, ὁ κεκαυμένην ἔχων τὴν πρόσοψιν. λέγεται γὰρ ἐκεῖσε σφοδρότερον αἴθειν τὸν ἥλιον ; Opp.~\textit{H.}~1.133 χαλκέα καὶ κορακῖνον ἐπώνυμον αἴθοπι χροιῇ ; Opp.~\textit{H.}~1.510 τέρπονται, σάργων τε γένος καὶ κόσσυφος αἴθων ; Hsch.~α~1874 αἴθωνα· [μέλανα] πυρώδη (Λ 548) ; Hsch.~α~1902 αἴθων· μέλας (Ο 690)}
\lp{246 γεγένημαι]}{Ar.~\textit{Thesm.}~846 ἰλλὸς γεγένημαι προσδοκῶν· ὁ δ᾽ οὐδέπω ; Ar.~\textit{Eq.}~945 ἀνὴρ γεγένηται τοῖσι πολλοῖς τοὐβολοῦ ; Ar.~\textit{Eq.}~1255 ἀνὴρ γεγένησαι δι᾽ ἐμέ ; Ar.~\textit{Eccl.}~551 ἀτὰρ γεγένηται; — ναὶ μὰ Δί᾽ ; Ar.~\textit{Plut.}~148 δοῦλος γεγένημαι, διὰ τὸ μὴ πλουτεῖν ἴσως ; Ar.~\textit{Plut.}~652 ἃ νῦν γεγένηται ; Ar.~\textit{Plut.}~1043 πολιὰ γεγένησαι ταχύ γε νὴ τὸν οὐρανόν}
\lp{246 πάντα τὰ περὶ τὴν τράμιν]}{\textit{Test.} Suid.~τ~903 Τράμις: ὁ πρωκτός. Ἀριστοφάνης· αἰθὸς γεγένημαι πάντα τὰ περὶ τὸν τράμιν ; Hippon.~fr.~84~Bergk ἐκτίλλοι τις αὐτοῦ τὴν τράμιν ὑποργάσαι ; Luc.~\textit{Lexiph.}~2 τὰ ἀμφὶ τὴν τράμιν μαλακίζομαι ; Σ Ar.~\textit{Thesm.}~246 τράμις ὁ πρωκτός ; Hsch.~τ~1243 τράμις· τὸ ⟨τ⟩ρῆμα τῆς ἕδρας, ὁ ὄρ⟨ρ⟩ος (Hippon. fr. 114 a W. = 133 Degani). τινὲς ἔντερον. οἱ δὲ ἰσχίον ; Phot.~τ~415 τράμις· ὁ ὄρρος ἢ ἡ διαγραφή· ἢ τὸ τρῆμα τοῦ ἀρχοῦ ; Ar.~\textit{Thesm.}~256 ἴθι νυν κατάστειλόν με τὰ περὶ τὼ σκέλει ; Ar.~\textit{Thesm.}~647–648 ἰσθμόν τιν᾽ ἔχεις ὦνθρωπ᾽· ἄνω τε καὶ κάτω / τὸ πέος διέλκεις πυκνότερον Κορινθίων ; Telecl.~fr.~72~K.-A. (=~66~Kock ; Phot. p.~579,~8 Porson) τερπότραμις (-τράμις cod.)· ἡ τῶν ἀφροδισίων τέρψις. οὕτως Τηλεκλείδης ; Ruf.~\textit{Onom.}~101 καὶ ἡ διὰ μέσου γραμμὴ, τραμίς· οἱ δὲ ὄῤῥον ὀνομάζουσιν ; Σ Ar.~\textit{Ran.}~223 ὄρρον δέ, τὸν λεγόμενον ταῦρον, τὸ μεταξὺ τῶν διδύμων. Καλλίστρατος, τὴν ὀσφῦν καὶ τὸ ἱερὸν ὀστοῦν. Δίδυμος, τὴν τράμιν, οὐχ ὥς τινες τὸ ἰσχίον ; Σ Luc.~\textit{Lexiph.}~2 τὸ δὲ ἐδάρην τὸν ὄρρον ἀντὶ τοῦ ἀνεσύρην, ἐτρίβην τὸν ὑπόταυρον ἤτοι τὴν τράμιν ; Hsch.~δ~1392 διάτραμις· διερρωγὼς τὴν τράμιν, τὸν ἀπὸ ὀσχέου ἐπὶ τὰ ἰσχία τόπον (Stratt. fr. 84 K.–A.)· τράμις γὰρ τὰ ἰσχία ; Hsch.~ο~1315 ὀρροπύγιον· ὀρθοπύγιον (Ar. Nub. 158) ; Hsch.~ο~1316 ὄρρος· ἡ τράμις· οὐχ, ὥς τινες, τὸ ἰσχίον ; Poll.~2.173 τράμις, καὶ ὄρρος ; Poll.~2.184 λισπόπυγοι, ἐφ᾽ ᾧ μάλιστα Ἀθηναῖοι κωμῳδοῦνται. τὴν δὲ ὧδε ἔχουσαν, διατράμιν Στράττις ὠνόμασε ; Phot.~ο~526 ὄρρος· τράμις· — καὶ ὀρροπύγιον· οὕτως καὶ οἱ Ἴωνες· ὀρθοπύγιον δὲ παρ᾽ οὐδενὶ τῶν Ἑλλήνων ; EM~763.55 Τράμην, οἱ μὲν, τὸ τρῆμα τῆς ἕδρας· οἱ δὲ, τὸ ἀπὸ τοῦ τύπου μέχρι τῆς ἀρχῆς τοῦ βαλάνου χωρίον ; Erot. s.v.~τράμιν (p.~124 Klein = Archil. fr.~195~Bergk) Τράμιν, τὸν ὄρρον, ὅνπερ καὶ ἱππόταυρον (ἱπποταύριον Valckenaer) καλοῦμεν … μέμνηται καὶ Ἀρχίλοχος ; Ar.~\textit{Eq.}~424 ἀποκρυπτόμενος ἐς τὼ κοχώνα τοὺς θεοὺς ἀπώμνυν ; Ar.~\textit{Vesp.}~438 ὦ Κέκροψ ἥρως ἄναξ τὰ πρὸς ποδῶν Δρακοντίδη ; Ar.~\textit{Lys.}~92 δήλη ᾽στὶν οὖσα ταυταγὶ τἀντευθενί ; Hsch.~κ~3887 κοχώνα· τὰ ἰσχία, καὶ τὰ ὅμοια (Ar. Equ. 424) ; Hsch.~π~1763 περίνα· περίναιον. τὸ αἰδοῖον (Nic. ther. 586)}
\lp{247 μὴ φροντίσῃς]}{Ar.~\textit{Thesm.}~233 μὴ φροντίσῃ· ὡς εὐπρεπὴς φανεῖ πάνυ}
\lp{247 ἕτερος γὰρ αὐτὰ σπογγιεῖ]}{Ar.~\textit{Thesm.}~242 πρὶν ἀντιλαβέσθαι πρωκτὸν ἕτερον τῆς φλογός ; Ar.~\textit{Ran.}~482 ἀλλ᾽ οἶσε πρὸς τὴν καρδίαν μου σφογγιάν ; Ar.~\textit{Ran.}~487–490 πῶς δειλὸς ὅστις σφογγιὰν ᾔτησά σε; / οὐκ ἂν ἕτερός γ᾽ αὔτ᾽ ἠργάσατ᾽ ἀνήρ. / ἀλλὰ τί; / κατέκειτ᾽ ἂν ὀσφραινόμενος, εἴπερ δειλὸς ἦν· / ἐγὼ δ᾽ ἀνέστην καὶ προσέτ᾽ ἀπεψησάμην ; Pherecr.~fr.~58~K.-A. (=~53~Kock ; Ἐπιλήσμων ἢ Θάλαττα) τὸν ἱδρῶτα καὶ τὴν ἄρδαν ἀπ᾽ ἐμοῦ σπόγγισον ; \textit{Il.}~18.414 σπόγγῳ δ᾽ ἀμφὶ πρόσωπα καὶ ἄμφω χεῖρ᾽ ἀπομόργνυ ; Hsch.~π~1746 περικωνῆσαι· περιρ⟨ρ⟩ομβῆσαι. περιαγαγεῖν. σπογγίσαι (Ar. Vesp. 600). ἢ περιπισσῶσαι. Λάκωνες ; Ar.~\textit{Ach.}~463 δός μοι χυτρίδιον σφογγίῳ βεβυσμένον ; Ar.~\textit{Vesp.}~600 τὸν σφόγγον ἔχων ἐκ τῆς λεκάνης τἀμβάδι᾽ ἡμῶν περικωνεῖ ; Ar.~fr.~59~K.-A. (=~55~Kock ; Anagyros) ἀλλὰ πάντας χρὴ παραλοῦσθαι καὶ τοὺς σπόγγους ἐᾶν ; Crates fr.~17,~6–7~K.-A. (=~15~Kock ; Θηρία ; Ath.~6.268a) εἶθ᾽ ἀλάβαστος εὐθέως ἥξει μύρου / αὐτόματος ὁ σπόγγος τε καὶ τὰ σάνδαλα ; Aesch.~\textit{Ag.}~1329 βολαῖς ὑγρώσσων σπόγγος ὤλεσεν γραφήν ; D.~18.258 τὸ μέλαν τρίβων καὶ τὰ βάθρα σπογγίζων καὶ τὸ παιδαγωγεῖον κορῶν ; Hsch.~ε~1514 ἐκμάσσει· ⟨ἀποσπογγίζει⟩ ; Hsch.~π~2992 πόντια ῥάκ[κ]η· σπόγγοι, ἢ τὰ τούτων σπαράγματα (trag. ad. fr. 258 K.–S.) ; Hsch.~σ~1533 σπόγγος· σπογγίον (Aesch. Ag. 1329) ; Suid.~π~1354 Περιψεῖν: ἐκμάσσειν, σπογγίζειν. Ἀριστοφάνης· ἰδού, δέχου κέρκον λαγώ, ὥστε τὰς λήμας περιψεῖν ; Phot.~π~267 παραλοῦμαι· παροιμιακῶς· εἰώθεισαν γὰρ πρότερον ἐν τοῖς βαλανείοις οἱ πλούσιοι παραλούειν τοὺς πένητας ; Theopomp. Com. fr.~41,~1~K.-A. (=~40~Kock ; Παμφίλη ; Ath.~11.485b) σπόγγος, λεκάνη, πτερόν, λεπαστὴ πάνυ πυκνή ; \textit{Od.}~1.111 οἱ δ᾽ αὖτε σπόγγοισι πολυτρήτοισι τραπέζας ; Epict.~2.22.31 ὡς τὰ ὑποδήματα σπογγίζει τὰ ἑαυτοῦ ; Hsch.~α~8880 Ἀχίλλειων· Ἱππεῦσιν Ἀριστοφάνης (819) … τινὲς δὲ σπόγγον φασίν, οἷς τοὺς ὀφθαλμοὺς ἀποψῶνται ; Hsch.~π~2993 πόντι[ν]οι· σπόγγοι ; Hsch.~σ~1532 σπόγγια· σπογγώδη}
\lp{248 οἰμώξετ᾽ ἆρ᾽ εἴ τις]}{Ar.~\textit{Ach.}~325 δήξομἄρ᾽ ὑμᾶς ἐγώ ; Ar.~\textit{Nub.}~217 ἀλλ᾽ οὐχ οἷόν τε. — νὴ Δί᾽ οἰμώξεσθ᾽ ἄρα ; Ar.~\textit{Thesm.}~916–917 κλαύσετ᾽ ἄρα νὴ τὼ θεὼ / ὅστις σ᾽ ἀπάξει τυπτόμενος τῇ λαμπάδι ; Ar.~\textit{Plut.}~876 οἰμώξἄρα σύ ; Eur.~\textit{Hipp.}~1086 κλαίων τις αὐτῶν ἄρ᾽ ἐμοῦ γε θίξεται ; Ar.~\textit{Pax}~532 κλαύσἄρα σὺ ; Ar.~\textit{Thesm.}~211–212 τοῦτον μὲν μακρὰ / κλάειν κέλευ᾽ ; Ar.~\textit{Ran.}~178 ὡς σεμνὸς ὁ κατάρατος· οὐκ οἰμώξεται ; Luc.~\textit{DDeor.}~8.5 οἰμώξεται γάρ, ἤν τίς σε λυπεῖν θέλῃ ; Ar.~\textit{Pax}~466 οἰμώξεσθ᾽ οἱ Βοιωτοί ; Ar.~\textit{Av.}~1207 τί ποτ᾽ ἐστὶ τουτὶ τὸ κακόν; — οἰμώξει μακρά ; Eup.~fr.~328,~1~K.-A. (=~305~Kock) τίς οὑξεγείρας μ᾽ ἐστίν; οἰμώξει μακρά ; Alex.~fr.~115,~19~K.-A. (=~110~Kock ; Κρατεία ἢ Φαρμακοπώλης ; Ath.~3.107a) οἰμώξεται γὰρ νὴ Δί᾽. ἀλλ᾽ ἐγὼ σοφῶς ; Hsch.~ο~301 οἰμώζει· τὰ αὐτά}
\lp{248 τὸν ἐμὸν πρωκτὸν πλυνεῖ]}{Ar.~\textit{Ach.}~381 κἀκυκλοβόρει κἄπλυνεν, ὥστ᾽ ὀλίγου πάνυ ; Ar.~\textit{Vesp.}~603–604 πάντως γάρ τοι παύσει ποτὲ κἀναφανήσει / πρωκτὸς λουτροῦ περιγιγνόμενος τῆς ἀρχῆς τῆς περισέμνου ; Ar.~\textit{Plut.}~514 ἢ σκυτοτομεῖν ἢ πλινθουργεῖν ἢ πλύνειν ἢ σκυλοδεψεῖν ; Σ Ar.~\textit{Vesp.}~604 ὁ γὰρ πρωκτὸς πλυνόμενος περιγίνεται τῆς καθάρσεως, καὶ ἔτι μολύνεται … Εὐφρόνιος μὲν … καὶ γὰρ τὸ μέρος ἐκεῖνο πλυνόμενον ἔτι μολύνεται. Καλλίστρατος δέ φησι· παροιμία, πρωκτὸς λουτροῦ περιγίνῃ ; Hsch.~ε~5195 ἐπισμῇ· ἐπιτρίβει, ἀπὸ τοῦ σμήχειν ἢ τρίβειν ἢ πλύνειν· λοιδορεῖ, πλήττει ; Hsch.~π~2629 πλύνεται· βλασφημεῖται, λοιδορεῖται ; \textit{idem} Phot.~π~965 (λοιδορεῖται· αἰσχρῶς ὑβρίζεται) ; Hsch.~π~2631 πλύνον καταπλυντήριζε καὶ πλυνθήσομαι· Ἀττικοὶ ἐπὶ τῶν λοιδοριῶν λέγουσι (com. ad. fr. 413 K.–A.) ; Hsch.~π~4144 πρωκτὸς λουτροῦ περιγίνῃ· παροιμία, ὅταν τις μὴ δύνηται ἀπονίψασθαι, ἀλλ᾽ ἡ κοιλία αὐτῷ ἐπιφέρηται ; \textit{idem} Phot.~π~1450 (om. αὐτῷ ; add. ἐπὶ τῶν ἀνωφελῶν οὖν καὶ εἰκῇ πραττομένων ἐλέγετο), Suid.~π~2951 (lemma tantum), App. Prov.~4.61 (περιγίνεται ; om. παροιμία), Diogenian. Vind.~3.58 (περιγίνεται ; om. παροιμία ; add. λέγεται ἡ παροιμία ἐπὶ τῶν ἀνωφελῶς πραττομένων), Apostol.~14.78 (περιγίνεται ; ἐπὶ τῶν ἀνωφελῶν καὶ εἰκῇ πραττομένων ἐλέγετο· οἱ γὰρ παχύπρωκτοι καὶ προγάστορες οὐ δύνανται ἑαυτοὺς ἀπονίψασθαι εὐπετῶς), Erasm.~\textit{Adag.}~1.10.90 (n°~990) (id est Podex lotionem vincis) ; \textit{Test.} Moer.~π~79 Hansen (= p.~208,~15~Bekker) πλυνεῖς κατὰ τὴν πρώτην Ἀτθίδα· κναφεῖς κατὰ τὴν δευτέραν Ἀτθίδα ; Ar.~\textit{Eq.}~160 τί μ᾽ ὦγάθ᾽ οὐ πλύνειν ἐᾷς ; Ar.~\textit{Eq.}~1057 ἀλλ᾽ οὐκ ἂν μαχέσαιτο· χέσαιτο γάρ, εἰ μαχέσαιτο ; Hsch.~π~2628 πλύνει· καθαίρει (Ar. Plut. 166) ; Ar.~\textit{Nub.}~838 ὥσπερ τεθνεῶτος καταλόει μου τὸν βίον ; Ar.~\textit{Vesp.}~627 κἀγκεχόδασίν μ᾽ οἱ πλουτοῦντες ; Ar.~\textit{Av.}~68 Ἐπικεχοδὼς ἔγωγε Φασιανικός ; Ar.~\textit{Ran.}~308 ὁδὶ δὲ δείσας ὑπερεπυρρίασέ σου ; Hsch.~π~2630 πλυνοί· πύελοι, ἐν αἷς τὰς ἐσθῆτας ἔπλυνον· ἢ βόθρον, ὅπου πλύνουσι (Χ 153) ; Hsch.~ρ~506 ῥυπαίνει· σμήχει. καθαρίζει. πλύνει ; Hsch.~χ~34 χαλαίρυπος· ὁ τῶν πλυνομένων ἱματίων ῥύπος (Cratin. fr. 500 K.–A.)}
\lp{248 πλυνεῖ (labeur vain)]}{Zen.~6.48 Χαμαὶ ἀντλεῖς, Πλίνθον πλύνεις ; \textit{idem} Diogenian.~7.50 (Πλίνθον πλύνεις, Χαμαὶ ἀντλεῖς ; add. ἐπὶ τῶν ἀδυνάτων), App. Prov.~4.58 (Πλίνθον πλύνεις, Χαμαὶ ἀντλεῖς ; add. Φακὸν κόπτεις, Θάλασσαν ἀντλεῖς: ἐπὶ τῶν ἀδυνάτων), Greg. Cypr.~3.39 (Πλίνθον πλύνεις tantum ; add. ἐπὶ τῶν ἀδυνάτων), Plu.~\textit{De\,imposs.}~2 (Πλίνθον πλύνεις tantum), Erasm.~\textit{Adag.}~1.4.48 (n°~348) (Πλίνθον πλύνεις, id est Laterem lavas) ; Σ Ar.~\textit{Vesp.}~279 παροιμία δέ ἐστιν ἐπὶ τῶν ἀδυνάτων γενέσθαι τὸ λίθον ἕψεις, ὡς καὶ τὸ πλίνθον πλύνεις, καὶ χύτραν ποικίλλεις, καὶ εἰς ὕδωρ γράφεις, καὶ Αἰθίοπα λευκαίνεις, καὶ κατὰ θαλάττης σπείρεις, καὶ τὰ τοιαῦτα ; Phot.~ο~363 ὄνου πόκαι· ἐπὶ τῶν ἀνηνύτων καὶ τῶν μὴ ὄντων λέγεται ἡ παροιμία ὑπὸ τῶν Ἀττικῶν· ὥσπερ αἱ τοιαῦται· „πλίνθον πλύνειν“· „ἀσκὸν τίλλειν“· „χύτραν ποικίλλειν“ ; \textit{idem} Phot.~ο~360 (ὄνου πόκοι· παροιμία ἐπὶ τῶν ἀνηνύτων ; χύτραν πλύνειν), Suid.~ο~399 (Ὄνου πόκαι: ἐπὶ τῶν ἀνηνύτων tantum) ; \textit{AP}~11.428 (Ps.-Lucien) εἰς τί μάτην νίπτεις δέμας Ἰνδικόν; ἴσχεο τέχνης· / οὐ δύνασαι δνοφερὴν νύκτα καθηλιάσαι ; Luc.~\textit{Ind.}~28 κατὰ τὴν παροιμίαν Αἰθίοπα σμήχειν ἐπιχειρῶ ; Aesop.~13 Halm (=~11 Chambry) πάντα μὲν αὐτῷ προσῆγε τὰ ῥύμματα, πᾶσι δὲ λουτροῖς ἐπειρᾶτο φαιδρύνειν· καὶ τὸ μὲν χρῶμα μεταβαλεῖν οὐκ ἔσχε, νοσεῖν δὲ τῷ πονεῖν παρεσκεύασεν ; Zen.~1.46 Αἰθίοπα σμήχων· ἐπὶ τῶν μάτην ἀνήνυτον ἔργον ἐπιτελούντων ; \textit{idem} Diogenian.~1.45 (σμήχεις ; δῆλον τὸ τοιοῦτον. Ὅμοια, Χύτραν ποικίλλεις, ᾠὸν τίλλεις), Plu.~\textit{De\,imposs.}~7 (σμήχεις tantum), Suid.~αι~125 (σμήχειν ; ἐπὶ τῶν μάτην πονούντων), Tosi n°~1893 (σμήχεις), Erasm.~\textit{Adag.}~1.4.50 (n°~350) (σμήχεις, id est Aethiopem lavas ; add. et Αἰθίοπα λευκαίνεις, id est Aethiopem dealbas) ; Hsch.~π~2601 πλινθεύεται· ἐξαπατᾶται. ἐπὶ ἀναισθήτων ; LXX \textit{Je.}~13.23 … ἀλλάξεται Αἰθίοψ τὸ δέρμα αὐτοῦ καὶ πάρδαλις τὰ ποικίλματα αὐτῆς}
\lp{249 Ἀγάθων, ἐπειδὴ]}{Ar.~\textit{Thesm.}~218 Ἀγάθων σὺ μέντοι ξυροφορεῖς ἑκάστοτε}
\lp{249 σαυτὸν ἐπιδοῦναι]}{Ar.~\textit{Eq.}~739–740 σαυτὸν δὲ λυχνοπώλαισι καὶ νευρορράφοις / καὶ σκυτοτόμοις καὶ βυρσοπώλαισιν δίδως ; Ar.~\textit{Thesm.}~213 ἄγε νυν ἐπειδὴ σαυτὸν ἐπιδίδως ἐμοί ; Ar.~\textit{Plut.}~781 ἄκων ἐμαυτὸν τοῖς πονηροῖς ἐπεδίδουν ; Ar.~\textit{Thesm.}~204–205 δοκῶν γυναικῶν ἔργα νυκτερείσια / κλέπτειν ὑφαρπάζειν τε θήλειαν Κύπριν}
\lp{249 φθονεῖς]}{Ar.~\textit{Thesm.}~252 λαμβάνετε καὶ χρῆσθ᾽· οὐ φθονῶ ; Eur.~\textit{HF}~333 κοσμεῖσθ᾽ ἔσω μολόντες· οὐ φθονῶ πέπλων ; Eur.~\textit{IT}~503 τί δὲ φθονεῖς τοῦτο; ἦ φρονεῖς οὕτω μέγα ; Eur.~\textit{Med.}~63 τί δ᾽ ἔστιν, ὦ γεραιέ; μὴ φθόνει φράσαι ; Aesch.~\textit{PV}~626 ἀλλ᾽ οὐ μεγαίρω τοῦδέ σοι δωρήματος ; Hsch.~φ~424 φθονεῖς· ζηλοῖς, βασκαίνεις (α 346) ; Eur.~fr.~703~N.² (\textit{Telephus}) μή μοι φθονήσητ᾽, ἄνδρες Ἑλλήνων ἄκροι, / εἰ πτωχὸς ὢν τέτληκ᾽ ἐν ἐσθλοῖσιν λέγειν ; \textit{Od.}~11.380–381 οὐκ ἂν ἐγώ γε / τούτων σοι φθονέοιμι καὶ οἰκτρότερ᾽ ἄλλ᾽ ἀγορεύειν ; Pl.~\textit{Thg.}~125a ἔπειτα φθονεῖς τε αὐτῷ καὶ οὐκ ἐθέλεις πέμπειν ; Hsch.~α~335 ἄγασθαι· φθονεῖν ; Hsch.~β~290 βασκαίνει· λυπεῖ. φθονεῖ. μέμφεται (Dem. 8, 19) ; Hsch.~φ~425 φθονερός· ζηλωτής, ζηλότυπος}
\lp{250 ἀλλ᾽ … γοῦν]}{Ar.~\textit{Lys.}~877 μὴ δῆτ᾽, ἀλλὰ τῷ γοῦν παιδίῳ ; Ar.~\textit{Thesm.}~710 ἀλλ᾽ οὖν ἥκεις γ᾽ ὅθεν οὐ φαύλως γ᾽ ; Hsch.~γ~865 γοῦν· τοίνυν}
\lp{250 ἱμάτιον]}{Ar.~\textit{Thesm.}~261 φέρ᾽ ἔγκυκλον· — τουτὶ λάβ᾽ ἀπὸ τῆς κλινίδος ; Ar.~\textit{Thesm.}~214 ἀπόδυθι τουτὶ θοἰμάτιον. — καὶ δὴ χαμαί ; Ar.~\textit{Eccl.}~535–536 εἶτ᾽ οὐ τὸ σαυτῆς ἱμάτιον ἐχρῆν σ᾽ ἔχειν; / ἀλλ᾽ ἔμ᾽ ἀποδύσασ᾽ ἐπιβαλοῦσα τοὔγκυκλον ; Ar.~\textit{Lys.}~113 ἔγωγ᾽ ἂν ⟨οὖν⟩ κἂν εἴ με χρείη τοὔγκυκλον ; Ar.~\textit{Thesm.}~499 ὡς ἡ γυνὴ δεικνῦσα τἀνδρὶ τοὔγκυκλον ; Ar.~\textit{Eccl.}~332–333 τὸ κροκωτίδιον ἀμπισχόμενος οὑνδύεται. / τὸ δ᾽ ἱμάτιόν σου ποῦ ᾽στιν; — οὐκ ἔχω φράσαι ; Hsch.~ι~619 ἵματα· ἱμάτια}
\lp{250 χρῆσον ἡμῖν τουτῳί]}{Ar.~\textit{Ach.}~414–415 ἀλλ᾽ ἀντιβολῶ πρὸς τῶν γονάτων σ᾽ Εὐριπίδη, / δός μοι ῥάκιόν τι τοῦ παλαιοῦ δράματος ; Ar.~\textit{Thesm.}~219 χρῆσόν τί νυν ἡμῖν ξυρόν ; Eur.~\textit{Hel.}~1248 τί σοι παράσχω δῆτα τῷ τεθνηκότι ; Ar.~\textit{Av.}~812 φέρ᾽ ἴδω, τί δ᾽ ἡμῖν τοὔνομ᾽ ἔσται τῇ πόλει ; Ar.~\textit{Thesm.}~428 νῦν οὖν ἐμοὶ τούτῳ δοκεῖ ; Eur.~\textit{El.}~191–192 καὶ παρ᾽ ἐμοῦ χρῆσαι πολύ- / πηνα φάρεα δῦναι ; Alciphr.~2.12 (III~15) δάνεισον οὖν μοι τούτων τοὺς περιττούς … εἰ οὖν δέοιο, προθύμως λάμβανε. τὸ γὰρ κοινὰ τὰ τῶν φίλων οὐχ ἥκιστα τοῖς ἀγροῖς ἐμφιλοχωρεῖν ὀφείλει ; Hsch.~κ~2828 κιχρήσει· δανείσει ; Suid.~δ~60 Δανεῖσαι· δοτικῇ. τούτου τὸ χρῆσαι αἱρετώτερον· τὸ μὲν γὰρ χρῆσαι ἐπὶ φίλων, τὸ δὲ δανεῖσαι πρὸς τοὺς τυχόντας ; Eur.~\textit{Or.}~735 συγκατασκάπτοις ἂν ἡμᾶς· κοινὰ γὰρ τὰ τῶν φίλων ; Pl.~\textit{Hp.Mi.}~364b μή σοι ἐμποδὼν εἴην ἐρωτῶν τῇ ἐπιδείξει ; Suid.~αι~377 Αἰτούμενος: κιχράμενος. Μένανδρος Ὕμνιδι· οὐ πῦρ γὰρ αἰτῶν, ἀλλὰ λοπάδα αἰτούμενος. τὸ μὲν αἰτῶν, εἰσαεὶ ἕξων, τὸ δὲ αἰτούμενος, λαβών, αὖθις ἀποδώσων ; Zen.~4.79 Κοινὰ τὰ φίλων: Τίμαιος φησὶν ὅτι προςιόντας Πυθαγόρᾳ μαθητὰς περὶ τὴν Ἰταλίαν ἔπειθεν ὁ φιλόσοφος κοινὰς τὰς οὐσίας ποιεῖσθαι ; Tosi n°~433 Communia esse amicorum inter se omnia}
\lp{251 στρόφιον]}{Ar.~\textit{Lys.}~931 τὸ στρόφιον ἤδη λύομαι. μέμνησό νυν ; Ar.~\textit{Thesm.}~139 τί λήκυθος καὶ στρόφιον; ὡς οὐ ξύμφορον ; Ar.~\textit{Thesm.}~255 αἶρε νῦν στρόφιον ; Ar.~\textit{Thesm.}~638 χάλα ταχέως τὸ στρόφιον ὦναίσχυντε σύ ; Cic.~\textit{Har.\,resp.}~44 P. Clodius a crocota, a mitra, a muliebribus soleis purpureisque fasceolis, a strophio, a psalterio, a flagitio, a stupro est factus repente popularis ; Hsch.~σ~2044 στρόφιον· ἡ στρογγύλη ζώνη (Ar. Lys. 931. Thesm. 638. fr. 664, 1 K.–A.) ; Ar.~fr.~332,~4~K.-A. (=~320,~4~Kock) στρόφι᾽, ὀπισθοσφενδόνην ; Ar.~fr.~664~K.-A. (=~647~Kock ; Poll.~7.67) ἀλλὰ τὸ στρόφιον λυθὲν / τὰ κάρυά μου ᾽ξέπιπτεν ; \textit{idem} Poll.~7.67 (praem. Τὸ μὲν γυναικεῖον ζώνιον οὕτω τε ἐκάλουν καὶ στροφίον ὠνόμαζον, ὡς Ἀριστοφάνης ; στροφίον ; ἐξέπιπτεν) ; Pherecr.~fr.~106~K.-A. (Λῆροι) ὄχθοιβος· λῶμα. Φερεκράτης Λήροις … μίτραν ἁλουργῆ, στρόφιον, ὄχθοιβον, κτένα ; Catull.~64.65 non tereti strophio lactentis vincta papillas ; Hsch.~α~6298 ἀπόδεσμος· κόσμιόν τι γυναικεῖον, περικεφάλαιον ; Poll.~7.66 ἄντικρυς δὲ τὸ νῦν καλούμενον ὑπὸ τῶν γυναικῶν στηθόδεσμον, εὕροις δ᾽ ἂν ὀνομαζόμενον ἀπόδεσμον ἐν Θεσμοφοριαζούσαις Ἀριστοφάνους ; Phot.~σ~639 στρόφιον· τὸ στρογγύλον ζωνάριον. Ἀριστοφάνης (Lys. 931)· „τὸ στρόφιον ἤδη λύομαι“ ; \textit{idem} Suid.~σ~1221 (λυόμενον), EM~730.56 ; Plaut.~\textit{Aul.}~516 strophiarii astant, astant semul sonarii ; Hsch.~ζ~236 ζώνιον· γυναικεῖον ζῶμα (Ar. Lys. 72) ; Hsch.~τ~35 ταινία⟨ι⟩· στέμματα τῶν ἱερέων, διαδήματα ἀρχιερατικά, ἢ ζῶναι, ἢ στέφανοι ; Poll.~5.96 καὶ στρόφιον καὶ ὀπισθοσφενδόνην παρ᾽ Ἀριστοφάνους ; Phot.~ο~744 ὀχθοίβους· τὰ λώματα· ἔστι δὲ περὶ τὸ στῆθος τοῦ χιτῶνος ἁλουργὲς πρόσραμμα ; Suid.~β~110 παροιμία ἐπὶ τῶν ἀνομοίων. Ἀριστοφάνης· τίς ἡ στολή; … τί λήκυθος καὶ στρόφιον; ὡς ἀξύμφορον}
\lp{251 οὐ γὰρ ταῦτά γ᾽ ὡς οὐκ ἔστ᾽ ἐρεῖς]}{Aesch.~\textit{PV}~298–299 οὐ γάρ ποτ᾽ ἐρεῖς ὡς Ὠκεανοῦ / φίλος ἐστὶ βεβαιότερός σοι ; Pl.~\textit{Smp.}~216a καὶ ταῦτα, ὦ Σώκρατες, οὐκ ἐρεῖς ὡς οὐκ ἀληθῆ ; Soph.~\textit{OC}~1095–1096 τῷ σκοπῷ μὲν οὐκ ἐρεῖς / ὡς ψευδόμαντις}
